\documentclass[11pt,twoside,openright]{memoir}
% Repository typography: EB Garamond, loaded immediately after
% \documentclass per ecosystem invariant XII.  `localtheorems' keeps
% platonic.sty's own theorem environments.
\usepackage[localtheorems]{raeez-math-template}
\usepackage{platonic}
\input{integrated_macros}
\title{\texorpdfstring{\textbf{A-infinity Chiral Algebras and 3D Holomorphic--Topological QFT}\\[0.55em]\Large Volume II of the programme: Hamiltonian jets, shifted cotangent fields, universal branes, and the open--closed map}{A-infinity Chiral Algebras and 3D Holomorphic--Topological QFT: Volume II of the programme: Hamiltonian jets, shifted cotangent fields, universal branes, and the open--closed map}}
\author{Raeez Lorgat}
\date{26 July 2026}
\hypersetup{pdftitle={Mixed Holomorphic-Topological Theory},pdfauthor={Raeez Lorgat},pdfsubject={Hamiltonian jets, shifted cotangent fields, universal branes, and open-closed comparison},pdfkeywords={A-infinity chiral algebra, holomorphic-topological field theory, Hamiltonian jet, shifted cotangent, open-closed map}}
\begin{document}
\frontmatter
\maketitle
\thispagestyle{empty}
\cleardoublepage
\chapter*{Abstract}
\addcontentsline{toc}{chapter}{Abstract}

A bulk observable acts on boundary conditions and on the operators between them. Recovering bulk observables from this action requires a complex of coherent boundary actions and a comparison constructed from the field theory. This volume begins with finite Hamiltonian jets on the formal symplectic plane. Their adjoint traces vanish, so the corresponding finite shifted-cotangent action satisfies the quantum master equation. The next calculation compares the algebraic boundary quantization with its relative Hochschild cochains. For a finite-dimensional complex Lie algebra, antisymmetrization identifies the relative Hochschild complex of its Rees enveloping algebra with the adjoint Chevalley complex whose differential is multiplied by the quantization parameter $\hbarq$. This quasi-isomorphism holds over every truncated coefficient ring $\C[\hbarq]/(\hbarq^{N+1})$ and passes to the $\hbarq$-adic limit.

The quadratic calculation determines an associative relation that the Lie brackets alone do not detect. For the Moyal product $\star$ with $\{q,p\}=1$, set $e=q^2/2$, $f=-p^2/2$, and $h=-qp$. Then
\[
 e\star f+f\star e+\frac12h\star h=-\frac38\hbarq^2.
\]
The corresponding Casimir quotient of the Rees enveloping algebra of $\mathfrak{sl}_2$ is the even Weyl algebra. Scalar reduction and an augmentation give continuous stable trace maps of commutative algebras; monoidal module fusion uses a separately specified completed bialgebra. Extending these constructions to Hamiltonian jets and mixed open--closed operations requires compatible cochain transition maps, continuous limit identifications, and operadic refinements. Geometric and physical realizations additionally require their comparison maps and orientation, descent, and quantization data. The finite calculations thus determine concrete algebraic constraints on the proposed comparisons.

\chapter*{Operations and reconstruction}
\addcontentsline{toc}{chapter}{Operations and reconstruction}
A boundary algebra records compositions of boundary operators. Its derived centre records actions compatible with these compositions, including their higher homotopies. A selected boundary module instead has its own derived endomorphism algebra. These two constructions answer different questions: coherent action on all boundaries and reconstruction from a chosen boundary.

The field-theoretic problem starts with independently defined bulk observables and boundary conditions. It asks for an action map into the appropriate geometric centre. Its domain includes the support, position dependence, and allowed collisions of the insertions. Quantization must deform this map together with its source and target. An elementary example shows why separate quantizations are insufficient. Let $k$ be a field. Over the dual numbers $k[\varepsilon]/(\varepsilon^2)$, the free rank-two algebra with relation $x^2=\varepsilon$ deforms $k[x]/(x^2)$. The augmentation $x\mapsto0$ cannot lift to the undeformed coefficient algebra: a lift would send $x$ to $\varepsilon c$, whose square is zero. Thus the action itself has a deformation obstruction.

For a finite-dimensional complex Lie algebra, the relative Rees comparison gives a quasi-isomorphism from relative Hochschild cochains to the Chevalley complex whose differential is scaled by the Rees parameter. Quadratic Hamiltonians generate the even Weyl algebra through the oscillator Casimir quotient. Scalar reduction and an augmentation give continuous stable trace maps of commutative algebras. Monoidal module fusion uses a separately specified coassociative, counital bialgebra in the chosen completed coefficient category.

The Hamiltonian jet limit requires compatible cochain transition maps and limit identifications. Higher mixed open--closed operations require the stated operadic refinements. A Hochschild centre and the Drinfeld centre of a monoidal module category are separate comparison targets. Hall, quantum-double, and physical bulk identifications require their actual comparison morphisms and the relevant orientation, descent, pairing, and analytic hypotheses. A proper filtration specifies an inverse-limit carrier; compatibility of the operations with that limit remains part of each comparison theorem.

\chapter*{Notation and normalization}
\addcontentsline{toc}{chapter}{Notation and normalization}
The ground field has characteristic zero.  Stable categories are presentable and idempotent-complete.  Tensor products preserve colimits separately.  Completions are separated inverse limits of finite weight or proper-height quotients.  The Igusa normalization is
\[
 \Dtheta=64\Dmon,\qquad \ChiTen=\Dmon^2,
\]
with $\Dmon$ monic.  Physical charge is additive; Gram degree is carried through its polarization cocycle.  Bar, factorization, and Hopf coproducts retain distinct notation and distinct geometric meanings.

\tableofcontents
\mainmatter
\chapter*{Hamiltonian jets and boundary quantization}
\addcontentsline{toc}{chapter}{Hamiltonian jets and boundary quantization}

Quantization of a symplectic boundary changes the product of observables while retaining the classical Poisson bracket to first order. The formal symplectic plane gives a concrete setting in which to compare Hamiltonian operations with operations on a quantized brane. Pointed Hamiltonians vanish to order at least two; translations belong to a separately adjoined affine sector. Finite Hamiltonian jets and their shifted cotangent fields provide the local algebraic models.

The universal boundary brane leads to a Rees enveloping algebra. Its relative Hochschild complex is compared with a Chevalley complex whose differential carries the same Rees parameter. Duflo quantization gives the separate invariant comparison, and the quadratic Moyal calculation determines the oscillator Casimir quotient. These finite constructions distinguish the underlying cochain comparison from higher operations on it.

The all-arity open--closed problem asks for compatible closed and mixed operations on the boundary algebra. Its stated refinement supplies the operadic maps and coherence homotopies. Passing through the Hamiltonian jet limit also requires actual cochain transition systems and their continuous limit identifications. A geometric brane realization must supply the boundary module and compare its centre action with the chosen observable or trace maps; global realization further requires descent and quantization data.

\cleardoublepage
\phantomsection
\addcontentsline{toc}{part}{Mixed Holomorphic--Topological Deligne Theory}
\chapter*{Book entrance: Mixed Holomorphic--Topological Deligne Theory}
\addcontentsline{toc}{chapter}{Book entrance: Mixed Holomorphic--Topological Deligne Theory}
\begin{summarybox}
Hamiltonian jets, shifted cotangent fields, universal branes, Duflo quantization, and the all-arity open--closed map.
\end{summarybox}
\part{Finite Hamiltonian geometry}

\chapter{The formal symplectic plane}
\section{Hamiltonians and the Poisson bracket}
Let $V=\C^2$ with coordinates $(z_1,z_2)$ and symplectic form
\[
 \omega=\dd z_1\wedge\dd z_2.
\]
Let $\mathfrak m=(z_1,z_2)$.  The pointed formal Hamiltonian Lie algebra is
\[
 \h=\mathfrak m^2\subset\C\llbracket z_1,z_2\rrbracket,
 \qquad
 \{f,g\}=\partial_{z_1}f\,\partial_{z_2}g-
          \partial_{z_2}f\,\partial_{z_1}g.
\]
It consists of Hamiltonians whose vector fields fix the chosen Darboux point.  The shifted order filtration $F^p\h=\mathfrak m^{p+2}$ is complete and obeys
\[
 \{F^p\h,F^q\h\}\subset F^{p+q}\h.
\]
The affine Hamiltonian algebra $\h^{\mathrm{aff}}=\C\llbracket z_1,z_2\rrbracket/\C$ also contains translations.  Its linear sector is adjoined separately when a moving brane or center-of-mass coordinate is required.

Write $R_r=\Sym^r(V^\vee)$ for the homogeneous polynomials of degree $r$.  The bracket obeys
\[
 \{R_r,R_s\}\subset R_{r+s-2}.
\]
For $d\ge2$ define the finite jet quotient
\[
 \h_{\le d}=\bigoplus_{r=2}^{d}R_r,
\]
where every bracket component of degree greater than $d$ is set equal to zero.

\begin{theorem}[Finite Hamiltonian jet algebra]
The truncated bracket makes $\h_{\le d}$ a finite-dimensional Lie algebra.  Its dimension is
\[
 \dim\h_{\le d}=\frac{d(d+3)}2-2.
\]
The quadratic part is
\[
 R_2\cong\mathfrak{sl}_2,
\]
and each $R_r$ is the irreducible $\mathfrak{sl}_2$-module $\Sym^r(\C^2)$.
\end{theorem}
\begin{proof}
The Jacobiator of the polynomial Poisson bracket vanishes before truncation.  Projection to degrees at most $d$ preserves this identity because the discarded terms form a Lie ideal.  The dimension is
\[
 \sum_{r=2}^{d}(r+1)
 =\frac{d(d+3)}2-2.
\]
Quadratic Hamiltonians act linearly on $V$ and preserve $\omega$; this identifies $R_2$ with $\mathfrak{sp}(V)=\mathfrak{sl}_2$.  The standard action on homogeneous polynomials gives the final statement.
\end{proof}

\section{Unimodularity}
\begin{theorem}[Unimodularity of every finite jet]
For every $d\ge2$ and every $x\in\h_{\le d}$,
\[
 \Tr_{\h_{\le d}}(\ad_x)=0.
\]
\end{theorem}
\begin{proof}
A homogeneous Hamiltonian of degree greater than two raises polynomial degree.  Its adjoint action is strictly upper triangular in the decomposition $\bigoplus_{r=2}^{d}R_r$, hence has trace zero.  A quadratic Hamiltonian belongs to $\mathfrak{sl}_2$.  Its action on each finite-dimensional $\mathfrak{sl}_2$-module $R_r$ has trace zero.  Linearity gives the result.
\end{proof}

\begin{corollary}[Canonical finite BV measure]
The linear odd symplectic space $T^*[-1]\h_{\le d}[1]$ has a translation-invariant Berezinian for which the cotangent BF Hamiltonian satisfies the quantum master equation.
\end{corollary}
\begin{proof}
Choose dual linear coordinates on the shifted cotangent space and use their constant Berezinian.  The classical term vanishes by Jacobi, while the BV divergence is the adjoint trace.  Unimodularity makes that trace zero.
\end{proof}

\chapter{The shifted cotangent theory}
\section{The Lie algebra and its functions}
Put
\[
 \g_d=T^*[-1]\h_{\le d}
 =\h_{\le d}\ltimes\h_{\le d}^\vee[-1],
\]
where $\h_{\le d}$ acts on the dual by the coadjoint action and the shifted dual is abelian.

\begin{proposition}[Chevalley--Poisson identification]
There is a canonical isomorphism of commutative dg algebras
\[
 C^*_{\CE}(\g_d)
 \cong
 C^*\!\left(\h_{\le d},\Sym(\h_{\le d})\right).
\]
The differential on the right is the Chevalley differential for the adjoint action on $\Sym(\h_{\le d})$.  Under the identification $\Sym(\h_{\le d})=\Ocal(\h_{\le d}^*)$, it is the linear Poisson differential.
\end{proposition}
\begin{proof}
The dual of $\h_{\le d}^\vee[-1]$ shifted by $[-1]$ is $\h_{\le d}$ in degree zero.  Chevalley cochains therefore contain polynomial functions on $\h_{\le d}^*$ and exterior ghosts for $\h_{\le d}$.  The semidirect bracket produces the coadjoint action, which is the Hamiltonian action of the linear Poisson tensor.
\end{proof}

\section{The BV action}
Choose dual bases $e_a$ and $e^a$.  Let $c^a$ be the ghost coordinates and $b_a$ the cotangent coordinates.  The finite action is
\[
 S_d=\frac12 b_c f^c{}_{ab}c^ac^b.
\]
The canonical BV bracket satisfies
\[
 \frac12\{S_d,S_d\}=0
\]
by Jacobi.

\begin{theorem}[Exact finite quantum master equation]
With the translation-invariant BV Laplacian,
\[
 \Delta S_d=\sum_b\Tr(\ad e_b)c^b=0.
\]
Consequently
\[
 \frac12\{S_d,S_d\}+\hbarq\Delta S_d=0
\]
for every $d$.
\end{theorem}
\begin{proof}
Contracting one field and one antifield in the cubic Hamiltonian gives the divergence of the adjoint vector field.  The preceding unimodularity theorem makes every coefficient vanish.
\end{proof}

\chapter{The universal boundary brane}
\section{Polarization}
Place the cotangent theory on a half-plane.  The standard Lagrangian boundary condition fixes the cotangent field and leaves the $\h_{\le d}$ field free.  Boundary observables form the Rees enveloping algebra
\[
 U_\hbar(\h_{\le d})
 =T(\h_{\le d})\llbracket\hbarq\rrbracket/
 \angles{xy-yx-\hbarq[x,y]}.
\]
The Poincar\'e--Birkhoff--Witt symmetrization gives a filtered vector-space isomorphism
\[
 \operatorname{sym}_\hbar:\Sym(\h_{\le d})\llbracket\hbarq\rrbracket
 \xrightarrow{\sim}U_\hbar(\h_{\le d}).
\]

\begin{theorem}[Boundary quantization]
The factorization algebra of boundary observables of the finite cotangent theory is locally constant and its value on an interval is $U_\hbar(\h_{\le d})$.  Fusion of adjacent intervals is multiplication in the enveloping algebra.
\end{theorem}
\begin{proof}
Canonical quantization turns the linear boundary phase space $\h_{\le d}^*$ into the Gutt star algebra.  For a linear Poisson bracket the Gutt product is transported by PBW to $U_\hbar(\h_{\le d})$.  Excision identifies ordered interval fusion with multiplication.
\end{proof}

\part{Relative Hochschild comparison}

\chapter{Hochschild cochains of the universal brane}
\section{The coefficient ring and the Chevalley differential}
The Rees relation scales both the Lie bracket and its adjoint action.
Its Chevalley differential must retain the same parameter.
For example, take the quadratic jet $\flie=\mathfrak{sl}_2$ with
$[h,e]=2e$.  Modulo $\hbarq$, the enveloping algebra is commutative,
so the Hochschild differential of $e$ vanishes.
The unscaled Chevalley differential instead satisfies
$(d_{\CE}e)(h)=2e$.

Let $\flie$ be a finite-dimensional complex Lie algebra.
For $N\ge0$, put
\[
 R_N=\C[\hbarq]/(\hbarq^{N+1}),\qquad
 \flie_N=R_N\tensor_\C\flie,\qquad
 [x,y]_N=\hbarq[x,y],
\]
and
\[
 A_N=U_{R_N}(\flie_N)
 =T_{R_N}(\flie_N)/(xy-yx-\hbarq[x,y]).
\]
Here $N$ denotes the order in $\hbarq$, independently of a Chan--Paton rank.
Write $\bar A_N$ for the augmentation ideal and $S=\Sym_\C\flie$.
PBW symmetrization is the $R_N$-module isomorphism
\[
 \operatorname{sym}_N:S\tensor R_N\longrightarrow A_N,
 \qquad
 x_1\cdots x_r\longmapsto
 \frac1{r!}\sum_{\sigma\in S_r}x_{\sigma(1)}\cdots x_{\sigma(r)}.
\]
The source has its ordinary symmetric product.
Use normalized relative Hochschild cochains
\[
 C^n_{\HH}(A_N/R_N,A_N)
 =\operatorname{Hom}_{R_N}(\bar A_N^{\tensor_{R_N}n},A_N)
\]
with differential
\[
\begin{split}
 (bF)(a_1,\ldots,a_{n+1})
 ={}&a_1F(a_2,\ldots,a_{n+1})\\
 &+\sum_{i=1}^n(-1)^iF(a_1,\ldots,a_i a_{i+1},\ldots,a_{n+1})\\
 &+(-1)^{n+1}F(a_1,\ldots,a_n)a_{n+1}.
\end{split}
\]
Let $d_{\CE}$ act on
$C^n(\flie,S)=\operatorname{Hom}_\C(\Lambda^n\flie,S)$
by the adjoint Chevalley differential.

\begin{theorem}[Relative Rees comparison]\label{thm:vol2-relative-rees}
Antisymmetrization defines a quasi-isomorphism of $R_N$-complexes
\[
 \operatorname{Alt}_N:
 C^\bullet_{\HH}(A_N/R_N,A_N)
 \longrightarrow
 \bigl(C^\bullet(\flie,S)\tensor R_N,\hbarq d_{\CE}\bigr),
\]
given by
\[
 (\operatorname{Alt}_NF)(x_1,\ldots,x_n)
 =\operatorname{sym}_N^{-1}
 \sum_{\sigma\in S_n}\operatorname{sgn}(\sigma)
 F(x_{\sigma(1)},\ldots,x_{\sigma(n)}).
\]
These maps commute with coefficient reduction in $N$.
\end{theorem}
\begin{proof}
Write $A=A_N$, $R=R_N$, and $A^e=A\tensor_R A^{\op}$.
In homological degree $n$, set
\[
 K_n=A\tensor_R\Lambda_R^n\flie_N\tensor_R A.
\]
Its augmentation is multiplication.
For $X=x_1\wedge\cdots\wedge x_n$, let $X_{\widehat i}$ and
$X_{\widehat i\widehat j}$ omit the indicated entries without reordering.
Define
\[
\begin{split}
 \partial(a\tensor X\tensor b)
 ={}&\sum_i(-1)^{i+1}
 \bigl(ax_i\tensor X_{\widehat i}\tensor b
       -a\tensor X_{\widehat i}\tensor x_i b\bigr)\\
 &+\hbarq\sum_{i<j}(-1)^{i+j}
 a\tensor[x_i,x_j]\wedge X_{\widehat i\widehat j}\tensor b.
\end{split}
\]
The normalized relative bimodule bar has terms
$\mathcal B_n=A\tensor_R\bar A^{\tensor_R n}\tensor_R A$.
Define
\[
 \iota_n(a\tensor X\tensor b)
 =\sum_{\sigma\in S_n}\operatorname{sgn}(\sigma)
       a[x_{\sigma(1)}|\cdots|x_{\sigma(n)}]b.
\]
The outer bar faces give the first sum in $\partial$.
Pair adjacent transpositions in the internal faces and use
$xy-yx=\hbarq[x,y]$ to obtain the second sum.
For two inputs this reads
\[
\begin{split}
 b\bigl(1[x|y]1-1[y|x]1\bigr)
 ={}&x[y]1-y[x]1+1[x]y-1[y]x\\
 &-\hbarq\,1[[x,y]]1.
\end{split}
\]
Thus $b\iota=\iota\partial$.
PBW makes $\flie_N$ a free summand of $\bar A$.
Since every factorial is invertible in $R$, $\iota$ is injective.
The identity $b^2=0$ then gives $\partial^2=0$.

Give $K_n$ filtration degree equal to its outer PBW degrees plus $n$.
The associated graded is the diagonal Koszul complex of
$\Sym_R(\flie_N)$.
Its differential uses differences of left and right polynomial generators.
A linear change of variables makes these independent polynomial variables.
They form a regular sequence even when $R$ has nilpotents.
For an augmented cycle, lift a primitive of its highest-degree symbol
and subtract its differential.
The filtration degree decreases, so this procedure terminates.
Hence $K_\bullet\to A$ is a free bimodule resolution.
The relative bar is also free over $A^e$, with its usual unit-insertion contraction.
The map $\iota$ lifts the identity of $A$.
Projectivity constructs its chain homotopy inverse degree by degree.
Applying $\operatorname{Hom}_{A^e}(-,A)$ gives a cochain homotopy equivalence.

The induced differential on
$\operatorname{Hom}_R(\Lambda_R^n\flie_N,A)$ is
\[
\begin{split}
 (\delta_N\phi)(x_1,\ldots,x_{n+1})
 ={}&\sum_i(-1)^{i+1}
 [x_i,\phi(x_1,\ldots,\widehat{x_i},\ldots,x_{n+1})]_A\\
 &+\hbarq\sum_{i<j}(-1)^{i+j}
 \phi([x_i,x_j],x_1,\ldots,\widehat{x_i},
                         \ldots,\widehat{x_j},\ldots,x_{n+1}).
\end{split}
\]
Commuting $x$ past the factors of a symmetrized monomial gives
\[
 [x,\operatorname{sym}_N(f)]_A
 =\hbarq\operatorname{sym}_N(\operatorname{ad}_x f).
\]
Both sums therefore become $\hbarq$ times the ordinary Chevalley differential.
Precomposition with $\iota$ gives $\operatorname{Alt}_N$.
All formulas commute with coefficient reduction.
\end{proof}

\section{The formal coefficient limit}
Put $R=\C\llbracket\hbarq\rrbracket$ and $A=\varprojlim_N A_N$.
PBW identifies its underlying module with $S\llbracket\hbarq\rrbracket$.
Each coefficient of $\hbarq$ is a polynomial.
Define
\[
 A^{\widehat\tensor_R n}=\varprojlim_N A_N^{\tensor_{R_N}n},
 \qquad
 C^\bullet_{\HH,\mathrm{cont}}(A/R,A)
 =\varprojlim_N C^\bullet_{\HH}(A_N/R_N,A_N).
\]
Equivalently, these are continuous $R$-multilinear normalized cochains
on the completed relative tensors.

\begin{corollary}[Formal relative Rees comparison]\label{cor:vol2-formal-relative-rees}
The compatible antisymmetrizations induce a quasi-isomorphism
\[
 C^\bullet_{\HH,\mathrm{cont}}(A/R,A)
 \longrightarrow
 \bigl(C^\bullet(\flie,S)\llbracket\hbarq\rrbracket,
                    \hbarq d_{\CE}\bigr).
\]
\end{corollary}
\begin{proof}
Coefficient reduction is surjective in each cochain degree.
For Hochschild cochains, lift outputs on the PBW input basis.
The Chevalley exterior input has finite rank, so its outputs also lift.
The comparison cones are acyclic and their transition maps are degreewise surjective.
The long exact sequence shows that every transition kernel is acyclic.
Choose a primitive of a compatible cycle at level zero.
Lift it one level and correct its differential by a primitive in the transition kernel.
Iteration gives compatible primitives at every order.
The inverse-limit cone is therefore acyclic.
\end{proof}

The coefficient ring is part of this result.
Absolute $\C$-linear cochains also detect variations of the coefficient parameter.
For $\flie=0$, the continuous absolute derivation $\partial_{\hbarq}$
on $R$ has no relative counterpart.
For nonabelian $\flie$, the associated graded for the $\hbarq$ filtration
has zero differential, whereas $d_{\CE}$ need not vanish.

\chapter{The Kontsevich--Duflo map}
\section{Linear Poisson formality}
For $\flie=\h_{\le d}$, put $A_d=U_\hbar(\h_{\le d})$ and
\[
 \mathcal P_d=
 \bigl(C^*(\h_{\le d},\Sym\h_{\le d})
       \llbracket\hbarq\rrbracket,\hbarq d_{\CE}\bigr).
\]
This is the polynomial polyvector complex on $\h_{\le d}^*$ with
differential $[\hbarq\pi_d,-]$ for the linear Poisson tensor $\pi_d$.
The relative comparison gives $C^\bullet_{\HH,\mathrm{cont}}(A_d/R,A_d)
\to\mathcal P_d$ as a map of complexes.
Higher operations require a formality map with specified operadic structure.

Fix Kontsevich's formality morphism $\mathcal U$ on the affine space
$\h_{\le d}^*$.
Its twist by $\hbarq\pi_d$ defines a star product $\star_d$ and an
$\Linf$ quasi-isomorphism on the degree-shifted complexes.
Let $T_d:(\Sym\h_{\le d}\llbracket\hbarq\rrbracket,\star_d)
\to A_d$ be the algebra identification fixing the linear generators.
This identification includes the equivalence between the Kontsevich and PBW products.
PBW symmetrization alone identifies the underlying modules.
The first Taylor coefficient of the twisted comparison is
\[
 \Phi_d(v)=(T_d)_*
 \sum_{r\ge0}\frac1{r!}
 \mathcal U_{r+1}
 (\underbrace{\hbarq\pi_d,\ldots,\hbarq\pi_d}_{r},v).
\]
It has source $\mathcal P_d$ and target
$C^\bullet_{\HH,\mathrm{cont}}(A_d/R,A_d)$.
At each order in $\hbarq$, polynomial inputs have polynomial outputs.
Its reduction modulo $\hbarq$ is the Hochschild--Kostant--Rosenberg map,
with the normalized antisymmetrization convention.
That map is a quasi-isomorphism for the polynomial algebra.
The complete $\hbarq$ filtration then proves that $\Phi_d$ is a quasi-isomorphism.
The higher Taylor coefficients supply the $\Linf$ identities after shifting by $[1]$.
These statements use the affine formality construction
\cite{KontsevichDQ}*{Sections~6 and~8.3}.

On invariant polynomials, the standard Duflo normalization gives
\[
 \operatorname{Duf}_\hbar
 =\operatorname{sym}_\hbar\circ j^{1/2}(\hbarq\partial),
 \qquad
 j^{1/2}(x)=
 \det_{\h_{\le d}}^{1/2}
 \left(\frac{\sinh(\operatorname{ad}_x/2)}
 {\operatorname{ad}_x/2}\right).
\]
It is an algebra isomorphism from
$(\Sym\h_{\le d})^{\h_{\le d}}\llbracket\hbarq\rrbracket$
to $Z(A_d)$, by the Duflo theorem
\cite{KontsevichDQ}*{Section~8.3.4}.
The square root has constant term one.
This assertion concerns invariant degree-zero cohomology.
It does not identify antisymmetrization with a homotopy-$\Etwo$ map.

\section{The boundary action}
The brace operad acts on
$(C^\bullet_{\HH,\mathrm{cont}}(A_d/R,A_d),A_d)$.
Its closed color is the relative derived centre.
A homotopy-$\Etwo$ refinement means a morphism or zigzag
\[
 \widetilde\Phi_d:\mathcal P_d
 \simeq C^\bullet_{\HH,\mathrm{cont}}(A_d/R,A_d)
\]
in complete $R$-linear homotopy-$\Etwo$ algebras,
with the chosen Gerstenhaber structure on $\mathcal P_d$.
A compatible Swiss-cheese refinement also specifies the mixed operations
on the pair and their coherence homotopies.

\begin{corollary}[Open--closed structure with a refinement]\label{cor:vol2-open-closed-refinement}
Given these refinements, the pair $(\mathcal P_d,A_d)$ has the transported
homotopy Swiss-cheese action.
Its open product is the enveloping product.
The specified closed and mixed Taylor coefficients determine its
closed bracket and boundary action.
\end{corollary}
\begin{proof}
Transport the two-colored operations through the supplied equivalence
of complete operadic algebras.
The coherence homotopies are part of this equivalence.
\end{proof}

\chapter{The formal Hamiltonian limit}
\section{Weight completion and comparison data}
The Hamiltonian Lie algebra and its boundary algebra have quotient systems
\[
 \h=\varprojlim_d\h_{\le d},\qquad
 \widehat U_\hbar(\h)=\varprojlim_d U_\hbar(\h_{\le d}).
\]
A Lie quotient does not by itself define a covariant map between
Chevalley cochains with adjoint coefficients.
Nor does an algebra quotient define a map between all Hochschild
cochains with values in the algebra itself.
A cochain must descend through the input quotient and the output quotient.
Consequently the coefficient limit of
Corollary~\ref{cor:vol2-formal-relative-rees} is distinct from this jet limit.

A \emph{compatible jet comparison} consists of the following data.
Choose complete $R$-linear cochain models $P_d,H_d$ with transition maps,
levelwise comparisons $P_d\simeq\mathcal P_d$ and
$H_d\simeq C^\bullet_{\HH,\mathrm{cont}}(A_d/R,A_d)$,
and compatible quasi-isomorphisms $P_d\to H_d$.
Require degreewise surjective transitions, or an explicit derived
inverse-limit construction with the same exactness property.
Specify identifications
\[
 \varprojlim_dP_d\simeq\mathcal P_\infty,
 \qquad
 \varprojlim_dH_d\simeq
 C^\bullet_{\HH,\mathrm{cont}}
 (\widehat U_\hbar(\h)/R,\widehat U_\hbar(\h)),
\]
where $\mathcal P_\infty$ is the proposed continuous Hamiltonian
polyvector complex with differential $[\hbarq\pi,-]$.
The topology on the right combines the stated Hamiltonian weight
completion with the $\hbarq$ completion.
The identifications must use this same topology.
For higher operations, also supply compatible homotopy-$\Etwo$ and
Swiss-cheese refinements whose operations converge in that topology.

\begin{theorem}[Continuous comparison from compatible jets]\label{thm:vol2-compatible-jets}
A compatible jet comparison induces a continuous quasi-isomorphism
\[
 \Phi_\infty:\mathcal P_\infty
 \xrightarrow{\sim}
 C^\bullet_{\HH,\mathrm{cont}}
 (\widehat U_\hbar(\h)/R,\widehat U_\hbar(\h)).
\]
Compatible homotopy-$\Etwo$ and Swiss-cheese refinements give the
corresponding complete open--closed operations.
\end{theorem}
\begin{proof}
For surjective transitions, the cones of the levelwise maps are
acyclic inverse systems with surjective maps in each degree.
The primitive-lifting argument of
Corollary~\ref{cor:vol2-formal-relative-rees} proves exactness of their limit.
For a derived inverse limit, use its invariance under levelwise equivalences.
Apply the supplied identifications of the two limits.
Compatibility and convergence of the operadic maps give the final assertion.
\end{proof}

The finite PBW comparison constructs the relative centre at each fixed jet.
A continuous field-theory comparison additionally requires the jet datum
and a map from the specified field observables to $\mathcal P_\infty$.
Universal graph formulas alone do not provide the cochain transition maps.

\begin{principle}[Universal brane principle]
The boundary algebra $\widehat U_\hbar(\h)$ is the universal formal Hamiltonian brane.  Finite matrix branes are representation-theoretic images of this algebra.  Its complete relative derived centre is the algebraic target for bulk observables.  A field-theory identification requires the continuous comparison data stated above.
\end{principle}


\part{The finite-jet laboratory}

\chapter{The first three Hamiltonian jets}
\section{Quadratic Hamiltonians}
Set
\[
 e=\frac12z_1^2,
 \qquad
 f=-\frac12z_2^2,
 \qquad
 h=-z_1z_2.
\]
A direct differentiation gives
\[
 \{h,e\}=2e,
 \qquad
 \{h,f\}=-2f,
 \qquad
 \{e,f\}=h.
\]
Thus $\h_{\le2}=R_2$ is $\mathfrak{sl}_2$ with its standard Chevalley basis.  The finite cotangent theory for $d=2$ is the linear BV theory of $T^*[-1]\mathfrak{sl}_2$ and its boundary algebra is $U_\hbar(\mathfrak{sl}_2)$.

\section{Cubic Hamiltonians}
Let
\[
 v_j=z_1^{3-j}z_2^j,
 \qquad 0\le j\le3.
\]
Then
\[
 \{h,v_j\}=(3-2j)v_j,
 \qquad
 \{e,v_j\}=jv_{j-1},
 \qquad
 \{f,v_j\}=(3-j)v_{j+1},
\]
with the boundary terms understood as zero.  Hence $R_3$ is the irreducible highest-weight-three module.

\begin{proposition}[The cubic jet algebra]
The degree-three jet algebra is the semidirect product
\[
 \h_{\le3}\cong\mathfrak{sl}_2\ltimes\Sym^3(\C^2),
\]
where the second summand is abelian.
\end{proposition}
\begin{proof}
The bracket of two cubic polynomials has degree four and vanishes in the degree-three quotient.  The bracket with a quadratic polynomial is the displayed $\mathfrak{sl}_2$ action.
\end{proof}

\section{The first nonlinear layer}
Write $w_k=z_1^{4-k}z_2^k$, $0\le k\le4$.  In $\h_{\le4}$ the cubic modes satisfy
\begin{equation}\label{eq:vol2-cubic-quartic}
 \{v_i,v_j\}=3(j-i)w_{i+j-1},
\end{equation}
whenever the index lies between $0$ and $4$; the coefficient makes every remaining case vanish.

\begin{calculation}[Two brackets]
\[
 \{z_1^3,z_2^3\}=9z_1^2z_2^2=9w_2,
 \qquad
 \{z_1^2z_2,z_1z_2^2\}=3z_1^2z_2^2=3w_2.
\]
The two independent cubic pairs reach the same quartic channel with amplitudes in the ratio $3:1$.
\end{calculation}
\begin{proof}
Differentiate each monomial in the two symplectic coordinates and substitute in the Poisson formula.  The first pair gives $3z_1^2\cdot3z_2^2$.  The second gives $2z_1z_2\cdot2z_1z_2-z_1^2\cdot z_2^2=3z_1^2z_2^2$.
\end{proof}

\begin{theorem}[The quartic jet presentation]
The Lie algebra $\h_{\le4}$ is generated by $e,f,h,v_0,v_1,v_2,v_3$.  Its relations are the $\mathfrak{sl}_2$ relations, the irreducible action on the $v_j$, formula \eqref{eq:vol2-cubic-quartic}, and the induced $\mathfrak{sl}_2$ action on $R_4$.  The quartic subspace is abelian in the quotient.
\end{theorem}
\begin{proof}
Quadratic and cubic generators produce every quartic monomial through \eqref{eq:vol2-cubic-quartic}.  Brackets involving a quartic and a cubic have degree five, and brackets of two quartics have degree six, so both vanish after projection.  Jacobi follows from the polynomial Poisson bracket before projection.
\end{proof}

\chapter{Moyal quantization in coordinates}
\section{The bidifferential formula}
For polynomials on $\C^2$, the Moyal product is
\begin{equation}\label{eq:vol2-moyal}
 f\star g
 =\sum_{n\ge0}\frac1{n!}\left(\frac{\hbarq}{2}\right)^n
 \sum_{r=0}^n(-1)^r\binom nr
 (\partial_1^{n-r}\partial_2^rf)
 (\partial_1^r\partial_2^{n-r}g).
\end{equation}
Every polynomial pair gives a finite sum.

\begin{theorem}[Quadratic exactness]
For polynomials $f,g$ of degree at most two,
\[
 [f,g]_\star=\hbarq\{f,g\}.
\]
The quadratic generators therefore define a Lie representation of the Rees-scaled $\mathfrak{sl}_2$.  Their associative algebra has an additional Casimir relation.
\end{theorem}
\begin{proof}
The commutator retains the odd powers of the Poisson bidifferential.  Derivatives of order three vanish on quadratics, so the first-order term is the entire commutator.
\end{proof}

\begin{proposition}[The oscillator quotient]\label{prop:vol2-oscillator-quotient}
Let $R_N=\C[\hbarq]/(\hbarq^{N+1})$ and let $W_N=R_N[q,p]$
have the Moyal product for $\{q,p\}=1$.
Let $U_N$ have generators $E,F,H$ with relations
\[
 [H,E]=2\hbarq E,\qquad [H,F]=-2\hbarq F,\qquad [E,F]=\hbarq H.
\]
Set $\Omega=EF+FE+H^2/2$.
The assignments
\[
 E\longmapsto e=q^2/2,\qquad
 F\longmapsto f=-p^2/2,\qquad
 H\longmapsto h=-qp
\]
induce an algebra isomorphism
\[
 U_N/(\Omega+3\hbarq^2/8)
 \xrightarrow{\sim} W_N^{(q,p)\mapsto(-q,-p)}.
\]
These compatible isomorphisms also hold after $\hbarq$-adic completion.
If linear Hamiltonians are retained, the generated associative algebra is the full Weyl algebra.
\end{proposition}
\begin{proof}
Quadratic exactness proves the three commutator relations.
The associative products are
\[
\begin{split}
 e\star f&=-q^2p^2/4-\hbarq qp/2-\hbarq^2/8,\\
 f\star e&=-q^2p^2/4+\hbarq qp/2-\hbarq^2/8,\\
 h\star h&=q^2p^2-\hbarq^2/4.
\end{split}
\]
Thus $\Omega$ maps to $-3\hbarq^2/8$.
Its commutator with each of $E,F,H$ vanishes by the defining relations.
Give these generators PBW degree one.
Give an even monomial $q^ap^b$ degree $(a+b)/2$.
The leading symbol of $\Omega+3\hbarq^2/8$ is
$C=2EF+H^2/2$.
It is a non-zero-divisor in $R_N[E,F,H]$, since $2C$ is monic in $H$.
Centrality and this property give
\[
 \operatorname{gr}\bigl(U_N/(\Omega+3\hbarq^2/8)\bigr)
 =R_N[E,F,H]/(C).
\]
Indeed, every element of the two-sided ideal is a multiple of the central generator.
Its leading symbol is the product of the two leading symbols.
The quotient has basis $E^aH^bF^c$ with $a,c\ge0$ and $b=0,1$.
Their leading symbols map to unit multiples of the distinct monomials
$q^{2a+b}p^{2c+b}$, precisely the even monomials.
Hence the associated-graded map is an isomorphism.
Lifting highest-degree terms proves surjectivity and comparing them proves injectivity.
Every element at fixed $N$ has finite polynomial degree.
The formulas commute with reduction in $N$, so their inverse limit gives the formal isomorphism.
Finally $q,p$ generate $W_N$ and topologically generate its formal limit.
\end{proof}

The primitive coproduct on $U_N$ does not descend to this quotient.
Modulo $\hbarq$, it sends the classical relation to
\[
 \Delta(C)=C\tensor1+1\tensor C
   +2E\tensor F+2F\tensor E+H\tensor H.
\]
The cross-term is nonzero in the tensor square of $\C[E,F,H]/(C)$.
For example, evaluation at $(E,F,H)=(1,0,0)$ in the first factor
and $(0,1,0)$ in the second gives $2$.
Thus the oscillator quotient is an associative algebra, with no coproduct inherited from these primitive generators.

\begin{calculation}[Cubic correction]
For $f=z_1^3$ and $g=z_2^3$,
\[
 f\star g-g\star f
 =9\hbarq z_1^2z_2^2+\frac32\hbarq^3.
\]
The first term quantizes the quartic Hamiltonian channel.  The constant is the first wheel-free ordering correction visible beyond the quadratic Lie algebra.
\end{calculation}
\begin{proof}
In \eqref{eq:vol2-moyal}, the odd orders are $n=1$ and $n=3$.  The first contributes $9\hbarq z_1^2z_2^2$.  The third contributes $2(\hbarq/2)^3(3!)=3\hbarq^3/2$.
\end{proof}

\section{Weyl generators}
With the convention $\{q,p\}=1$, the symbols satisfy
$[q,p]_\star=\hbarq$.
Their Weyl-ordered differential operators are
\[
 \widehat q=q,\qquad \widehat p=-\hbarq\partial_q,
 \qquad [\widehat q,\widehat p]=\hbarq.
\]
The abstract Weyl algebra is defined by this commutator relation.
Weyl ordering transports its multiplication to the Moyal product.
The differential-operator representation realizes these relations over
$\C[\hbarq]$ and its formal completion.
The finite quotient argument above uses the abstract $R_N$-algebra,
so it does not require a faithful differential-operator representation at nilpotent $\hbarq$.

\chapter{Chevalley coordinates and the Poisson differential}
Choose a basis $e_a$ of $\h_{\le d}$ with structure constants $[e_a,e_b]=f^c{}_{ab}e_c$.  Let $c^a$ be the odd Chevalley coordinate and let $x_a$ be the even coordinate on $\h_{\le d}^*$.  Then
\begin{equation}\label{eq:vol2-ce-coordinate}
 d_{\CE}c^c=-\frac12f^c{}_{ab}c^ac^b,
 \qquad
 d_{\CE}x_a=f^c{}_{ba}c^b x_c.
\end{equation}
The first equation encodes the Lie bracket.  The second is the adjoint action on $\Sym(\h_{\le d})$, viewed as functions on the dual.

\begin{proposition}[Hamiltonian interpretation]
For a polynomial $F(x)$,
\[
 d_{\CE}F
 =c^a\{x_a,F\}_{\mathrm{lin}}.
\]
Consequently $d_{\CE}^2=0$ is the Jacobi identity for the linear Poisson tensor.
\end{proposition}
\begin{proof}
The linear Poisson bracket is $\{x_a,x_b\}=f^c{}_{ab}x_c$.  The Leibniz rule gives \eqref{eq:vol2-ce-coordinate} on every polynomial.  The square is one half of the Schouten bracket of the Poisson tensor with itself.
\end{proof}

\section{The \texorpdfstring{$\mathfrak{sl}_2$}{sl2} invariant}
For the quadratic jet, identify the coordinate functions with $e,f,h$.  The invariant
\[
 q=2ef+\frac12h^2
\]
is $d_{\CE}$-closed.  Its formal Duflo image is $\Omega+\hbarq^2/2$, where $\Omega=ef+fe+h^2/2$ in the Rees enveloping algebra.  This is the degree-zero invariant comparison.  Under the oscillator quotient of Proposition~\ref{prop:vol2-oscillator-quotient}, its image is $\hbarq^2/8$.

\chapter{Residue duality and the formal current algebra}
\section{The continuous dual}
Let
\[
 \rho_{ab}=z_1^{-a-1}z_2^{-b-1}\,\dd z_1\dd z_2,
 \qquad a,b\ge0.
\]
The Grothendieck residue pairing is
\[
 \operatorname{Res}_0\left(z_1^iz_2^j\rho_{ab}\right)
 =\delta_{ia}\delta_{jb}.
\]
The continuous dual of $\C\llbracket z_1,z_2\rrbracket$ is the direct sum of the $\rho_{ab}$.  The continuous dual of the pointed algebra $\h=\mathfrak m^2$ is the span of the modes with $a+b\ge2$.  The two modes of total degree one pair with the translation sector of $\h^{\mathrm{aff}}$.

\begin{proposition}[Coadjoint current formula]
For $f=z_1^iz_2^j$ and $\rho_{ab}$,
\[
 \ad_f^*(\rho_{ab})(g)
 =-\rho_{ab}(\{f,g\}).
\]
In the residue basis this is a finite linear combination of modes of bidegree
\[
 (a-i+1,b-j+1).
\]
\end{proposition}
\begin{proof}
The definition is the coadjoint action.  The Poisson bracket lowers the total degree of $fg$ by two, so residue conservation fixes the displayed bidegree.  The Poisson formula contains two derivative terms.
\end{proof}

The shifted cotangent current algebra
\[
 \g=\h\ltimes\h^\vee_{\mathrm{cont}}[-1]
\]
is therefore a complete topological Lie algebra with a continuous invariant pairing.  Its finite jet quotients are precisely the $\g_d$ used above.

\chapter{All-arity open--closed operations}
\section{Hochschild braces}
For Hochschild cochains $F\in C^p(A,A)$ and $G_i\in C^{q_i}(A,A)$, the brace
\[
 F\{G_1,\ldots,G_r\}
\]
is the signed sum over ordered insertions of the $G_i$ into the inputs of $F$.  The brace identities state that a nested insertion is the sum over all ways of distributing the inner insertions among the outer slots.

\begin{theorem}[Brace realization with a refinement]
For $A=A_d$, supply the homotopy-$\Etwo$ and mixed refinements of
Corollary~\ref{cor:vol2-open-closed-refinement}.
They identify the selected homotopy Gerstenhaber operations on
$\mathcal P_d$ with the relative Hochschild centre operations.
The chosen mixed refinement also identifies their action on the boundary algebra.
\end{theorem}
\begin{proof}
The relative Hochschild complex has its brace operations.
The supplied operadic equivalence preserves their coherent homotopies.
An equivalence of complexes alone would not imply this conclusion.
\end{proof}

\section{The first mixed identities}
A relative Hochschild $1$-cocycle $D$ of an algebra concentrated in degree zero
is an $R$-linear derivation:
\[
 D(ab)=D(a)b+aD(b).
\]
For two such cocycles $D,E$, their Gerstenhaber bracket acts by
\[
 [D,E](a)=D(E(a))-E(D(a)).
\]
For a graded algebra and homogeneous derivations, these formulas become
\[
\begin{split}
 D(ab)&=D(a)b+(-1)^{|D||a|}aD(b),\\
 [D,E](a)&=D(E(a))-(-1)^{|D||E|}E(D(a)).
\end{split}
\]
Higher Hochschild cochains act through multilinear evaluations and braces.
They need not act by derivations or carry a coproduct defining a diagonal action.
A compatible mixed refinement transports these operations, including their homotopies.

\begin{corollary}[Continuous open--closed structure with compatible jets]
Under the jet comparison and mixed refinement hypotheses of
Theorem~\ref{thm:vol2-compatible-jets}, the pair
$(\mathcal P_\infty,\widehat U_\hbar(\h))$
has the corresponding complete homotopy Swiss-cheese action.
Its closed color is equivalent to the continuous relative derived centre.
\end{corollary}
\begin{proof}
The supplied finite operations commute with the transition maps
and converge in the specified completion.
Their limit acts on the two limit objects.
The comparison theorem identifies the closed complex with continuous relative Hochschild cochains.
\end{proof}

\chapter{Finite Chan--Paton ranks}
\section{Rank one}
For $N=1$, the commuting equation is automatic.  The boundary coordinate algebra is $\C[x,y]$, and
\[
 J_1(f)=f(x,y).
\]
The necklace bracket reduces to the ordinary Poisson bracket.  Restriction to $f\in\mathfrak m^2$ gives the defining pointed-Hamiltonian representation.  The two linear Hamiltonians act by translations and recover the affine center-of-mass extension.

\section{Rank two and Cayley--Hamilton reduction}
For a $2\times2$ matrix $X$,
\[
 X^2-(\Tr X)X+(\det X)I=0.
\]
Hence every trace containing two consecutive powers of $X$ reduces to traces with $X$-degree at most one after adjoining $\Tr X$ and $\det X$.  The same statement holds for $Y$.  On the commuting locus, simultaneous diagonalization identifies the invariant algebra with the symmetric polynomials in two ordered pairs.

\begin{proposition}[First finite-rank trace relation]
For commuting $2\times2$ matrices,
\[
 p_{20}=p_{10}^2-2\det X,
 \qquad
 p_{02}=p_{01}^2-2\det Y,
\]
and
\[
 p_{11}^2-p_{20}p_{02}
 =-\bigl(x_1y_2-x_2y_1\bigr)^2.
\]
The stable generators remain independent until a relation forced by the rank is reached.
\end{proposition}
\begin{proof}
Write the joint eigenvalues as $(x_1,y_1)$ and $(x_2,y_2)$ and expand both sides.
\end{proof}

\section{Stable passage}
The transition from rank $N$ to rank $N+1$ adds a zero eigenvalue.  Each fixed weighted degree stabilizes once $N$ exceeds that degree.  The inverse limit is therefore the completed multisymmetric algebra.  The stable trace is an exact quotient of the universal centre on the polynomial trace sector.


\part{Global holomorphic symplectic descent}

\chapter{Fedosov descent}
\section{The Weyl algebroid}
All Hochschild cochains and Morita centres in this chapter are continuous
and relative to $R=\C\llbracket\hbarq\rrbracket$.
The symplectic globalization uses its own sheafwise formality and
Fedosov resolution data, independently of the finite Hamiltonian comparison.
Let $(S,\omega)$ be a holomorphic symplectic surface.  The bundle of formal Weyl algebras has fibre the completed Weyl algebra of a formal symplectic plane.  A holomorphic symplectic connection and a Fedosov one-form define a flat connection
\[
 D_F=\nabla+\frac1\hbarq[\gamma,-]
\]
whose flat sections form a deformation-quantization algebroid $\mathcal W_{S,\hbar}$.

\begin{theorem}[Globalized formality on a holomorphic symplectic surface]
Fix a stable formality morphism and a compatible homotopy-$\Etwo$ refinement.  There is a sheafwise homotopy-$\Etwo$ quasi-isomorphism
\[
 \Phi_S:
 \left(\PV_S\llbracket\hbarq\rrbracket,[\hbarq\pi,-]\right)
 \xrightarrow{\sim}
 C^\bullet_{\HH}(\mathcal W_{S,\hbar},\mathcal W_{S,\hbar}),
\]
where $\pi=\omega^{-1}$.  A change of holomorphic symplectic connection or Fedosov one-form gives a gauge-equivalent resolution and a homotopic quasi-isomorphism.  Holomorphic symplectomorphisms act equivariantly on the resulting homotopy class.
\end{theorem}
\begin{proof}
Fedosov resolutions replace both sheaves by acyclic sheaves of formal fibrewise polyvectors and polydifferential operators.  The fixed covariant formality morphism acts fibrewise and intertwines the Fedosov differentials.  A change of connection or Fedosov one-form is a gauge transformation in the pronilpotent resolution and therefore gives a homotopy.
\end{proof}

\begin{corollary}[Global mixed-HT centre]
The deformation-quantization category $\Perf(\mathcal W_{S,\hbar})$ has Morita centre
\[
 Z_{\mathrm{Mor}}\bigl(\Perf(\mathcal W_{S,\hbar})\bigr)
 \simeq
 \R\Gamma\!
 \left(S,\PV_S\llbracket\hbarq\rrbracket,[\hbarq\pi,-]\right).
\]
For a K3 surface this computes the symplectic Poisson centre in the stated sheafwise model.  A Hamiltonian BF identification additionally requires a comparison from its Chevalley observables to this symplectic polyvector complex.
\end{corollary}
\begin{proof}
The Morita centre of perfect modules over an algebroid is the derived global section complex of its Hochschild cochains.  Apply the globalized formality equivalence and then derived global sections.  A K3 surface is holomorphic symplectic, so the same sheafwise calculation applies whenever the chosen globalization data are available.
\end{proof}

\chapter{Factorization and descent}
\section{The algebraic mixed-HT realization}
The sheaf
\[
 \mathcal P_S=
 \left(\PV_S\llbracket\hbarq\rrbracket,[\hbarq\pi,-]\right)
\]
is a sheaf of complete homotopy-$\Etwo$ algebras.  Dunn additivity gives its locally constant factorization realization in two oriented topological directions.  The holomorphic sheaf variable retains the dependence on $S$.

\begin{theorem}[Algebraic factorization realization]
The homotopy-$\Etwo$ sheaf $\mathcal P_S$ determines a complete mixed holomorphic--topological factorization algebra
\[
 \Acal^{\mathrm{alg}}_{S,\mathrm{bulk}}
\quad\text{on}\quad
 \R^2_{\mathrm{top}}\times S.
\]
Its universal boundary category is $\Perf(\mathcal W_{S,\hbar})$, and the bulk--boundary morphism is a local equivalence
\[
 \Acal^{\mathrm{alg}}_{S,\mathrm{bulk}}
 \xrightarrow{\sim}
 Z^{\mathrm{der}}_{\mathrm{Mor}}
 \bigl(\Perf(\mathcal W_{S,\hbar})\bigr).
\]
On a formal Darboux chart this map is the symplectic polyvector-to-Weyl Hochschild comparison.  Comparing it with the pointed Hamiltonian centre additionally requires a map between those two polyvector complexes.
\end{theorem}
\begin{proof}
The factorization envelope of a complete homotopy-$\Etwo$ algebra is locally constant on oriented two-disks.  Sheafwise application to $\mathcal P_S$ gives the first assertion.  The Fedosov algebroid quantizes the symplectic Poisson tensor on $S$.  Globalized formality identifies its Hochschild cochains with $\mathcal P_S$, and Morita invariance identifies Hochschild cochains with the centre of the perfect-module category.  The construction is local on $S$ and therefore respects factorization descent.
\end{proof}

\section{BV realization}
The cyclic Lie sheaf of formal Hamiltonians has the shifted-cotangent action
\[
 S_{\mathrm{BF}}(A,B)
 =\int_{\R^2\times S}
 \left\langle B,(\dd+\dbar)A+\frac12[A,A]\right\rangle.
\]
A comparison of its classical observables with $\Acal^{\mathrm{alg}}_{S,\mathrm{bulk}}$ must specify a map from Hamiltonian Chevalley cochains to the symplectic polyvector complex.  The finite Lie algebra identification alone does not supply that map.  A compatible gauge fixing, extension of graph weights to compactified configurations, and solution of the quantum master equation produce a perturbative factorization realization of the same algebraic bulk--boundary map.  The finite jet zero modes carry the exact quantum measure constructed above.

\part{Matrix branes and trace geometry}

\chapter{The derived commuting stack}
Let $X,Y\in\mathfrak{gl}_N$ and let $\psi\in\mathfrak{gl}_N[1]$.  The Koszul differential
\[
 QX=QY=0,
 \qquad
 Q\psi=[X,Y]
\]
resolves the commuting equation.  The simultaneous conjugation action gives the boundary cdga
\[
 A_N=C^*_{\CE}\!\left(
 \mathfrak{gl}_N,
 \Sym(\mathfrak{gl}_N\oplus\mathfrak{gl}_N\oplus\mathfrak{gl}_N[1])
 \right).
\]

\begin{proposition}[Representation of the universal brane]
Substitution of matrices defines a Lie morphism
\[
 \rho_N:\h\longrightarrow\Der(A_N),
 \qquad
 f\longmapsto\delta_f,
\]
which integrates to an algebra morphism
\[
 \widehat U_\hbar(\h)\longrightarrow\End(A_N)\llbracket\hbarq\rrbracket,
 \qquad f\longmapsto\hbarq\,\delta_f.
\]
Thus the $N$-brane stack is a representation of the universal boundary algebra.
\end{proposition}
\begin{proof}
For a polynomial Hamiltonian $f$, cyclic differentiation gives the matrix vector field
\[
 \delta_f(X)=-\partial_Y^{\mathrm{cyc}}f(X,Y),
 \qquad
 \delta_f(Y)=\partial_X^{\mathrm{cyc}}f(X,Y).
\]
The cyclic Euler identity implies $\delta_f([X,Y])=0$.  Hence $\delta_f$ commutes with the Koszul differential $Q\psi=[X,Y]$.  The same formulas are equivariant for simultaneous conjugation, so $\delta_f$ descends to a derivation of the Chevalley--Eilenberg quotient $A_N$.  The trace-moment-map calculation gives
\[
 [\delta_f,\delta_g]=\delta_{\{f,g\}},
\]
first for polynomial Hamiltonians and then continuously for the formal completion.  This proves that $\rho_N$ is a continuous Lie morphism.  The assignment $f\mapsto\hbarq\delta_f$ satisfies
\[
 [\hbarq\delta_f,\hbarq\delta_g]
 =\hbarq^2\delta_{\{f,g\}}
 =\hbarq\,(\hbarq\delta_{\{f,g\}}),
\]
which is the Rees relation.  The universal property of the completed Rees enveloping algebra gives the displayed morphism.  It is the universal boundary action evaluated in the $N$-dimensional Chan--Paton representation.
\end{proof}

\chapter{Necklace Hamiltonians}
Let
\[
 F=\C\angles{x,y},
 \qquad
 F_{\mathrm{cyc}}=F/[F,F].
\]
For a cyclic word $f$, define cyclic derivatives by cutting the word at every occurrence of the selected letter.  Put
\[
 \delta_f(x)=-\partial_y^{\mathrm{cyc}}f,
 \qquad
 \delta_f(y)=\partial_x^{\mathrm{cyc}}f.
\]

\begin{theorem}[Necklace Lie algebra]
The bracket
\[
 \{f,g\}_{\mathrm{neck}}
 =\left[
 \partial_x^{\mathrm{cyc}}f\,
 \partial_y^{\mathrm{cyc}}g-
 \partial_y^{\mathrm{cyc}}f\,
 \partial_x^{\mathrm{cyc}}g
 \right]_{\mathrm{cyc}}
\]
makes $F_{\mathrm{cyc}}/\C$ a Lie algebra.  The assignment $f\mapsto\delta_f$ is a Lie morphism.  It preserves the moment map:
\[
 \delta_f([x,y])=0.
\]
\end{theorem}
\begin{proof}
The cyclic Euler identity is
\[
 [x,\partial_x^{\mathrm{cyc}}f]+[y,\partial_y^{\mathrm{cyc}}f]=0.
\]
It gives the moment-map statement.  The canonical double Poisson bracket on the free algebra descends to the displayed necklace bracket; its Jacobi identity gives the first assertion.  Evaluation on generators gives $[\delta_f,\delta_g]=\delta_{\{f,g\}_{\mathrm{neck}}}$.
\end{proof}

\begin{theorem}[Trace moment map]
For
\[
 J_N(f)=\Tr f(X,Y),
\]
one has
\[
 \{J_N(f),J_N(g)\}_{\mathcal M_N}
 =J_N\!\left(\{f,g\}_{\mathrm{neck}}\right),
\]
where $\mathcal M_N=\mathfrak{gl}_N^2$ carries $\Tr(\dd X\wedge\dd Y)$.
\end{theorem}
\begin{proof}
The differential of a trace is the matrix evaluation of the cyclic derivative.  Contracting the two differentials with the inverse trace symplectic form gives the necklace formula.
\end{proof}

\chapter{Stable traces}
On the semisimple commuting locus write the joint eigenvalues as $(x_i,y_i)$.  Define
\[
 p_{ab}^{(N)}=\sum_{i=1}^{N}x_i^ay_i^b.
\]

\begin{theorem}[Stable multisymmetric Poisson algebra]
The inverse stable trace algebra is the completed polynomial algebra
\[
 \Pcal_\infty=
 \C\llbracket p_{ab}\mid a,b\ge0,\ a+b>0\rrbracket.
\]
Its Poisson bracket is
\[
 \{p_{ab},p_{cd}\}
 =(ad-bc)p_{a+c-1,b+d-1}.
\]
\end{theorem}
\begin{proof}
The stable fundamental theorem for multisymmetric functions gives polynomial generators $p_{ab}$.  The bracket follows from $\{x_i,y_j\}=\delta_{ij}$ and the Leibniz rule.
\end{proof}

\begin{remark}
The stable trace algebra is a quotient of the universal closed algebra after representation and scalarization.  The universal centre retains Hochschild classes, stabilizer directions, and higher central operations that vanish under the trace.
\end{remark}

\chapter{The two-brane cone}
Separate center-of-mass and relative coordinates.  Put
\[
 u=x_1-x_2,
 \qquad
 v=y_1-y_2,
 \qquad
 E=\frac12u^2,
 \quad
 F=-\frac12v^2,
 \quad
 H=-uv.
\]
Then
\[
 H^2+4EF=0.
\]

\begin{proposition}[The $A_1$ Kirillov--Kostant cone]
With $\{u,v\}=2$,
\[
 \{H,E\}=4E,
 \qquad
 \{H,F\}=-4F,
 \qquad
 \{E,F\}=2H.
\]
After the evident rescaling this is the $\mathfrak{sl}_2$ Lie--Poisson bracket on the nilpotent cone.
\end{proposition}
\begin{proof}
The three identities follow by direct differentiation.  The quadratic equation is the determinant-zero equation for the traceless matrix
\[
 \begin{pmatrix}H&2E\\2F&-H\end{pmatrix}.
\]
\end{proof}

\part{Wheel corrections and exact quantum structure}

\chapter{The Duflo wheel}
First set the coefficient parameter to $1$ in the polynomial Rees algebra.
For the resulting ordinary $\mathfrak{sl}_2$ enveloping algebra, take generators
\[
 [h,e]=2e,
 \qquad
 [h,f]=-2f,
 \qquad
 [e,f]=h.
\]
The invariant polynomial and the quadratic Casimir are
\[
 q=2ef+\frac12h^2,
 \qquad
 \Omega=ef+fe+\frac12h^2.
\]

\begin{calculation}[The zero-point shift]
The Duflo map gives
\[
 \operatorname{Duf}(q)=\Omega+\frac12.
\]
On the irreducible highest-weight-$n$ module,
\[
 \left(\Omega+\frac12\right)\big|_{V_n}
 =\frac{(n+1)^2}{2}\id.
\]
\end{calculation}
\begin{proof}
Write $x=ue+vf+wh$.  Direct adjoint matrices give
$\Tr(\operatorname{ad}_x^2)=8uv+8w^2$.
Thus the quadratic term of $j^{1/2}(\partial)$ is
$(\partial_e\partial_f+\partial_h^2)/6$.
Applied to $q=2ef+h^2/2$, it gives $(2+1)/6=1/2$.
Higher differential terms annihilate this quadratic polynomial.
Symmetrization sends $q$ to $\Omega$.
On a highest-weight vector, $ef$ acts by $n$, $fe$ by zero,
and $h^2/2$ by $n^2/2$.
Centrality gives the eigenvalue on the entire irreducible module.
\end{proof}

Restoring the Rees parameter gives $\operatorname{Duf}_\hbar(q)=\Omega+\hbarq^2/2$.  The polynomial calculation permits specialization at $1$; arbitrary formal power series do not admit that specialization.  Higher wheels assemble into $j^{1/2}$ on higher-degree invariants.

\chapter{The local centre with comparison data}
Assume a compatible jet comparison and a complete homotopy-$\Etwo$
refinement as in Theorem~\ref{thm:vol2-compatible-jets}.
Put
\[
 \mathcal A=\widehat U_\hbar(\h),\qquad
 \Acal_{\mathrm{bulk}}^{\mathrm{alg}}
 =\operatorname{Fact}_{\Etwo}(\mathcal P_\infty).
\]
The factorization envelope uses the two oriented topological directions.
The formal holomorphic coefficient remains at the Darboux point.
Let $\Ccal_\partial$ be the category of perfect complete $\mathcal A$-modules
in the chosen continuous $R$-linear Morita theory.
Assume its centre is computed by
$C^\bullet_{\HH,\mathrm{cont}}(\mathcal A/R,\mathcal A)$.

\begin{theorem}[Local derived-centre comparison]
These data give an equivalence of complete homotopy-$\Etwo$ factorization algebras
\[
 \Acal_{\mathrm{bulk}}^{\mathrm{alg}}
 \xrightarrow{\sim}
 Z^{\mathrm{der}}_{\mathrm{Mor},R}(\Ccal_\partial).
\]
At the formal insertion point the map is the selected refinement of $\Phi_\infty$.
Every continuous boundary module has the induced derived centre action.
\end{theorem}
\begin{proof}
The compatible jet theorem gives the complete homotopy-$\Etwo$ comparison.
The factorization envelope preserves this equivalence.
The assumed continuous Morita model identifies its target with the centre.
A centre element is a derived natural endomorphism of the identity functor,
so evaluation gives the action on each boundary module.
\end{proof}

\begin{corollary}[Global K3 comparison with descent data]
Let $S$ be a K3 surface with holomorphic symplectic form.
Suppose the local comparisons extend to a sheafwise equivalence with the
relative Hochschild centre of $\Perf(\mathcal W_{S,\hbar})$.
Suppose further that the selected cyclic and mixed refinements extend
and satisfy descent in the same completed category.
Then derived global sections give the corresponding global centre,
cyclic, and open--closed comparisons.
\end{corollary}
\begin{proof}
Derived global sections preserve equivalences of sheaves.
Apply them to the supplied sheafwise maps and their coherence homotopies.
\end{proof}

\chapter{Cyclic classes, stable traces, and Lie homology}
A trace of the unit records the representation dimension.
For $B=\C$, one has $\operatorname{Tr}_N(1)=N$, so ordinary traces
cannot define a rank-compatible map on all of $HC_0(B)$.
A scalar-reduced trace removes this dependence after choosing an augmentation.

Let $B$ be a unital associative $\C$-algebra.
The affine representation scheme $\Rep_NB$ represents unital algebra maps
$B\to\operatorname{Mat}_N(D)$ for commutative $\C$-algebras $D$.
Put
\[
 \mathcal R_N=\Ocal(\Rep_NB)^{GL_N},\qquad
 \mathfrak c_B=HC_0(B)=B/[B,B],
\]
where $[B,B]$ denotes the linear span of commutators.
The ordinary trace is the linear map
\[
 \operatorname{Tr}_N:\mathfrak c_B\longrightarrow\mathcal R_N,
 \qquad
 \operatorname{Tr}_N([b])(\rho)=\operatorname{Tr}(\rho(b)).
\]
Cyclicity makes it independent of the representative, and conjugation
invariance puts its image in $\mathcal R_N$.

For Lie and Poisson compatibility, specify a Lie bracket on $\mathfrak c_B$,
a Poisson bracket on each $\mathcal R_N$, and the identities
\begin{equation}\label{eq:vol2-trace-bracket-datum}
 \{\operatorname{Tr}_N\beta,\operatorname{Tr}_N\gamma\}_N
 =\operatorname{Tr}_N[\beta,\gamma]_{\mathrm{cyc}}
 \qquad(\beta,\gamma\in\mathfrak c_B).
\end{equation}
A Hamiltonian comparison also specifies a Lie map
$j:\h\to\mathfrak c_B$ for this bracket.
These are additional data on $B$.
The preceding necklace bracket and trace moment map supply the bracket
and trace data for $B=\C\angles{x,y}$ with its canonical symplectic bracket.
The map $j$ remains an additional comparison.
An arbitrary unital algebra does not carry these structures by definition.

For rank stabilization, choose an augmentation $\epsilon:B\to\C$.
It defines the rank-stabilization morphism
\[
 s_N:\Rep_NB\longrightarrow\Rep_{N+1}B,
 \qquad \rho\longmapsto\rho\oplus\epsilon.
\]
Restriction gives $s_N^*:\mathcal R_{N+1}\to\mathcal R_N$.
Define the reduced cyclic space and trace by
\[
 \overline{\mathfrak c}_B=B/(\C1+[B,B]),
 \qquad
 t_N(\overline b)(\rho)
 =\operatorname{Tr}(\rho(b))-N\epsilon(b).
\]
Let $\mathfrak m_N\subset\mathcal R_N$ be the ideal of functions
vanishing at the representation $b\mapsto\epsilon(b)I_N$.
The completed invariant rings and stable target are
\[
 \widehat{\mathcal R}_N
 =\varprojlim_{r\ge1}\mathcal R_N/\mathfrak m_N^r,
 \qquad
 \mathcal R_\infty^{\wedge}
 =\varprojlim_N\widehat{\mathcal R}_N.
\]
Use the inverse-limit topologies.
Write $J$ for the augmentation ideal in $\Sym_\C\overline{\mathfrak c}_B$
and put
\[
 \widehat\Sym\overline{\mathfrak c}_B
 =\varprojlim_{r\ge1}
   \Sym_\C\overline{\mathfrak c}_B/J^r.
\]

The stable matrix Lie algebra is
$\mathfrak{gl}_\infty B=\varinjlim_N\mathfrak{gl}_N B$
under block-zero Lie embeddings.
For its homology, use trivial coefficients $\C$.
Let $V_B=\bigoplus_{n\ge1}HC_{n-1}(B)$, with $HC_{n-1}(B)$ placed
in homological degree $n$.
The notation $\Sym_{\mathrm{gr}}V_B$ denotes the free graded-commutative
algebra, with the Koszul sign rule.

\begin{theorem}[Cyclic traces and stable Lie homology]
Assume the compatible jet comparison and its homotopy Gerstenhaber refinement.
For the augmented algebra and completion above, there are maps
\[
 \mathcal P_\infty
 \xrightarrow{\ \Phi_\infty\ }
 C^*_{\HH,\mathrm{cont}}(\mathcal A/R,\mathcal A),
 \qquad
 \operatorname{Tr}_N:\mathfrak c_B\longrightarrow\mathcal R_N,
\]
\[
 \widehat\Sym\overline{\mathfrak c}_B
 \longrightarrow\mathcal R_\infty^{\wedge},
 \qquad
 \Sym_{\mathrm{gr}}V_B
 \xrightarrow{\sim}
 H_*^{\CE}(\mathfrak{gl}_\infty B;\C).
\]
The first preserves the selected homotopy Gerstenhaber structure.
The second is linear and, under the bracket data
\eqref{eq:vol2-trace-bracket-datum}, is a Lie map.
Under those data, its composite with $j$ preserves the Hamiltonian bracket.
The third is a continuous homomorphism of commutative algebras,
and the fourth is an isomorphism of graded Hopf algebras.
These maps are natural for morphisms preserving the stated data.
\end{theorem}
\begin{proof}
The first arrow is the chosen comparison of
Theorem~\ref{thm:vol2-compatible-jets}.
Its operadic refinement gives the asserted Gerstenhaber compatibility.
The trace formula above constructs the second arrow.
Equation~\eqref{eq:vol2-trace-bracket-datum} supplies its Lie compatibility,
and the Lie property of $j$ gives the composite assertion.

The augmentation annihilates commutators.
Both terms of $t_N$ vanish on a commutator, and their values on a scalar
multiple of the unit cancel.
Thus $t_N$ is well defined on $\overline{\mathfrak c}_B$.
At the trivial representation its value is zero, so its image lies in $\mathfrak m_N$.
Moreover,
\[
\begin{split}
 (s_N^*t_{N+1})(\overline b)(\rho)
 &=\operatorname{Tr}(\rho(b))+\epsilon(b)-(N+1)\epsilon(b)\\
 &=t_N(\overline b)(\rho).
\end{split}
\]
The point $N\epsilon$ maps to $(N+1)\epsilon$, so
$s_N^*(\mathfrak m_{N+1}^r)\subset\mathfrak m_N^r$.
Hence restriction defines the transition maps between the completed rings.
The symmetric-algebra extension of $t_N$ sends $J^r$ into $\mathfrak m_N^r$.
It therefore induces a ring homomorphism on each corresponding quotient.
Taking the inverse limit in $r$ gives a continuous map to
$\widehat{\mathcal R}_N$.
The displayed stabilization identity makes these maps compatible in $N$.
Their inverse limit is the third arrow.

The diagonal Lie map
$\mathfrak{gl}_\infty B\to
 \mathfrak{gl}_\infty B\oplus\mathfrak{gl}_\infty B$
induces the homology coproduct.
Block sum in the opposite direction induces the homology product,
using the K\"unneth isomorphism over $\C$.
The Loday--Quillen--Tsygan theorem identifies primitive homology
in degree $n$ with $HC_{n-1}(B)$.\footnote{J.-L. Loday and D. Quillen,
\emph{Cyclic homology and the Lie algebra homology of matrices},
Commentarii Mathematici Helvetici \textbf{59} (1984), 565--591,
Theorem~6.2, pp.~584--586.}
It identifies the full stable homology with the free graded-commutative
Hopf algebra on these primitives, giving the fourth arrow.
The trace formulas, the augmentation-preserving rank maps, and stable
Lie homology are natural.
Naturality of the first arrow requires preservation of the selected jet comparison.
\end{proof}

Scalar reduction need not preserve the supplied Poisson bracket.
For example, for $B=\C[x,y]$ with $\{x,y\}=1$ and
$\epsilon(x)=\epsilon(y)=0$, traces on $N$ points satisfy
\[
 \{\operatorname{Tr}_Nx,\operatorname{Tr}_Ny\}=N,
 \qquad t_N(\overline1)=0.
\]
A reduced Poisson comparison consequently requires additional compatibility,
such as vanishing of $\epsilon$ on the selected cyclic brackets,
or a restriction to a Hamiltonian subalgebra with that property.
Only multiplicativity and continuity are asserted for the third arrow.
For $B=\C$, the reduced cyclic space is zero and that arrow is the identity on $\C$.

\chapter{Synthesis}
At each fixed Hamiltonian jet, the Rees enveloping algebra has relative
Hochschild complex equivalent to the Chevalley complex with differential
$\hbarq d_{\CE}$.
PBW antisymmetrization proves this over every $R_N$ and in the formal coefficient limit.
The Duflo theorem gives the separate multiplicative comparison on invariant functions.
Quadratic Moyal Hamiltonians generate the oscillator quotient of the enveloping algebra.
Their associative Casimir relation prevents the primitive coproduct from descending.

The Hamiltonian jet limit requires the compatible cochain systems of
Theorem~\ref{thm:vol2-compatible-jets}.
With its operadic refinement and continuous Morita model, the centre acts
on every supplied continuous boundary module $M$ through
\[
 \mathcal P_\infty
 \xrightarrow{\ \Phi_\infty\ }
 C^\bullet_{\HH,\mathrm{cont}}(\mathcal A/R,\mathcal A)
 \longrightarrow\RHom_{\mathcal A}(M,M).
\]
The last arrow evaluates a derived natural endomorphism at $M$.
A matrix-brane realization must construct such a module and compare its
centre action with the selected trace functions.
A commuting square with scalar representation invariants requires these
additional comparison maps.
Likewise, Fedosov descent concerns symplectic polyvectors and a Weyl algebroid.
Its relation to Hamiltonian Lie cochains requires the specified map between their carriers.

\part{Technical appendices to the preceding book}
\chapter{Signs in the shifted cotangent complex}
Let $e_a$ have degree zero and let $e^a[-1]$ have degree $-1$.  The Chevalley generators dual to them have degrees $1$ and $0$.  Thus
\[
 C^*_{\CE}(T^*[-1]\flie)
 =\Sym\bigl(\flie^\vee[-1]\oplus\flie\bigr).
\]
The degree-zero polynomial factor is $\Sym\flie=\Ocal(\flie^*)$.  The differential on a polynomial $p$ is
\[
 (d_{\CE}p)(x)=\ad_xp,
\]
which is the linear Poisson differential.

\chapter{A direct Jacobi check for Hamiltonian jets}
For monomials
\[
 f=z_1^az_2^b,
 \qquad
 g=z_1^cz_2^d,
\]
one has
\[
 \{f,g\}=(ad-bc)z_1^{a+c-1}z_2^{b+d-1}.
\]
Substitution of three monomials gives a scalar multiple of
\[
 \det((a,b),(c,d))\,\det((a+c-1,b+d-1),(e,f)) + \mathrm{cyclic}.
\]
Writing the three coefficients in terms of the $2\times2$ determinants of the exponent vectors shows pairwise cancellation.  The coordinate-free proof is $[\pi,\pi]_{\mathrm{Sch}}=0$ for $\pi=\partial_{z_1}\wedge\partial_{z_2}$.

\chapter{Bibliographic note}
The finite affine formality map and the separate symplectic globalization use Kontsevich formality and its Fedosov resolution \cite{KontsevichDQ,DolgushevCov,CalaqueDolgushevHalbout}.  The relative Hochschild model is the Cartan--Eilenberg comparison, with both Chevalley terms scaled by $\hbarq$.  Its finite-base resolution proof is Theorem~\ref{thm:vol2-relative-rees}.  Jet limits and operadic refinements require the additional data stated above.  The all-arity centre action is the Deligne--Tamarkin brace action \cite{McClureSmith,KontsevichSoibelmanDeligne}.  The global deformation-quantization category is the holomorphic Weyl algebroid.  The matrix trace and necklace formulas are the noncommutative symplectic moment-map construction.

\input{chapters/Volume_I_Ordered_Chiral_Geometry.tex}
\chapter{Enveloping bars and planar path sequences}
\label{ch:enveloping-audit}

The dimension of a bar chain space and the dimension of its homology are different invariants. For an enveloping algebra, the Chevalley comparison gives a Lie-theoretic model for the homology. This makes it possible to compare Motzkin and Riordan path counts with the ordered bar on a specified carrier.

Bar length and Chevalley degree are homological degrees: $H_n$ denotes degree-$n$ homology. The cohomology of the dual Chevalley complex is written $H^n$. For finite-dimensional weight pieces their dimensions agree.

\section{Chevalley comparison}

\begin{theorem}[Bar--Chevalley comparison]
\label{thm:bar-chevalley-spine}
Let $\g$ be a Lie algebra object in $\FactD(X)$ for the pointwise product,
$K$-flat, positively graded and complete, with every total weight admitting
finitely many decompositions into positive weights, and suppose internal PBW
holds for $U(\g)$. There is a natural quasi-isomorphism of conilpotent
coalgebras in $\FactD(X)$
\[
 \BarPerp\bigl(U(\g)\bigr)\ \xrightarrow{\ \sim\ }\ C_*(\g),
\]
compatible with the factorization coefficient and the PBW filtration.
\end{theorem}
\begin{proof}
Filter $U(\g)$ by PBW degree and both complexes by total word length. The
associated graded enveloping algebra is $\Sym(\g)$, and the bar of a symmetric
algebra is the suspended symmetric coalgebra on $\g$, which is the associated
graded of Chevalley chains. The comparison is an isomorphism on the first page;
finite decomposition gives a bounded exhaustive spectral sequence in each
weight, and completeness promotes it to a quasi-isomorphism. Every step uses
only internal tensor products and the bracket, so the map stays in $\FactD(X)$.
\end{proof}

\begin{calculation}[Finite-dimensional simple case]
\label{calc:simple-spine}
For a simple Lie algebra of rank $r$ with exponents $m_1,\dots,m_r$, the
Chevalley cohomology is the exterior algebra on primitive classes of degrees
$2m_i+1$. The corresponding homology Poincar\'e polynomials are
\[
 \BarPerp U(\mathfrak{sl}_2):\ 1+t^3,
 \qquad
 \BarPerp U(\mathfrak{sl}_3):\ 1+t^3+t^5+t^8,
 \qquad
 \BarPerp U(\mathfrak g_2):\ 1+t^3+t^{11}+t^{14}.
\]
The corresponding Betti sequences for $\mathfrak{sl}_2$ and $\mathfrak{sl}_3$ are
$1,0,0,1$ and $1,0,0,1,0,1,0,0,1$.
For a finite-dimensional Lie algebra, the Chevalley chain dimensions are
$\binom{\dim\g}{n}$.

For an ordinary augmented algebra $A$ over $\mathbb C$, let
$I=\ker(A\to\mathbb C)$.  Its normalized two-sided bar with trivial
coefficients has homological chain spaces
\[
 B_0(A)=\mathbb C,\qquad B_n(A)=I^{\otimes n}\quad(n\ge1).
\]
Here $n$ is bar length, independently of any internal weight or PBW degree.
If $I$ has finite dimension $d$, then $\dim B_n(A)=d^n$; in particular,
an augmentation ideal of dimension three gives the chain counts $3^n$.
In a cohomological convention that places $B_n$ in degree $-n$, its homology
is denoted $H^{-n}$ instead of $H_n$.

For $A=U(\mathfrak{sl}_2)$, PBW gives the basis
$e^a h^b f^c$, with $a,b,c\ge0$, where
$[h,e]=2e$, $[h,f]=-2f$, and $[e,f]=h$.
The augmentation ideal has all these basis elements except $1$.
Consequently $B_n(A)$ is countably infinite-dimensional for every $n\ge1$.
Its total PBW filtration has finite-dimensional associated graded pieces,
with generating function
\[
 \sum_{d\ge0}\dim\operatorname{gr}^{\mathrm{PBW}}_d B_n(A)\,z^d
 =\bigl((1-z)^{-3}-1\bigr)^n.
\]
In particular, $3^n$ is the dimension in the least possible total PBW degree
$d=n$. These least-degree dimensions are $3,9,27,81,243$ for
$n=1,2,3,4,5$; the full bar chain spaces are infinite-dimensional. The quotient of
indecomposables is zero: the perfectness of $\mathfrak{sl}_2$ implies $I/I^2=0$.

The Chevalley chain model is different: its degree-$n$ space is
$\Lambda^n\mathfrak{sl}_2$, of dimensions $1,3,3,1$ in degrees $0,1,2,3$.
With the convention $d(x\wedge y)=[x,y]$, the map
$d:\Lambda^2\mathfrak{sl}_2\to\mathfrak{sl}_2$ sends
\[
 e\wedge h\longmapsto-2e,\qquad
 e\wedge f\longmapsto h,\qquad
 h\wedge f\longmapsto-2f,
\]
and is an isomorphism.  The next differential vanishes, since
\[
 d(e\wedge h\wedge f)
 =-2e\wedge f-h\wedge h-2f\wedge e=0.
\]
Together with $d:\mathfrak{sl}_2\to\mathbb C$ equal to zero, this gives
Chevalley homology dimensions $1,0,0,1$.  The ordinary bar--Chevalley
comparison identifies these homology groups with those of $U(\mathfrak{sl}_2)$;
it does not identify the two chain spaces.

\end{calculation}

\section{The Virasoro row}

The transverse presentation of the Virasoro family is the enveloping algebra of
the positive Witt algebra
\[
 L_1=\bigoplus_{n\ge1}\mathbb C\,e_n,
 \qquad
 [e_i,e_j]=(j-i)e_{i+j},
\]
to which Theorem~\ref{thm:bar-chevalley-spine} applies: $L_1$ is positively
graded with one-dimensional weight spaces, so every weight has finitely many
decompositions.

\begin{theorem}[Pentagonal rigidity]
\label{thm:pentagonal-spine}
The bar homology of $U(L_1)$ satisfies
\[
 \dim H_n\bigl(\BarPerp U(L_1)\bigr)=2
 \qquad\text{for every }n\ge1,
\]
one dimension in each of the weights $(3n^2-n)/2$ and $(3n^2+n)/2$. Its graded
Euler character is
\begin{equation}
\label{eq:witt-euler-spine}
 \chi_q\bigl(\BarPerp U(L_1)\bigr)=\prod_{n\ge1}\bigl(1-q^n\bigr),
\end{equation}
and \eqref{eq:witt-euler-spine} is accounted for degree by degree with no
residual freedom.
\end{theorem}
\begin{proof}
By Theorem~\ref{thm:bar-chevalley-spine} the bar of $U(L_1)$ is quasi-isomorphic to $C_*(L_1)$.
Each weight subcomplex is finite-dimensional, so its homology has the same
dimensions as the weightwise dual cohomology. Goncharova's calculation gives
two dimensions in each degree $n\ge1$, in the weights $(3n^2\mp n)/2$. Identity \eqref{eq:witt-euler-spine} is the
denominator product applied to $L_1$, where every root space is one dimensional
and even, so $\smult(L_1)_n=1$. Euler's pentagonal number theorem expands the
product as $\sum_{m\in\Z}(-1)^mq^{m(3m-1)/2}$, whose support is exactly
$\{(3n^2\pm n)/2\}$ with coefficient $(-1)^n$ at each of $(3n^2-n)/2$ and
$(3n^2+n)/2$. Comparing with $\sum_n(-1)^n\dim H_{n,w}$, each weight of the Euler
character is saturated by a single class in the degree Goncharova predicts, so
no further classes and no cancellation are available.
\end{proof}

For $n=1,2,3,4$, the nonzero homology groups occupy the weight pairs
$(1,2)$, $(5,7)$, $(12,15)$, $(22,26)$, one dimension in each.
The corresponding nonzero coefficients of $\prod_{n\ge1}(1-q^n)$ occur at
$0,1,2,5,7,12,15,22,26$ with signs $+,-,-,+,+,-,-,+,+$.

\section{Planar path sequences}

A Motzkin path is a lattice path from $(0,0)$ to $(n,0)$, staying at
nonnegative height, with steps $(1,1)$, $(1,0)$, and $(1,-1)$.
Let $M(n)$ count all such paths of length $n$, and let $R(n)$ count those with
no horizontal step at height zero.  Both counts include the empty path:
$M(0)=R(0)=1$, while $M(1)=1$ and $R(1)=0$.
For $n\ge1$ consider the two sequences
\[
 a_n=M(n+1)-M(n)=1,2,5,12,30,76,\ldots,
 \qquad
 b_n=R(n+3)=3,6,15,36,91,\ldots.
\]

\begin{proposition}[Motzkin and Riordan identities]
\label{prop:one-sequence-spine}
For every $n\ge0$,
\[
 M(n)=R(n)+R(n+1),\qquad
 M(n+1)-M(n)=R(n+2)-R(n).
\]
The ordinary generating functions of the two path counts, and hence those
of $a_n$ and $b_n$, belong to
$\mathbb Q(z,\sqrt{1-2z-3z^2})$.
\end{proposition}
\begin{proof}
Work in the formal power-series ring $\mathbb Q[[z]]$, and write
$M(z)=\sum_{n\ge0}M(n)z^n$ and $R(z)=\sum_{n\ge0}R(n)z^n$.
A nonempty Motzkin path either starts with a horizontal step followed by
an arbitrary Motzkin path, or has a unique first-return decomposition
$U P D Q$, where $U$ and $D$ are the up and down steps,
$P$ and $Q$ are arbitrary Motzkin paths, and $P$ is drawn
one unit above the horizontal axis.  Thus
\[
 M=1+zM+z^2M^2.
\]
A nonempty path counted by $R$ has no initial horizontal step.  Its first
return similarly has the form $U P D Q$, with $P$ arbitrary and $Q$ counted
by $R$.  Consequently
\[
 R=1+z^2MR,\qquad R=(1-z^2M)^{-1}.
\]
The inverses here exist because their denominators have constant term one.
Using $z^2M^2=M-1-zM$, compute
\[
 (1+zM)(1-z^2M)
 =1+zM-z^2M-z^3M^2=1+z.
\]
It follows that $(1+z)R=1+zM$.  Comparing coefficients of $z^{n+1}$ gives
$R(n+1)+R(n)=M(n)$ for all $n\ge0$.  Subtracting consecutive identities
gives the difference formula.

Solving the quadratic equation for the formal series with constant term
one gives
\[
 M(z)=\frac{1-z-\sqrt{1-2z-3z^2}}{2z^2},\qquad
 R(z)=\frac{1+zM(z)}{1+z},
\]
where the square root has constant term one.  The apparent pole in the
formula for $M$ is removable as a formal series.  Finite index shifts and
differences preserve the indicated quadratic function field, proving the
last assertion.  These identities are consequences solely of the path
decompositions.
\end{proof}

\begin{remark}[Low-degree homology]
Theorem~\ref{thm:pentagonal-spine} gives homology dimensions $2,2,2$ in degrees $1,2,3$, whereas $a_1,a_2,a_3$ are $1,2,5$. These triples agree only in degree two. The degree-one discrepancy already excludes equality of the full homology dimension sequences.
\end{remark}

\begin{theorem}[Enveloping bars and planar path counts]
\label{thm:carrier-spine}
Let $B_\bullet(A)$ be the ordinary normalized bar over $\mathbb C$, with
trivial coefficients and homological degree given by bar length.  With
$a_n=M(n+1)-M(n)$ and $b_n=R(n+3)$ indexed by $n\ge1$ as above,
\[
 \dim H_1 B_\bullet(U(\mathfrak{sl}_2))=0\ne b_1=3,
 \qquad
 \dim H_1 B_\bullet(U(L_1))=2\ne a_1=1,
\]
where $L_1=\bigoplus_{m\ge1}\mathbb C e_m$ and
$[e_i,e_j]=(j-i)e_{i+j}$.
Hence neither sequence is the bar homology dimension sequence of the
corresponding enveloping algebra.  Neither is its full bar chain dimension
sequence: in both cases $B_n$ is countably infinite-dimensional for every
$n\ge1$.
\end{theorem}
\begin{proof}
For any augmented algebra $A$ with augmentation ideal $I$, the bar
differential out of $B_1=I$ is zero and the image from $B_2=I\otimes I$ is
$I^2$.  Therefore
\[
 H_1B_\bullet(A)=I/I^2.
\]
For an ordinary Lie algebra $\mathfrak g$, the inclusion of generators
induces an isomorphism
\[
 \mathfrak g/[\mathfrak g,\mathfrak g]\ \xrightarrow{\ \sim\ }
 I/I^2,\qquad I=\ker(U(\mathfrak g)\to\mathbb C).
\]
Indeed the generators span $I/I^2$, and
$[x,y]=xy-yx$ vanishes there.  For injectivity, equip
$\mathbb C\oplus\mathfrak g/[\mathfrak g,\mathfrak g]$ with the square-zero
product on its second summand.  The quotient map from $\mathfrak g$ to that
summand is a Lie map into the commutator Lie algebra of this associative
algebra.  The universal property of $U(\mathfrak g)$ extends it to an
augmented algebra map, which kills $I^2$ and provides an inverse on the
displayed quotient.

The brackets $[h,e]=2e$, $[h,f]=-2f$, and $[e,f]=h$ show that
$\mathfrak{sl}_2$ equals its commutator algebra.  For $L_1$, every nonzero
bracket has weight at least three.  Conversely, for $m\ge3$,
\[
 e_m=\frac{1}{m-2}[e_1,e_{m-1}].
\]
Thus $[L_1,L_1]=\bigoplus_{m\ge3}\mathbb C e_m$, and its abelianization
has basis the classes of $e_1,e_2$, in weights one and two.  This proves
both degree-one assertions without using any higher-degree calculation.

PBW gives a countably infinite basis of the augmentation ideal of either
enveloping algebra.  In the $L_1$ case the basis consists of nonempty
ordered monomials in the countable set $e_1,e_2,\ldots$, each monomial
having finite length.  Finite tensor powers of either augmentation ideal
are therefore countably infinite-dimensional.  The path counts are finite
in every degree, so they cannot count these full chain spaces.
\end{proof}

In particular, a complex whose homology dimensions are $b_n$ cannot be
quasi-isomorphic to $B_\bullet(U(\mathfrak{sl}_2))$, and one whose homology
dimensions are $a_n$ cannot be quasi-isomorphic to $B_\bullet(U(L_1))$,
with the stated degree convention.  A discrepancy of chain dimensions
alone would not imply this conclusion; the degree-one homology discrepancy
does.

The weight grading gives a different, finite comparison.  An ordered PBW
monomial of weight $w$ specifies a partition of $w$, so for $n\ge1$,
\[
 \sum_{w\ge0}\dim B_n(U(L_1))_w q^w
 =\left(\prod_{m\ge1}(1-q^m)^{-1}-1\right)^n.
\]
Each coefficient is finite.  Thus a weightwise dimension, a dimension
after a specified weight cutoff, and an ungraded dimension are distinct
invariants.  In a weight completion, dimensions must likewise be interpreted
weightwise; the countable algebraic dimensions above refer to the ordinary
direct-sum enveloping algebras.

\begin{remark}[Chiral coefficients and transverse multiplication]
Independently of these scalar calculations, a Virasoro chiral coefficient does not determine the transverse bar: a fixed coefficient may be combined with free, square-zero, enveloping, or quantum-plane transverse algebras whose bars differ, while the coefficient and the central charge are unchanged. A chain-level comparison with an ordered bar therefore also requires a specified transverse multiplication and a quasi-isomorphism from the proposed complex to that bar.
\end{remark}

\chapter{The obstruction constitution}
\label{ch:obstructions}

Each statement below names what is \emph{not} determined, and therefore what
must be supplied. Without them the preceding theorems admit false corollaries.

\begin{proposition}[Where a mixed distributive law can live]
\label{prop:lambda-sector-spine}
The sector on which a mixed distributive law between $\tensor^!$ and $\starCh$
is geometrically defined is the common-decomposition sector, the image of the
aligned idempotent $e_{\mathrm{al}}$. Off that sector the two products admit no
interchange, only the aligned retraction, so no distributive law can be
written.
\end{proposition}
\begin{proof}
A distributive law requires a natural transformation interchanging the two
products in a fixed direction, subject to the coherence conditions of each. The
canonical comparison in the additive model is the aligned inclusion $\iota$ with
retraction $\pi$, and $\pi\iota=\id$ while $\iota\pi=e_{\mathrm{al}}\ne\id$ in
general. A natural transformation subject to both coherence conditions
therefore exists only where $e_{\mathrm{al}}$ acts as the identity.
\end{proof}

This locates a mixed law without constructing one. No example in the present
corpus exhibits one; Zamolodchikov--Faddeev algebras are the natural candidate,
since their spectral variable belongs to the chiral direction while their word
belongs to the transverse direction.

\begin{theorem}[No quantum group from dimension loss]
\label{nogo:dimension-loss-spine}
The dimension of the kernel of $V^{\otimes n}\to\Sym^nV$ determines no
$R$-matrix, associator, or quantum group. That kernel carries no
multiplication, coproduct, monodromy, or coherence equation.
\end{theorem}
\begin{proof}
The kernel is a representation of $\Sigma_n$ and nothing more. Two algebras
with the same underlying graded vector space, hence the same kernel dimensions
in every arity, may have different multiplications and inequivalent
representation categories; a free algebra and a quadratic quotient with the
same Hilbert series is an example. A quantum-group structure is a coproduct
together with coherence data, none of which is a function of a dimension.
\end{proof}

\begin{theorem}[No modular object from topology alone]
The identity $\partial^2=0$ on $\overline{\mathcal M}_{g,n}$ assigns no
coefficient operations to boundary strata. A modular master element exists only
after cyclic contractions and clutching-compatible coefficient maps have been
supplied.
\end{theorem}

\begin{theorem}[No bulk from a boundary centre]
The derived centre is terminal among algebraic central actions on the boundary
algebra. It does not detect a decoupled bulk sector, so a physical bulk is not
determined by the boundary chiral algebra unless a bulk-to-centre comparison is
supplied and proved to be an equivalence.
\end{theorem}
\begin{proof}
Terminality is the universal property of the centre. It is a statement about
the boundary category and is insensitive to any sector acting trivially on it,
so a bulk with trivial boundary action is invisible.
\end{proof}

\begin{theorem}[No classification by central charge]
Central charge determines the Virasoro anomaly coefficient. It does not
determine the chiral operation, the representation category, the braid
monodromy, or the higher-genus conformal blocks. Distinct holomorphic theories
at equal central charge exist.
\end{theorem}

\begin{theorem}[No identification of anomalies of different types]
\label{nogo:anomalies-spine}
A current-algebra level, a Virasoro central charge, a determinant-line
curvature, a BRST ghost-number anomaly, and a Borcherds-product weight are
invariants in different theories, living in different deformation complexes.
They may be related by a theorem in a particular family. They are not one
universal scalar, and no identity among them follows from numerical agreement.
\end{theorem}
\begin{proof}
Each is a class in a distinct complex: an Atiyah-algebra extension class, a
Virasoro cocycle, the curvature of a determinant line, a class in $H^1$ or
$H^2$ of a BRST convolution complex, and a weight in the character lattice of
an automorphic form. A map between two of these complexes is additional data;
in its absence, equality of numbers is a coincidence of values, not an
identification of classes.
\end{proof}

\begin{theorem}[No unnamed comparison of bars]
An asserted equality $\BarPerp(A)=\BarAssCh(A)$ has no meaning without a
functor carrying the pointwise multiplication problem to the chiral-convolution
multiplication problem. The two have different underlying gradings, counting
interval subdivisions and collision trees respectively. Agreement of Euler
characteristics is not a comparison map.
\end{theorem}

\begin{theorem}[A scalar sequence is not homology data]
\label{nogo:scalar-spine}
A sequence of integers becomes bar homology only after a chain complex, a
differential, a grading convention, and a rank computation or quasi-isomorphism
are specified. Matching initial values, discriminants, or growth rates
constructs no such map.
\end{theorem}
\begin{proof}
Chain-space dimensions and homology dimensions differ by the ranks of the
differentials, which are not functions of the chain dimensions.
Calculation~\ref{calc:simple-spine} exhibits an algebra whose bar chain
dimensions are $3,9,27,81,243$ and whose bar homology dimensions are $1,0,0,1$;
Proposition~\ref{prop:one-sequence-spine} exhibits two sequences with the same
discriminant and growth rate that are transforms of one another and of neither
homology.
\end{proof}

Theorem~\ref{nogo:anomalies-spine} together with the central-charge statement
retracts any programme treating a tuple of scalars drawn from these sources as
a single classifying invariant, and any trace identity asserted among them
without a named comparison morphism.

\cleardoublepage
\phantomsection
\addcontentsline{toc}{part}{The Igusa Square Root}
\chapter*{Book entrance: The Igusa Square Root}
\addcontentsline{toc}{chapter}{Book entrance: The Igusa Square Root}
\begin{summarybox}
The Jacobi seed, Borcherds denominator, root superalgebra, local--wrapped currents, quantum group, Hall model, and Fredholm realization.
\end{summarybox}
\part{The automorphic determinant}

\chapter{The protected Jacobi input}
\section{The K3 elliptic genus}
Let $S$ be a K3 surface. The right-moving supersymmetric index of a unitary $\mathcal N=(4,4)$ sigma model is the elliptic genus
\[
 Z_{K3}(\tau,z)=2\phiweak(\tau,z).
\]
The normalized weak Jacobi form has weight zero, index one, and Fourier expansion
\[
 \phiweak(\tau,z)=\sum_{n\ge0}\sum_{\ell\in\Z}f(n,\ell)q^nr^\ell.
\]
In Eichler--Zagier normalization,
\[
 \phiweak(\tau,z)=4\sum_{a=2}^{4}
 \frac{\vartheta_a(\tau,z)^2}{\vartheta_a(\tau,0)^2}
 \in J^{\mathrm w}_{0,1}
\]
\cite{EZ}. The coefficient depends on the discriminant $4n-\ell^2$ and on $\ell$ modulo $2$.

\begin{theorem}[Exact low Fourier coefficients]
The first four $q$-rows are
\[
\begin{aligned}
\phiweak={}&r^{-1}+10+r\\
&+q\bigl(10r^{-2}-64r^{-1}+108-64r+10r^2\bigr)\\
&+q^2\bigl(r^{-3}+108r^{-2}-513r^{-1}+808-513r+108r^2+r^3\bigr)\\
&+q^3\bigl(-64r^{-3}+808r^{-2}-2752r^{-1}+4016\\
&\hspace{7.5em}{}-2752r+808r^2-64r^3\bigr)+O(q^4).
\end{aligned}
\]
In particular,
\[
 f(0,0)=10,
 \qquad \frac{f(0,0)}2=5.
\]
\end{theorem}
\begin{proof}
Insert the theta-series expansions into the preceding quotient formula. Work with $t=q^{1/8}$ and $x=r^{1/2}$; each quotient is an exact formal power series in $t$ with Laurent-polynomial coefficients in $x$. The fractional powers cancel in the sum. Division through $t^{24}$ gives the displayed rows. The arithmetic calculation is reproduced in Appendix~\ref{app:theta-computation}.
\end{proof}

\section{The virtual half-index}
For every $(n,\ell)$ choose a finite-dimensional super vector space $U_{n,\ell}$ satisfying
\[
 \sdim U_{n,\ell}=f(n,\ell).
\]
The integer \(f(n,\ell)\) is its superdimension. In the ordinary
Grothendieck group of finite-dimensional super vector spaces, the
even and odd dimensions remain separate. The signed integer is their
image under superdimension. The determinant depends only on that image.

\begin{lemma}[Fock character and Berezinian]
Let $V=\bigoplus_{\gamma}V_\gamma$ be a locally finite graded super vector space. Then
\[
 \sch\bigl(\Sym(V_{\bar0})\tensor\bigwedge(V_{\bar1})\bigr)
 =\prod_\gamma(1-x^\gamma)^{-\sdim V_\gamma},
\]
and
\[
 \sdet_V(1-x)=\prod_\gamma(1-x^\gamma)^{\sdim V_\gamma}.
\]
\end{lemma}
\begin{proof}
An even one-dimensional summand contributes the geometric series $(1-x^\gamma)^{-1}$. An odd one-dimensional summand contributes $1-x^\gamma$. Multiplication over a homogeneous basis proves both identities.
\end{proof}

\chapter{The Borcherds product}
\section{The effective chamber}
Define
\[
\begin{aligned}
\Gammaeff={}&\set{(n,\ell,m)\in\Z^3:m>0,\ n\ge0}\\
&\cup\set{(n,\ell,0):n>0}
 \cup\set{(0,\ell,0):\ell<0}.
\end{aligned}
\]
The ordering is lexicographic in $(m,n,-\ell)$ and is precisely the ordering selected by the cusp chamber.
For $\gamma=(n,\ell,m)$ put
\[
 V_\gamma=U_{nm,\ell},
 \qquad \sdim V_\gamma=f(nm,\ell).
\]

\begin{definition}[Monic virtual determinant]
The monic determinant of the half-index is
\[
 \Dmon(Z)=q^{1/2}r^{1/2}s^{1/2}
 \sdet_V(1-x)
 =q^{1/2}r^{1/2}s^{1/2}
 \prod_{\Gammaeff}(1-q^nr^\ell s^m)^{f(nm,\ell)}.
\]
\end{definition}

\begin{theorem}[Borcherds--Gritsenko--Nikulin product]
The preceding product converges normally on the chosen cusp chamber and extends to a holomorphic genus-two modular form of weight five with a character of order two. It is the monic Igusa form $\Dmon$. The theta-normalized form is
\[
 \Dtheta=2^6\Dmon.
\]
\end{theorem}
\begin{proof}
The singular theta lift of the vector-valued form corresponding to $\phiweak$ produces the product, divisor, weight, and automorphic character. The weight is $f(0,0)/2=5$. The genus-two identification with the Igusa form and the product formula are the Gritsenko--Nikulin theorems \cite{Borcherds95,Borcherds98,GN97,GN98b}. The product of the ten even theta constants has coefficient $2^6$ at $q^{1/2}r^{1/2}s^{1/2}$, fixing the scalar.
\end{proof}

\section{The theta constant normalization}
Let $\mathfrak E$ be the ten even genus-two characteristics. Define
\[
 \Dtheta(Z)=\prod_{\mathfrak m\in\mathfrak E}\vartheta[\mathfrak m](Z,0).
\]
The cusp leading terms of the ten factors contribute six factors of $2$ and total exponent $(1/2,1/2,1/2)$. Thus
\[
 [q^{1/2}r^{1/2}s^{1/2}]\Dtheta=64.
\]
This coefficient selects a canonical scalar normalization of the automorphic section.

\begin{proposition}[Ten-factor leading-term calculation]
Group the ten even characteristics by the first half-characteristic $a\in(\Z/2)^2$. Their leading contributions in the chosen cusp chamber are
\[
\begin{array}{c|c|c|c}
a&\#&\text{scalar}&\text{monomial}\\ \hline
(0,0)&4&1&1\\
(1,0)&2&2&q^{1/8}\\
(0,1)&2&2&s^{1/8}\\
(1,1)&2&2&q^{1/8}r^{1/4}s^{1/8}.
\end{array}
\]
Multiplication gives
\[
 1^4\,2^2\,2^2\,2^2\,
 q^{1/2}r^{1/2}s^{1/2}
 =64q^{1/2}r^{1/2}s^{1/2}.
\]
\end{proposition}
\begin{proof}
Insert each characteristic into the genus-two theta series and retain the unique terms of least order in the lexicographic cusp valuation $(m,n,-\ell)$. For $a=(0,0)$ the zero lattice vector contributes one. For each nonzero $a$, the two shortest lattice representatives are exchanged by sign and contribute the scalar two. Their quadratic exponents are the entries in the final column. Summing the exponents and multiplying the six factors of two gives the formula.
\end{proof}

\begin{proposition}[Boundary factor of the Borcherds product]
The $m=0$ part of the monic product is
\[
 \prod_{n>0,\ell\in\Z}(1-q^nr^\ell)^{f(0,\ell)}
 \prod_{\ell<0}(1-r^\ell)^{f(0,\ell)}
 =(1-r^{-1})\prod_{n\ge1}
 (1-q^nr^{-1})(1-q^n)^{10}(1-q^nr).
\]
Consequently the first Fourier--Jacobi coefficient is
\[
 q^{1/2}r^{1/2}(1-r^{-1})
 \prod_{n\ge1}(1-q^nr^{-1})(1-q^n)^{10}(1-q^nr)
 =\eta(\tau)^9\vartheta_{11}(\tau,z),
\]
where the odd theta function is fixed by the exact convention
\[
 \vartheta_{11}(\tau,z)
 =-q^{1/8}r^{-1/2}
 \prod_{n\ge1}(1-q^{n-1}r)(1-q^nr^{-1})(1-q^n).
\]
\end{proposition}
\begin{proof}
The $q^0$ row of $\phiweak$ has exactly three nonzero coefficients:
$f(0,-1)=1$, $f(0,0)=10$, and $f(0,1)=1$. Substitution into the $m=0$ factors gives the first formula. The Jacobi triple product and $\eta(\tau)=q^{1/24}\prod_{n\ge1}(1-q^n)$ give the second, including its leading sign and monomial.
\end{proof}

\begin{corollary}[The scalar square]
The monic weight-ten form is
\[
 \ChiTen=\Dmon^2=2^{-12}\Dtheta^2.
\]
Its automorphic character is trivial.
\end{corollary}
\begin{proof}
The identity $\Dtheta=2^6\Dmon$ gives the scalar factor.  The character of $\Dmon$ has order two, so its square is trivial.
\end{proof}

\chapter{Divisor, character, and the meaning of the square root}
The diagonal Humbert divisor is
\[
 H_1=\set{Z\in\HHtwo:z=0}/\operatorname{Sp}_4(\Z).
\]

\begin{theorem}[Diagonal divisor]
The divisor of $\Dmon$ is the $\operatorname{Sp}_4(\Z)$-orbit of $H_1$, with multiplicity one. The form transforms by the Maass character $\nu_5$ of order two, and $\ChiTen=\Dmon^2$ is scalar-valued.
\end{theorem}
\begin{proof}
This is the defining diagonal-divisor property of the Igusa square root and the first row of the Cl\'ery--Gritsenko classification \cite{Igusa62,Igusa64,GC}.
\end{proof}

Real roots of the denominator algebra define Weyl reflections. Imaginary roots contribute root spaces and product exponents. This division of labor is intrinsic to a Borcherds algebra: the Weyl group is generated by real simple roots, while the imaginary simple roots enter the correction term of the denominator identity.

\section{Automorphic section, orientation line, and Pfaffian operator}
An automorphic square root is a section of a line bundle with character. An orientation is a square root of a determinant line. A Pfaffian operator is a skew family whose regularized Pfaffian is a section of that orientation line. These form a sequence of successive refinements:
\[
\begin{gathered}
 \text{automorphic section}
 \longrightarrow \text{determinant-line identification}\\
 \longrightarrow \text{oriented square-root line}
 \longrightarrow \text{Fredholm family and Pfaffian section}.
\end{gathered}
\]

\begin{proposition}[Logical separation of the four levels]
The product theorem constructs the first object in the sequence. A determinant-line comparison constructs the second. A coherent square root of that line constructs the third. A first-order operator requires a symbol, domain, ellipticity or Fredholm condition, and regularization, and thereby produces the fourth refinement.
\end{proposition}
\begin{proof}
Each stage adds a specified geometric structure. A determinant realization adds an isomorphism from a geometric determinant line to the automorphic line. An oriented realization adds a coherent square-root line and its parity refinement. An operator realization adds a principal symbol, a Fredholm domain, a spectrum, and a regularization whose Pfaffian is the resulting section.
\end{proof}

The title names the monic automorphic section $\Dmon$, characterized by the exact identity $\Dmon^2=\ChiTen$. Part~IV formulates an analytic Pfaffian realization as a further geometric refinement.

\part{The denominator algebra}

\chapter{The Lorentzian root lattice}
Let
\[
 L=\Z f_2\oplus\Z f_3\oplus\Z f_{-2}
\]
with bilinear form
\[
 (f_2,f_{-2})=-1,
 \qquad (f_3,f_3)=2,
\]
and all remaining basis pairings zero. Its signature is $(2,1)$. Define the type-II root lattice
\[
 L_{II}=2\Z f_2\oplus\Z f_3\oplus2\Z f_{-2}
 \cong U(4)\oplus\angles{2}.
\]
The Gram-index map is
\[
 \alpha:\Gammaind=\Z^3\longrightarrow L_{II},
 \qquad
 \alpha(n,\ell,m)=2nf_2-\ell f_3+2mf_{-2}.
\]
Then
\[
 (\alpha(n,\ell,m),\alpha(n',\ell',m'))
 =-2\bigl(2nm'+2n'm-\ell\ell'\bigr),
\]
and
\[
 (\alpha(n,\ell,m),\alpha(n,\ell,m))
 =-2(4nm-\ell^2).
\]

\section{Real simple roots and Weyl vector}
Put
\[
 \delta_1=2f_2-f_3,
 \qquad
 \delta_2=2f_{-2}-f_3,
 \qquad
 \delta_3=f_3.
\]

\begin{proposition}[Type-II Cartan matrix]
The Gram matrix of the real simple roots is
\[
 \bigl((\delta_i,\delta_j)\bigr)_{i,j=1}^3
 =\begin{pmatrix}
 2&-2&-2\\
 -2&2&-2\\
 -2&-2&2
 \end{pmatrix}.
\]
The reflections
\[
 s_i(x)=x-(x,\delta_i)\delta_i
\]
are integral involutions preserving $L_{II}$ and its form.
\end{proposition}
\begin{proof}
Substitution into the basis pairing gives the matrix. Since $(\delta_i,\delta_i)=2$, the reflection formula has integral coefficients. Direct expansion gives $s_i^2=1$ and $(s_ix,s_iy)=(x,y)$.
\end{proof}

The Weyl vector is
\[
 \rho=f_2-\frac12f_3+f_{-2},
 \qquad (\rho,\delta_i)=-1.
\]
In the Gram-index variables, $e^{-\rho}=q^{1/2}r^{1/2}s^{1/2}$.

\chapter[The denominator algebra]{The Gritsenko--Nikulin Lie superalgebra}
The Gritsenko--Nikulin construction specifies a generalized
Kac--Moody Lie superalgebra \(\gIg\), with real simple roots
\(\delta_1,\delta_2,\delta_3\) and the imaginary correction
attached to \(\phiweak\). Its generators and relations supply the
bracket whose denominator is compared with the Borcherds product.

\begin{theorem}[Root supermultiplicities and denominator]
On the active positive support,
\[
 \smult\gIg_{\alpha(n,\ell,m)}=f(nm,\ell).
\]
In the Lorentzian coordinate convention of Gritsenko--Nikulin,
\[
 \den(\gIg)=\Dmon(2Z)=2^{-6}\Dtheta(2Z).
\]
\end{theorem}
\begin{proof}
The real roots and Weyl chamber are the type-II case of the Gritsenko--Nikulin construction. The imaginary simple-root correction is chosen so that its Weyl--Kac--Borcherds product has exponent $f(nm,\ell)$. The resulting product is the doubled-variable monic Igusa form \cite{GN97,GN98a,GN98b}.
\end{proof}

The theorem gives a Lie superalgebra with a bracket, Cartan, Weyl group, root parity, and denominator. The scalar product alone records the signed dimensions of the root spaces; the Gritsenko--Nikulin construction supplies the algebraic operations.

\chapter{The ordered bar as the denominator complex}
Let $\nIg$ be the positive nilpotent subalgebra of $\gIg$. Fix a proper integral height $h$ on the positive root cone and define
\[
 F^{>H}\nIg=\bigoplus_{h(\alpha)>H}\gIg_\alpha,
 \qquad
 \nIg^{(H)}=\nIg/F^{>H}\nIg.
\]
Height additivity makes $F^{>H}\nIg$ an ideal. Properness and local finiteness make $\nIg^{(H)}$ finite-dimensional and nilpotent.

\begin{theorem}[Height-complete Bar--Chevalley comparison]
For every $H$, antisymmetrization gives a natural quasi-isomorphism of conilpotent dg coalgebras
\[
 \Barc\bigl(U(\nIg^{(H)})\bigr)\simeq C_*(\nIg^{(H)}).
\]
The quotient maps preserve these comparisons and induce a quasi-isomorphism in the separated height-complete category
\[
 \Barc_h^\wedge\bigl(U(\nIg)\bigr)
 :=\varprojlim_H\Barc\bigl(U(\nIg^{(H)})\bigr)
 \simeq
 C_{*,h}^\wedge(\nIg)
 :=\varprojlim_H C_*(\nIg^{(H)}).
\]
Its root-graded Euler supercharacter is
\[
 \chi_{\mathrm{root}}\,C_{*,h}^\wedge(\nIg)
 =\prod_{\alpha>0}
 (1-e^{-\alpha})^{\sdim\gIg_\alpha}.
\]
\end{theorem}
\begin{proof}
For each $H$, the antisymmetrization map sends $\Sym(\nIg^{(H)}[1])$ to the normalized bar of $U(\nIg^{(H)})$. The PBW filtrations have the same associated-graded Koszul resolution of $\Sym(\nIg^{(H)})$, so antisymmetrization is a filtered quasi-isomorphism. Naturality gives a morphism of inverse systems. Every fixed root degree is visible in all sufficiently large height quotients and is unchanged thereafter; hence the inverse-limit map is a quasi-isomorphism coefficientwise. An even root vector becomes odd after the Chevalley shift and contributes $1-e^{-\alpha}$; an odd root vector becomes even and contributes its reciprocal. Multiplication over the finite-dimensional root spaces gives the product.
\end{proof}

Let $\mathbb T_{L_{II}}$ be the formal root torus and write $Z_{\mathrm{root}}$ for its logarithmic coordinate, normalized by
\[
 e^{-\alpha(n,\ell,m)}=Q^nR^\ell S^m,
 \qquad
 e^{-\rho}=Q^{1/2}R^{1/2}S^{1/2}.
\]

\begin{corollary}[Bar denominator identity]
On the formal root torus,
\[
 e^{-\rho}\chi_{\mathrm{root}}
 C_{*,h}^\wedge(\nIg)
 =\Dmon(Z_{\mathrm{root}}).
\]
The Gritsenko--Nikulin embedding of the root torus into the Siegel period domain sends $Z_{\mathrm{root}}$ to $2Z$ \cite{GN97,GN98b}. Hence
\[
 e^{-\rho}\chi_{\mathrm{root}}
 C_{*,h}^\wedge(\nIg)
 =\den(\gIg)=\Dmon(2Z)
\]
in Siegel period coordinates.
\end{corollary}
\begin{proof}
Multiply the Euler product in the preceding theorem by the Weyl monomial $e^{-\rho}$.  The denominator theorem identifies the result on the root torus, and the stated embedding gives the Siegel-coordinate formula.
\end{proof}

The denominator is therefore the Weyl-shifted Euler supercharacter of the height-complete interval bar of the positive enveloping algebra.

\chapter{Finite-height currents}
Retain the finite-height quotients $\nIg^{(H)}$ of the preceding chapter.  Define the current Lie conformal superalgebra
\[
 \Cur(\nIg^{(H)})=\kfield[\partial]\tensor\nIg^{(H)},
 \qquad [a_\lambda b]=[a,b],
\]
and its universal chiral enveloping algebra
\[
 V_H=U^{\ch}\!\left(\Cur_E(\nIg^{(H)})\right).
\]

\begin{theorem}[Height-complete current envelope]
The quotient maps induce a coefficientwise inverse limit
\[
 \widehat V_{\Delta_5}=\varprojlim_HV_H.
\]
Every $V_H$ is PBW-admissible as an associative chiral algebra, the current generators are primitive for the standard current coproduct, and
\[
 \sch\widehat V_{\Delta_5}
 =\prod_{j\ge1}\prod_{\alpha>0}
 (1-t^je^{-\alpha})^{-\sdim\gIg_\alpha}.
\]
Moreover,
\[
 \BarAssCh(V_H)
 \simeq
 \BarLieCh\!\left(\Cur_E(\nIg^{(H)})\right),
\]
and these equivalences pass to the separated height completion.
\end{theorem}
\begin{proof}
PBW orders the negative modes $a(-j)$, $j\ge1$.  Even modes contribute symmetric powers and odd modes contribute exterior powers, giving the product.  Proper height makes every coefficient stationary.  The current coproduct sends $a(z)$ to $a(z)\tensor1+1\tensor a(z)$.  Chiral antisymmetrization identifies the associative-chiral bar of the chiral enveloping algebra with the chiral Chevalley coalgebra; the complete separated PBW filtration is coefficientwise finite.
\end{proof}

\section{Classical local theory and exact quantum zero modes}
For every $H$, the finite-dimensional nilpotent Lie superalgebra $\nIg^{(H)}$ defines a cotangent BF theory on $E\times\R$.  Its local dg Lie algebra is
\[
 \Lcal_H=\Omega^{0,*}(E)\widehattensor\Omega^*(\R)\tensor\nIg^{(H)}
\]
and the classical action is
\[
 S_H(A,B)=\int_{E\times\R}\Omega_E\wedge
 \left\langle B,(\dbar+\dd_{\R})A+\frac12[A,A]\right\rangle .
\]

\begin{theorem}[Finite-height BF realization]
The action $S_H$ satisfies the classical master equation and defines a classical holomorphic--topological factorization algebra.  Harmonic reduction along a compact interval produces the finite-dimensional cotangent BV theory of $\nIg^{(H)}$.  The positive-root height grading makes every even adjoint endomorphism strictly height-increasing and every odd adjoint endomorphism odd; hence the modular supercharacter vanishes and the reduced theory satisfies the quantum master equation.

Choose the standard complementary Lagrangian boundary polarizations, fixing the $\nIg^{(H)}[1]$ coordinate at one endpoint and its cotangent coordinate at the other.  The algebraic boundary quantization, the interval amplitude between the two augmentation states, and the holomorphic boundary current envelope are then
\[
 U(\nIg^{(H)}),\qquad
 \mathbbm1\tensor^{\mathbb L}_{U(\nIg^{(H)})}\mathbbm1
 \simeq C_*(\nIg^{(H)}),\qquad V_H.
\]
The height quotient maps preserve the polarizations and yield the inverse systems
\[
 \widehat U_h(\nIg),\qquad
 C_{*,h}^{\wedge}(\nIg),\qquad
 \widehat V_{\Delta_5}.
\]
\end{theorem}
\begin{proof}
The classical master equation is the Jacobi identity paired with invariance of the cotangent pairing.  Elliptic descent gives the classical factorization algebra.  Harmonic reduction along the interval gives a finite-dimensional BV space.  If $x$ is even of positive root height, $\operatorname{ad}_x$ raises height and is nilpotent on the finite quotient, so its restrictions to the even and odd summands both have trace zero; if $x$ is odd, $\operatorname{ad}_x$ is odd and has supertrace zero.  Thus the modular character vanishes and the finite BV quantum master equation follows.

The chosen boundary polarization quantizes the linear Kirillov--Kostant boundary bracket to $U(\nIg^{(H)})$.  Gluing the two augmentation states gives the two-sided derived bar, and the PBW antisymmetrization identifies it with $C_*(\nIg^{(H)})$.  Holomorphic descendants generate $U^{\mathrm{ch}}(\operatorname{Cur}_E\nIg^{(H)})=V_H$.  All maps preserve root height, so they pass coefficientwise to the separated inverse limits.
\end{proof}

\part{The compact scalar character}

\chapter{Reduced curve counting on \texorpdfstring{$K3\times E$}{K3 x E}}
Let $S$ be a nonsingular projective K3 surface and $E$ an elliptic curve. For a primitive class $\beta_h\in H_2(S,\Z)$ with $\beta_h^2=2h-2$, let $N_{g,h,d}$ be the reduced disconnected Gromov--Witten invariant of class $(\beta_h,d)$ with the standard incidence insertion. Define
\[
 \mathcal Z(u,q,s)=
 \sum_{g,h,d}N_{g,h,d}u^{2g-2}q^{h-1}s^{d-1}.
\]

\begin{theorem}[Primitive Igusa cusp-form theorem]
Under $r=e^{iu}$, the partition function is the Laurent expansion
\[
 \mathcal Z(u,q,s)=-\frac1{\ChiTen(r,q,s)}.
\]
Equivalently, in theta normalization,
\[
 \mathcal Z=-2^{12}\Dtheta^{-2}.
\]
\end{theorem}
\begin{proof}
Oberdieck and Pixton identify the primitive reduced partition function with the Laurent expansion of $-\ChiTen^{-1}$ \cite{OP}; their normalization is the monic Borcherds product $\ChiTen$. The second formula is the identity $\ChiTen=2^{-12}\Dtheta^2$.
\end{proof}

The theorem identifies the numerical generating series in its stated
primitive sector. The Borcherds root algebra supplies a separate
algebraic character to compare with that series.

\begin{corollary}[Exact character comparison]
Let $\jmath:\HHtwo\to\HHtwo$ be the half-coordinate map $\jmath(Z)=Z/2$, and let
\[
 \Fock(\nIg)=
 \Sym(\nIg_{\bar0})\tensor\bigwedge(\nIg_{\bar1}).
\]
Then
\[
 \mathcal Z^{\red}_{S\times E}
 =-\jmath^*\!\left(
 e^{2\rho}\sch\bigl(\Fock(\nIg)^{\tensor2}\bigr)
 \right)
 =-\ChiTen^{-1}.
\]
\end{corollary}
\begin{proof}
Put
\[
 P=\prod_{\alpha>0}(1-e^{-\alpha})^{\sdim\gIg_\alpha}.
\]
The denominator theorem gives $\den(\gIg)=e^{-\rho}P=\Dmon(2Z)$, while PBW gives $\sch\Fock(\nIg)=P^{-1}$.  Hence
\[
 e^{2\rho}\sch\bigl(\Fock(\nIg)^{\tensor2}\bigr)
 =\den(\gIg)^{-2}.
\]
Pullback by $\jmath$ sends $\den(\gIg)$ to $\Dmon(Z)$, whose square is $\ChiTen$.  The reduced incidence convention supplies the overall minus sign in the Oberdieck--Pixton theorem.
\end{proof}

The equality compares the completed root-Fock supercharacter with the
numerical reduced generating series. Its proof uses the denominator
identity and the half-coordinate pullback. A map from compact geometric
states to this Fock space, preserving operators and their trace,
is an additional comparison.

\chapter{Extensions over the elliptic classifying stack}
Work on locally noetherian complex schemes with the fppf topology.
Let \(BE\) be the stack of right \(E\)-torsors. For a left
\(E\)-stack \(M\), an object of \([M/E](T)\) is a right
torsor \(P\to T\) and an equivariant map \(P\to M\).
Forgetting the map gives \([M/E]\to BE\). All products of
stacks below are two-fibre products.

\begin{proposition}[The common translation torsor]
\label{igusa:platonic-common-torsor}
There is a natural equivalence
\[
 [(M\times N)/E_{\mathrm{diag}}]
 \simeq[M/E]\times_{BE}[N/E].
\]
The equivalences for any finite number of factors commute with
regrouping and insertion of the terminal \(E\)-stack.
\end{proposition}
\begin{proof}
An object on the right is \((P,m),(Q,n)\) together with a
torsor isomorphism \(\alpha:P\to Q\). Send it to
\((P,m,n\alpha)\). Conversely, use a common torsor twice
and its identity isomorphism. A pair of torsor maps \(a,b\)
on the right satisfies \(b\alpha=\alpha'a\), so \(b\)
is determined by \(a\). The same construction identifies the
maps and their descent isomorphisms. It commutes with base change.
For several factors, every regrouping gives one torsor with the
same list of equivariant maps, proving all stated coherences.
\end{proof}

For \(M=N=E\) with translation, the diagonal quotient is
\(E\), with coordinate \(y-x\). Both separate quotients
are points. Their fibre product over \(BE\) retains the
relative translation, while their ordinary product loses it.

Let \(\mathfrak M\) be the stack of coherent sheaves on
\(X=S\times E\), with families flat over the parameter
scheme, and put \(\mathfrak Q=[\mathfrak M/E]\).
An object of \(\mathfrak Q(T)\) is a torsor \(P\to T\)
and a flat sheaf on \(S\times P\). Fix a polarization and
write \(\mathfrak Q_\alpha\) for a Hilbert-polynomial component.
Translation preserves these polynomials. Let
\(\mathfrak H_{\alpha,\beta}\) classify short exact sequences
\(0\to A\to F\to B\to0\) on a common \(S\times P\),
with all terms flat over \(T\) and polynomials
\(\alpha,\alpha+\beta,\beta\). Its span over \(BE\) is
\[
 \mathfrak Q_\alpha\times_{BE}\mathfrak Q_\beta
 \xleftarrow{p}\mathfrak H_{\alpha,\beta}
 \xrightarrow{q}\mathfrak Q_{\alpha+\beta}.
\]
The category of spans of fppf stacks over \(BE\) has composition
by fibre product and tensor product \(\times_{BE}\). Its unit
object is the identity map of \(BE\).

\begin{theorem}[Geometric Hall correspondence]
\label{igusa:platonic-common-torsor-hall}
For fixed \(\alpha,\beta\), \(q\) is representable by
algebraic spaces and proper. The spans form a coherently associative
unital Hall object in the category of spans over \(BE\).
Its unit is the zero-sheaf span
\(BE\leftarrow BE\to\mathfrak Q\).
\end{theorem}
\begin{proof}
Over a fixed \((P,F)\), the fibre of \(q\) classifies
quotients \(F\twoheadrightarrow B\) of polynomial \(\beta\).
Its kernel is flat, because \(F\) and \(B\) are flat.
An automorphism fixing the quotient and \(F\) is the identity.
Trivializing \(P\) fppf-locally identifies this fibre with the
fixed-polynomial Quot scheme, projective over the parameter scheme
by Grothendieck's theorem \cite[Theorem~3.2]{IgusaGrothendieckQuot1961}.
The quotient functor descends to an algebraic space, and properness
descends with it.

Both composites of three Hall spans classify
\(0\subset A\subset D\subset F\) on the same torsor.
One bracketing concatenates inclusions. The other starts with
\(A\subset F\) and a subobject \(B\subset F/A\), and sets
\(D=F\times_{F/A}B\). These constructions are inverse on
objects and isomorphisms. They preserve flatness and commute with
base change. For longer flags all regroupings identify the same
subobjects and quotients. Their comparison maps agree, proving the
pentagon and higher coherences. Finally \(\mathfrak Q_0=BE\).
The sequences with zero subobject or zero quotient give the unit
and its coherences.
\end{proof}

Properness here concerns the full coherent-sheaf stack on a fixed
pair of Hilbert-polynomial components. A stability restriction can
replace its fibre by an open part of a Quot scheme, which need not
be proper. A bounded charge window also needs closure under the
extensions and flags used in its operations. A reduced critical-coefficient realization must supply pullback,
pushforward, and orientation comparisons on these same flags. Its unit has
stabilizer \(E\). A coefficient theory for this quotient must
include that nonaffine stabilizer, or work relative to \(BE\),
on the unit and every flag.


\chapter{Additive physical charge and quadratic automorphic degree}
Let
\[
 \Lam=\widetilde H(S,\Z)
\]
be the Mukai lattice of signature $(4,20)$. The electric--magnetic charge lattice is
\[
 \Gammphys=\Lam\oplus\Lam,
 \qquad c=(Q,P).
\]
Its addition is the physical superselection law. Define the Gram map
\[
 \PiGram(Q,P)=
 \left(\frac{Q^2}{2},Q\cdot P,\frac{P^2}{2}\right)
 \in\Gammaind.
\]

\begin{proposition}[Polarization identity]
For $c_i=(Q_i,P_i)$,
\[
 \PiGram(c_1+c_2)=\PiGram(c_1)+\PiGram(c_2)+B(c_1,c_2),
\]
where
\[
 B(c_1,c_2)=
 \bigl(Q_1\cdot Q_2,
 Q_1\cdot P_2+Q_2\cdot P_1,
 P_1\cdot P_2\bigr).
\]
The cross-term is a symmetric vector-valued two-cocycle:
\[
 B(c_1,c_2)+B(c_1+c_2,c_3)
 =B(c_2,c_3)+B(c_1,c_2+c_3).
\]
\end{proposition}
\begin{proof}
Expand the three quadratic expressions. The cocycle identity is bilinearity of $B$.
\end{proof}

\section{The polarized central extension}
Define the abelian central extension
\[
 0\longrightarrow\Gammaind\longrightarrow
 \widetilde\Gamma_\Pi
 \longrightarrow\Gammphys\longrightarrow0,
 \qquad
 \widetilde\Gamma_\Pi=\Gammphys\times\Gammaind,
\]
with multiplication
\[
 (c,\lambda)\cdot(d,\mu)
 =\bigl(c+d,\lambda+\mu+B(c,d)\bigr).
\]
The cocycle identity makes this multiplication associative. The polarization identity makes the graph map
\[
 j_\Pi:\Gammphys\longrightarrow\widetilde\Gamma_\Pi,
 \qquad c\longmapsto(c,\PiGram(c))
\]
a homomorphism.

Let $A=\bigoplus_{c\in\Gammphys}A_c$ be a charge-graded associative algebra, and let $\kfield[\Gammaind]$ have monomials $u^\lambda$. Define the polarized central-extension algebra
\[
 \mathcal R_\Pi(A)=A\tensor\kfield[\Gammaind]
\]
with product
\[
 (a_cu^\lambda)\star(b_du^\mu)
 =a_cb_d\,u^{\lambda+\mu+B(c,d)}.
\]
Give $a_cu^\lambda$ degree $(c,\lambda)\in\widetilde\Gamma_\Pi$.

\begin{theorem}[Graph-multiplicative polarization]
The product $\star$ is associative and $\widetilde\Gamma_\Pi$-graded. The map
\[
 \iota_\Pi:A\longrightarrow\mathcal R_\Pi(A),
 \qquad a_c\longmapsto a_cu^{\PiGram(c)}
\]
is an algebra homomorphism whose homogeneous degree is the graph element $j_\Pi(c)$.
\end{theorem}
\begin{proof}
The two parenthesizations differ by
\[
 B(c,d)+B(c+d,e)-B(d,e)-B(c,d+e),
\]
which vanishes by the cocycle identity. The polarization formula gives
\[
 \iota_\Pi(a_c)\star\iota_\Pi(b_d)
 =a_cb_d\,u^{\PiGram(c+d)}.
\]
Thus the graph of the quadratic Gram map is multiplicative inside the central extension.
\end{proof}

The same construction applies to Lie superalgebras and their enveloping algebras. It retains additive physical charge, Gram degree, and the interaction cross-term in one grading group.

\chapter{Point-local and wrapped sectors}
Let $p_E:S\times E\to E$ be projection.

\begin{lemma}[Positive elliptic degree is wrapped]
Let $C\subset S\times E$ be a proper one-dimensional subscheme whose pushforward $(p_E)_*[C]$ has positive degree. Then $p_E(C)=E$.
\end{lemma}
\begin{proof}
Properness makes $p_E(C)$ closed. Positive pushforward degree implies that at least one irreducible component of $C$ dominates a one-dimensional closed subset of $E$. Every one-dimensional closed subset of the irreducible curve $E$ is $E$ itself, so the image of that component, and hence the image of $C$, is all of $E$.
\end{proof}

Projection-finite objects define point-supported sectors on $\Ran(E)$. Positive elliptic degree defines wrapped sectors supported over the entire elliptic curve.

\begin{definition}[The local--wrapped factorization operad]
Let $\mathsf{Ch}_E$ denote the chosen chiral/factorization operad on $E$ and let $\mathsf{Ass}$ denote the ordered associative operad. Define the colored operad $\mathsf{Hyb}_E$ with colors
\[
 \mathsf l,\qquad \mathsf w_m\quad(m\ge1)
\]
by the following operation objects.
\begin{enumerate}
\item Operations with local output are
\[
 \mathsf{Hyb}_E(\mathsf l^{\otimes k};\mathsf l)=\mathsf{Ch}_E(k).
\]
\item For $r\ge1$ and $m=m_1+\cdots+m_r$, operations with wrapped output are
\[
 \mathsf{Hyb}_E
 (\mathsf l^{\otimes k},\mathsf w_{m_1},\ldots,\mathsf w_{m_r};\mathsf w_m)
 =\mathsf{Ch}_E(k)\tensor\mathsf{Ass}(r).
\]
All other operation objects are initial.
\end{enumerate}
Operadic composition uses chiral insertion and disjoint union in the first factor and ordered substitution in the second. The equality
\[
 (\gamma_1\circ\gamma_2)\tensor(a_1\circ a_2)
 = (\gamma_1\tensor a_1)\circ(\gamma_2\tensor a_2)
\]
is the mixed interchange law. The arities $(k,r)=(k,0),(0,2),(k,1)$ recover, respectively, local factorization, wrapped multiplication, and the action of local insertions on a wrapped sector.
\end{definition}

\begin{definition}[Hybrid local--wrapped algebra]
A hybrid local--wrapped algebra over $E$ is an algebra over $\mathsf{Hyb}_E$. It consists of a factorization algebra $\F_{\loc}$, a height- or charge-complete family $W=\prod_{m>0}W_m$, ordered products
\[
 \mu_{m_1,\ldots,m_r}:W_{m_1}\tensor\cdots\tensor W_{m_r}
 \longrightarrow W_{m_1+\cdots+m_r},
\]
and factorization-compatible actions of the local operation objects $\mathsf{Ch}_E(k)$ on every wrapped product. The operad axioms state associativity, local factorization, and compatibility of the two operations.
\end{definition}

A compact Hall theory of $S\times E$ produces a correspondence-valued $\mathsf{Hyb}_E$-algebra when disjoint-support direct sums, wrapped extensions, and the mixed extension correspondences satisfy the same base-change identities.


\part{The Borcherds algebra and its quantum group}

\chapter{The local--wrapped root algebra}
The positive cone is graded by $\gamma=(n,\ell,m)$, with elliptic degree $m$.  Set
\[
 \nIg_{\loc}=\bigoplus_{\substack{\alpha>0\\m(\alpha)=0}}\gIg_\alpha,
 \qquad
 \nIg_{\wrap}=\bigoplus_{\substack{\alpha>0\\m(\alpha)>0}}\gIg_\alpha.
\]

\begin{proposition}[Semidirect local--wrapped decomposition]
The space $\nIg_{\loc}$ is a Lie subsuperalgebra, $\nIg_{\wrap}$ is a Lie ideal, and
\[
 \nIg=\nIg_{\loc}\ltimes\nIg_{\wrap}.
\]
The proper height filtration preserves both summands.
\end{proposition}
\begin{proof}
The bracket adds root degrees.  The sum of two roots of elliptic degree zero again has degree zero, and addition of a nonnegative degree to a positive degree remains positive.  Height is additive on every nonzero bracket.
\end{proof}

Define the complete enveloping Hopf algebra and current envelope
\[
 \mathbb U_{\Delta_5}=\widehat U_h(\nIg),
 \qquad
 \mathbb V_{\Delta_5,E}:=\widehat V_{\Delta_5}=\varprojlim_HU^{\ch}\!\left(\Cur_E(\nIg^{(H)})\right).
\]
The adjoint action of $\nIg_{\loc}$ on $\nIg_{\wrap}$ extends by derivations to $\widehat U_h(\nIg_{\wrap})$.  Let $\# $ denote the corresponding completed smash product.

\begin{definition}[Current-generated local--wrapped operad]
The colored operad $\mathsf{Hyb}^{\cur}_{E,\Delta}$ is generated by the chiral current operations of $\Cur_E(\nIg_{\loc})$, the ordered multiplications of $\widehat U_h(\nIg_{\wrap})$, and the adjoint derivation action of the former on the latter, modulo the chiral Jacobi, associative, derivation, and semidirect-product relations.
\end{definition}

\begin{theorem}[Local--wrapped Borcherds decomposition]
Let $\nIg_{\loc}$ be the sum of positive root spaces of elliptic degree $m=0$ and let $\nIg_{\wrap}$ be the sum of those with $m>0$.  Then
\[
 \nIg=\nIg_{\loc}\ltimes\nIg_{\wrap},
\]
where $\nIg_{\loc}$ is a Lie subsuperalgebra and $\nIg_{\wrap}$ is a Lie ideal.  For the proper-height completion,
\[
 \widehat U_h(\nIg)
 \cong
 \widehat U_h(\nIg_{\wrap})\#\widehat U_h(\nIg_{\loc}),
\]
and
\[
 \widehat C_{*,h}(\nIg)
 \simeq
 \Tot C_*(\nIg_{\loc};\widehat C_{*,h}(\nIg_{\wrap})).
\]
The primitive Lie superalgebra of the completed enveloping algebra is the height completion $\widehat{\nIg}_h=\varprojlim_H\nIg^{(H)}$, and the Weyl-shifted Euler supercharacter of the completed Chevalley complex is $\Dmon(2Z)$.
\end{theorem}
\begin{proof}
Root addition adds elliptic degree.  Thus the $m=0$ summand is closed under the bracket, and every bracket with an $m>0$ root again has $m>0$.  PBW orders wrapped generators before local generators and gives the completed smash product.  The Chevalley complex of a semidirect product is the Hochschild--Serre total complex.  Primitive elements of a completed enveloping algebra are recovered height by height.  The denominator theorem supplies the last character identity.
\end{proof}

\begin{remark}
The theorem constructs a Borcherds root algebra with local and wrapped colors.  The word Hall is reserved for an algebra obtained from geometric extension correspondences.
\end{remark}

\chapter{The formal Igusa quantum group}
Let $\mathfrak b_+$ be the positive Borcherds Borel half.  Form the restricted graded dual $\mathfrak b_-:=\mathfrak b_+^{\vee,\mathrm{gr}}$ with bracket and cobracket transported by the invariant pairing, identify the two Cartan copies along their common radical quotient, and complete by a proper root height.  The completed Drinfeld double of this paired Borel datum is denoted $\gIg^{\mathrm{pair}}$.  Its positive and negative halves, diagonal Cartan, invariant form, and canonical tensor form a height-admissible Manin triple and hence a quasitriangular Lie superbialgebra.

\begin{definition}[Height-admissible Lie superbialgebra]
A root-graded Lie superbialgebra is height-admissible when every root space is finite-dimensional, every bounded-height region contains finitely many roots, and the bracket, cobracket, invariant tensor, and all iterated operations are coefficientwise finite in the height completion.
\end{definition}

\begin{proposition}[Height admissibility]
The paired Gritsenko--Nikulin Borcherds Lie superbialgebra $\gIg^{\mathrm{pair}}$ is height-admissible.
\end{proposition}
\begin{proof}
The Borcherds root spaces are finite-dimensional and the chosen height is proper.  The bracket and cobracket preserve total root degree.  In a bounded total height, finitely many decompositions of a root occur.  The canonical tensor is therefore a coefficientwise finite sum in every height quotient.
\end{proof}

\begin{theorem}[Quantized Borcherds superalgebra]
Choose the symmetrizable Borcherds--Cartan superdatum and multiplicity index set that present $\gIg$.  Hong's construction gives a quantized enveloping Hopf superalgebra
\[
 U_q(\gIg)
\]
over $\Q(q)$, generated by its Cartan elements and the real and imaginary simple-root generators, with triangular decomposition
\[
 U_q(\gIg)\cong U_q^-(\gIg)\tensor U_q^0(\gIg)\tensor U_q^+(\gIg).
\]
Its PBW root spaces have the same graded ranks as the classical Borcherds superalgebra.  The canonical pairing of the two Borel halves yields a universal $R$-matrix in the root-height completion of
\[
 U_q^+(\gIg)\widehattensor U_q^-(\gIg).
\]
Consequently the quantization preserves the root-multiplicity data that enter $\Dmon(2Z)$.
\end{theorem}
\begin{proof}
The Gritsenko--Nikulin algebra is a symmetrizable Borcherds--Kac--Moody Lie superalgebra.  Apply the generators-and-relations construction of the quantized Borcherds superalgebra.  The nondegenerate bilinear form, triangular decomposition, centre, and universal $R$-matrix are established by Hong \cite{Hong98}.  The PBW theorem identifies the associated graded root spaces with their classical counterparts.  Since the denominator exponents are the classical signed root multiplicities, they are unchanged by the flat root-graded deformation.
\end{proof}

\begin{remark}
The quantum group is attached to the Borcherds generalized Cartan superdatum.  A geometric Hall quantum group is identified with it through a Hall coproduct, pairing, completion, and primitive comparison.
\end{remark}

\chapter{A holomorphic--topological BF realization}
For every height bound $H$, let $\nIg^{(H)}$ be the finite-dimensional nilpotent quotient.  On $E\times\R$ form the shifted cotangent theory of
\[
 \Lcal_H=\Omega^{0,*}(E)\widehattensor\Omega^*(\R)\tensor\nIg^{(H)}
\]
with action
\[
 S_H(A,B)=\int_{E\times\R}\Omega_E\wedge
 \left\langle B,(\dbar+\dd_{\R})A+\frac12[A,A]\right\rangle .
\]

\begin{theorem}[Classical local and exact zero-mode realization]
For every $H$:
\begin{enumerate}
\item $S_H$ satisfies the classical master equation and defines a classical holomorphic--topological factorization algebra;
\item harmonic reduction along a compact interval gives the finite-dimensional cotangent BV theory of $\nIg^{(H)}$, whose positive-root grading makes its modular supercharacter zero and hence gives the quantum master equation;
\item after choosing the complementary Lagrangian boundary polarizations, the boundary algebra is $U(\nIg^{(H)})$, the interval amplitude is $\mathbbm1\tensor^{\mathbb L}_{U(\nIg^{(H)})}\mathbbm1\simeq C_*(\nIg^{(H)})$, and the holomorphic boundary current envelope is $V_H$.
\end{enumerate}
The transition maps yield a pro-object with boundary algebra $\widehat U_h(\nIg)$, interval complex $C_{*,h}^{\wedge}(\nIg)$, and current envelope $\widehat V_{\Delta_5}$.
\end{theorem}
\begin{proof}
The classical master equation is the Jacobi identity and invariance of the cotangent pairing.  The classical observables satisfy factorization descent.  Harmonic reduction is finite-dimensional.  An even adjoint operator raises positive-root height and is nilpotent on the finite quotient; an odd adjoint operator is odd.  Both have supertrace zero, so the modular term in the finite BV Laplacian vanishes.  The chosen polarization quantizes the boundary Lie--Poisson bracket to the enveloping algebra; the two augmentation states give the two-sided bar, identified with Chevalley chains by PBW.  Holomorphic descendants give the current algebra.  All structure maps preserve height.
\end{proof}

The theorem gives the precise field-theoretic realization used in this volume: a classical local theory together with an exact quantum protected zero-mode sector.

\chapter{The determinant and Fredholm square-root operators}
Let
\[
 \Vcal_H=\bigoplus_{\substack{\alpha>0\\h(\alpha)\le H}}\gIg_\alpha
\]
and define
\[
 \mathscr D_H|_{\gIg_\alpha}=(1-e^{-\alpha})\id.
\]

\begin{theorem}[Universal determinant section]
The finite-height Berezinian is
\[
 \Ber(\mathscr D_H)
 =\prod_{\substack{\alpha>0\\h(\alpha)\le H}}
 (1-e^{-\alpha})^{\sdim\gIg_\alpha}.
\]
The coefficientwise limit satisfies
\[
 e^{-\rho}\Ber(\mathscr D_\infty)=\Dmon(2Z).
\]
After pullback by $\jmath(Z)=Z/2$,
\[
 \jmath^*\!\left(e^{-\rho}\Ber(\mathscr D_\infty)\right)=\Dmon(Z),
 \qquad \Dmon(Z)^2=\ChiTen(Z).
\]
\end{theorem}
\begin{proof}
The Berezinian of a diagonal endomorphism is the product of its eigenvalues with exponent equal to the superdimension of the eigenspace.  Proper height makes every coefficient stationary.  The denominator theorem supplies the Weyl line and the root-to-Siegel coordinate map.
\end{proof}

\section{Trace-class realization on the root Hilbert superspace}
Choose homogeneous orthonormal bases of the finite-dimensional root spaces and set
\[
 \Hcal_{\Delta}=\widehat\bigoplus_{\alpha>0}\gIg_\alpha
\]
with its induced $\Z/2$ grading.  Let $\mathfrak C_{\mathrm{abs}}$ be the subchamber on which
\[
 \sum_{\alpha>0}\dim\gIg_\alpha\,
       \abs{e^{-\alpha(Z)}}<\infty.
\]
For $Z\in\mathfrak C_{\mathrm{abs}}$, define the trace-class diagonal operator
\[
 K_Z|_{\gIg_\alpha}=e^{-\alpha(Z)}\id,
 \qquad D_Z=1-K_Z.
\]
On each parity summand form the skew doubling
\[
 \mathbb A_{Z,\bar\epsilon}=
 \begin{pmatrix}0&D_{Z,\bar\epsilon}\\-D_{Z,\bar\epsilon}^{\mathsf t}&0\end{pmatrix}.
\]
Orient every finite-height truncation by the fixed root order and normalize its Pfaffian to $1$ at $K_Z=0$.  The oriented Fredholm super-Pfaffian is the stabilized ratio
\[
 \operatorname{sPf}_{F}(\mathbb A_Z)
 =\frac{\Pf_F(\mathbb A_{Z,\bar0})}
        {\Pf_F(\mathbb A_{Z,\bar1})}.
\]

\begin{theorem}[Polarized Fredholm realization]
Choose the protected parity refinement and the root-space Hilbert norms above.  For every $Z$ in the absolute-convergence chamber, the diagonal operator $K_Z|_{\gIg_\alpha}=e^{-\alpha(Z)}\id$ is trace class on the even and odd root Hilbert spaces.  For $D_Z=1-K_Z$,
\[
 \sdet_F(D_Z)
 =\frac{\det_F(D_Z|_{\bar0})}{\det_F(D_Z|_{\bar1})}
 =\prod_{\alpha>0}(1-e^{-\alpha(Z)})^{\sdim\gIg_\alpha}.
\]
After choosing the standard polarization of $\Hcal_\Delta\oplus\Hcal_\Delta^*$ and its orientation, the paritywise skew block operators $\mathbb A_{Z,\bar0}$ and $\mathbb A_{Z,\bar1}$ have an oriented Fredholm super-Pfaffian
\[
 \operatorname{sPf}_F(\mathbb A_Z)=\frac{\Pf_F(\mathbb A_{Z,\bar0})}{\Pf_F(\mathbb A_{Z,\bar1})}
 =\sdet_F(D_Z).
\]  Hence
\[
 e^{-\rho}\sdet_F(D_Z)=\Dmon(2Z).
\]
\end{theorem}
\begin{proof}
The defining summability condition for \(\mathfrak C_{\mathrm{abs}}\)
is precisely the trace norm of the diagonal operator, summed in both
parities.  The Fredholm determinant of $1-K_Z$ is the convergent product over eigenvalues; taking the even/odd ratio gives the superdeterminant.  On each finite root-height truncation the Pfaffian of the oriented skew doubling equals the determinant of $D_Z$.  Trace-norm convergence of the truncations passes this identity to the Fredholm limit.  The final equality is the Gritsenko--Nikulin denominator formula.
\end{proof}

\begin{remark}
The polarization and orientation in this theorem belong to the root Hilbert space.  A compact moduli-space orientation is a separate line bundle and is compared with this one by a specified intertwiner.
\end{remark}

This constructs the square-root operator on the root Hilbert superspace
where its stated summability condition holds. The nonemptiness estimate
below concerns the minimal virtual superspace, whose dimensions are
\(\lvert\sdim\gIg_\alpha\rvert\). Applying that estimate to the
full root superspace requires a bound on both parity dimensions.
A compact geometric realization requires an operator intertwiner and
a compatible map of determinant lines.

\chapter{What the automorphic character determines}
\begin{theorem}[Denominator data and reconstruction data]
The Borcherds product determines:
\begin{enumerate}
\item the signed root multiplicity $\sdim\gIg_\alpha$ in every active root degree;
\item the Weyl vector and denominator character;
\item the divisor and automorphic multiplier of the section.
\end{enumerate}
A Lie bracket is supplied by the Borcherds--Kac--Moody presentation, parity dimensions by a parity-refined root datum, Hall multiplication by extension correspondences, and an oriented square root by a square-root line with monodromy.  These reconstruction data refine the scalar product independently.
\end{theorem}
\begin{proof}
The first three items are read from the exponent, prefactor, and automorphy law.  Independence is witnessed by explicit pairs: two brackets on the same graded vector space have the same PBW character; an even--odd pair has signed dimension zero; distinct Hall correspondences may have the same Euler product; and twisting a square-root line by a sign local system preserves its square.  Hence each refined object is reconstructed from the additional datum named in the statement.
\end{proof}

\begin{recognition}[Complete primitive recognition]
Let $\mathfrak p$ be a positive graded Lie superalgebra and let $\Psi:\nIg\twoheadrightarrow\mathfrak p$ be a graded Lie map.  The map is an isomorphism if either:
\begin{enumerate}
\item the parity-refined Hilbert series of $U(\nIg)$ and $U(\mathfrak p)$ agree; or
\item each primitive degree has prescribed pure parity and the signed PBW characters agree.
\end{enumerate}
The induced height completions, Chevalley complexes, and current envelopes are then isomorphic.
\end{recognition}
\begin{proof}
Choose the first height containing a homogeneous kernel vector.  All decomposable PBW terms in that height are determined by lower heights and agree.  The primitive parity dimension therefore drops, contradicting the first hypothesis.  In the second case the logarithm of the PBW product recovers signed primitive dimensions inductively, and prescribed parity recovers ordinary dimensions.  Naturality of completion, Chevalley chains, and current envelopes gives the final assertion.
\end{proof}

\part{The universal Borcherds Hall and three-Calabi--Yau models}

\chapter{Finite simple-root windows}
Let $I_\Delta$ be the multiset of simple roots of the Gritsenko--Nikulin Borcherds superalgebra.  A repeated imaginary simple root is written once for each copy.  The parity map is
\[
 p:I_\Delta\longrightarrow\Z/2.
\]
Choose a proper height function
\[
 h:I_\Delta\longrightarrow\Z_{>0}
\]
whose sublevel sets are finite.  Put
\[
 I_H=\set{i\in I_\Delta:h(i)\le H}.
\]
The restriction of the invariant form gives a finite symmetric Borcherds--Cartan matrix
\[
\begin{aligned}
 A_H&=(a_{ij})_{i,j\in I_H},\\
 a_{ii}&=2 &&\text{for real }i,\\
 a_{ii}&\leq0 &&\text{for imaginary }i,\\
 a_{ij}&\leq0 &&\text{for }i\neq j.
\end{aligned}
\]
The evenness of the root lattice makes every imaginary diagonal entry even.

\begin{construction}[The Igusa Borcherds quiver]
Order $I_H$.  For $i<j$ place $-a_{ij}$ arrows from $i$ to $j$.  At an imaginary vertex $i$ place
\[
 g_i=1-\frac{a_{ii}}2
\]
loops.  A real vertex has loop number zero.  Denote the resulting quiver by $Q_{\Delta,H}$.
\end{construction}

\begin{proposition}[Euler form]
The symmetrized Euler form of $Q_{\Delta,H}$ is $A_H$:
\[
 \angles{e_i,e_j}+\angles{e_j,e_i}=a_{ij}.
\]
\end{proposition}
\begin{proof}
For $i\ne j$, the total number of arrows between the two vertices is $-a_{ij}$, so the symmetrized Euler entry is $a_{ij}$.  At $i$ the value is $2-2g_i=a_{ii}$.
\end{proof}

The inclusions $I_H\subset I_{H+1}$ identify $Q_{\Delta,H}$ with a full subquiver of $Q_{\Delta,H+1}$.  Representations supported on the old vertices form an exact extension-closed subcategory.  This gives compatible Hall embeddings.


\chapter{The real-root seed window}
The three real simple roots are
\[
 \delta_1=2f_2-f_3,
 \qquad
 \delta_2=2f_{-2}-f_3,
 \qquad
 \delta_3=f_3,
\]
with Gram matrix
\[
 A_{\mathrm{re}}=
 \begin{pmatrix}
 2&-2&-2\\
 -2&2&-2\\
 -2&-2&2
 \end{pmatrix}.
\]
Choose the order $\delta_1<\delta_2<\delta_3$.  The real-root quiver $Q_{\mathrm{re}}$ has two arrows from vertex $i$ to vertex $j$ for each $i<j$.

\begin{proposition}[Real-root Hall seed]
The symmetrized Euler form of $Q_{\mathrm{re}}$ is $A_{\mathrm{re}}$.  The Hall composition algebra generated by its three simple representations is the positive quantum Kac--Moody algebra with Cartan matrix $A_{\mathrm{re}}$.
\end{proposition}
\begin{proof}
The diagonal Euler entries are $2$.  The total number of arrows between distinct vertices is two, hence the symmetrized off-diagonal entries are $-2$.  The Ringel--Green Hall theorem identifies the composition algebra with the positive quantum algebra defined by that symmetric Cartan matrix.
\end{proof}

The three vertices already contain the Weyl group and the reflective divisor.  Imaginary vertices add the automorphic correction.  The finite-height tower therefore begins with a completely explicit hyperbolic Hall algebra and enlarges it by one finite collection of imaginary simple roots at each stage.

\section{The first Ginzburg differential}
Let $a_{ij}^{(1)},a_{ij}^{(2)}:i\to j$ be the two arrows for $i<j$.  Add reverse arrows $(a_{ij}^{(r)})^*:j\to i$ of degree $-1$ and loops $t_i$ of degree $-2$.  The canonical three-Calabi--Yau completion has
\[
 d a_{ij}^{(r)}=0,
 \qquad
 d(a_{ij}^{(r)})^*=0,
\]
and
\begin{equation}\label{eq:vol3-ginzburg-differential}
 d t_i=
 \sum_{a:s(a)=i}aa^*
 -\sum_{a:t(a)=i}a^*a.
\end{equation}
The identity $d^2=0$ follows immediately.  Formula \eqref{eq:vol3-ginzburg-differential} is the derived moment-map equation for the doubled real-root quiver.

\begin{calculation}[Vertex one]
At vertex $1$,
\[
 d t_1=
 \sum_{r=1}^2
 \left(
 a_{12}^{(r)}(a_{12}^{(r)})^*
 +a_{13}^{(r)}(a_{13}^{(r)})^*
 \right).
\]
At vertex $3$,
\[
 d t_3=-\sum_{r=1}^2\left(
 (a_{13}^{(r)})^*a_{13}^{(r)}+
 (a_{23}^{(r)})^*a_{23}^{(r)}\right).
\]
The sum over vertices is $\sum_a(aa^*-a^*a)$.
It vanishes in the cyclic quotient and after matrix trace,
but generally not in the path algebra.
\end{calculation}
\begin{proof}
The outgoing arrows at vertex $1$ have targets $2$ and $3$.
The incoming arrows at vertex $3$ have sources $1$ and $2$.
The displayed formulas follow from these incidence lists.
Each arrow contributes one commutator to the sum over vertices.
The trace of a commutator vanishes, giving the asserted trace identity.
\end{proof}


\chapter{Hall realization of the positive algebra}
Let $\mathcal A_H=\Rep_{\mathbb F_q}(Q_{\Delta,H})$.  Its twisted Hall product is
\[
 [M]*[N]
 =v^{\angles{\dim M,\dim N}}
 \sum_{[L]}
 \abs{\set{N'\subset L:N'\cong N,\ L/N'\cong M}}
 [L],
 \qquad v^2=q.
\]

The class $[0]$ is the unit: the required subobject is unique on either
side. Associativity follows by counting flags $N_1\subset N_2\subset L$
in two orders. The twist factors agree because the Euler form is
bilinear:
\[
 \langle\alpha,\beta\rangle+\langle\alpha+\beta,\gamma\rangle
 =\langle\beta,\gamma\rangle+\langle\alpha,\beta+\gamma\rangle.
\]
Thus this normalization defines the unital twisted Hall algebra on
isomorphism classes.

\begin{theorem}[Finite-window composition algebra]
At each vertex $i$ of $Q_{\Delta,H}$, choose the one-dimensional
representation $S_i$ supported at $i$, with every loop acting by zero.
Use the generic composition algebra generated by these classes across
finite fields, with the twisted subobject-count multiplication above.
After extending coefficients from $\mathbb Z[v,v^{-1}]$ to
$\mathbb C(v)$, this is the positive quantum generalized Kac--Moody
algebra for the even Borcherds--Cartan matrix $A_H$.
The parameter in Kang--Schiffmann is $v_{\mathrm{KS}}=v^{-1}$.
Their specified semisimple perverse-sheaf construction on nilpotent
representation varieties gives the corresponding canonical basis.
\end{theorem}
\begin{proof}
This is the selected-simple composition construction of
Kang--Schiffmann, Section~2.3 and Theorem~2.2, with one selected
zero-loop simple at each vertex. Their geometric realization uses
its split graded Grothendieck group over
$\mathbb Z[v_{\mathrm{KS}},v_{\mathrm{KS}}^{-1}]$
(Theorem~4.1), with the stated inversion of parameter
\cite{KangSchiffmann}. The theorem does not use every simple object
of the full representation category. A vertex with $c$ loops has
$q^c$ one-dimensional simple representations over $\mathbb F_q$,
whereas only the chosen zero-loop representation is the generator here.
\end{proof}

A semi-derived Hall realization of a full quantum group additionally
specifies two-periodic complexes, the invertible acyclic classes and
Cartan reduction \cite{LuBorcherdsHall,AxtellLee}.
Its coproduct, pairing and completion are separate from the positive
composition algebra used here.

\section{Parity and the quantized superdatum}
The Gritsenko--Nikulin denominator algebra carries a parity map on its simple-root multiset.  It is additive on the free lattice of labelled simple generators.
Projection to the Lorentzian degree can identify contributions of
different parity. Both parity dimensions must therefore be retained
in each projected root space.  Hong's construction applied to this Borcherds superdatum gives the positive Borel
\[
 U_q^+(\gIg^{(H)})
\]
and the full quantized Borcherds superalgebra $U_q(\gIg^{(H)})$.

\begin{theorem}[Ordinary and super presentations]
The ordinary quiver Hall presentation and Hong's quantum
superalgebra use the same specified Cartan matrix and labelled
simple generators. Their relation ideals and parity conventions
are different inputs. Equality of their Cartan gradings alone does
not identify their PBW characters.
\end{theorem}
\begin{proof}
The two presentations are those of the preceding Hall construction
and Hong's superalgebra construction. The difference already occurs
for a single isotropic generator. The enveloping algebra of a
one-dimensional even abelian Lie algebra is \(\mathbb C[x]\).
Declaring its generator odd only as a graded vector space leaves
a nonzero vector \(x^2\) in degree two. The enveloping algebra
of a one-dimensional odd abelian Lie superalgebra is
\(\bigwedge(\xi)\), since its defining relation is
\(2\xi^2=[\xi,\xi]=0\). Its degree-two component is zero.
Thus parity assignment to an ordinary PBW basis does not produce
the super relation ideal or its character. A comparison must
respect these relations degree by degree.
\end{proof}

\chapter{The three-Calabi--Yau completion}
Let $kQ_{\Delta,H}$ be the path algebra. Its algebraic
three-Calabi--Yau completion is
\[
 \Gamma_{\Delta,H}=\Pi_3(kQ_{\Delta,H}).
\]
Write $J$ for the ideal generated by all arrows in its graded quiver.
Its completed version is the separate topological dg algebra
$\widehat\Gamma_{\Delta,H}=\varprojlim_n\Gamma_{\Delta,H}/J^n$.
The algebraic Calabi--Yau theorem is used for $\Gamma_{\Delta,H}$;
a topological Calabi--Yau assertion for $\widehat\Gamma_{\Delta,H}$
requires its completed bimodule category.
This is the Ginzburg model with zero potential: a degree-minus-one
reverse arrow for every arrow and a degree-minus-two loop at each
vertex. The loop differential is the signed sum of incident
arrow--reverse-arrow products.

\begin{theorem}[Borcherds--Ginzburg projections]
For each finite quiver \(Q_{\Delta,H}\), its undeformed
three-Calabi--Yau completion \(\Gamma_{\Delta,H}\) is the
Ginzburg dg algebra with zero potential. If \(H\le H'\) and
\(Q_{\Delta,H}\) is the full subquiver on a subset of vertices,
there is a surjective dg algebra map
\[
 \pi_{H'H}:\Gamma_{\Delta,H'}\longrightarrow\Gamma_{\Delta,H}.
\]
It kills the added vertices, every arrow incident to them, the
corresponding reverse arrows, and the added degree-minus-two loops.
These maps compose strictly and extend continuously to the completed algebras.
\end{theorem}
\begin{proof}
The finite Calabi--Yau assertion is the Ginzburg model of Keller's
completion theorem \cite{KellerCY,GinzburgCY}. We concatenate
paths from left to right, so \(aa^*\) is based at the source
of \(a\). Original and reverse arrows have zero differential.
For the degree-minus-two loop \(t_i\), the differential is
\[
 d t_i=\sum_{s(a)=i}aa^*-\sum_{t(a)=i}a^*a.
\]
The projection deletes exactly the terms involving an added vertex,
leaving the same sum for the full subquiver. It therefore commutes
with the differential. It preserves path length and extends to the
completed path algebra. Successive projections kill the union of
their omitted vertices and arrows, proving strict composition.
\end{proof}

The evident inclusion of graded quivers does not give the reverse
dg map. A single old vertex has loop differential zero. Adding
an arrow \(a:i\to j\) changes that differential to \(aa^*\)
at \(i\). The map fixing the old loop therefore fails to commute
with the differential.

Write \(\Gamma_\Delta=(\Gamma_{\Delta,H},\pi_{H'H})_H\)
for this pro-system of dg algebras. It is separate from the
ordinary Hall embeddings obtained by extension by zero on
representations of a full subquiver. A direct system of
Calabi--Yau dg categories compatible with those Hall embeddings
requires additional dg functors. It does not follow from the
inclusion of the displayed Ginzburg generators.

The quantum group retains its separately specified root-height
completion. Its transition maps must preserve the Cartan generators,
coproduct, and pairing. Quotienting the positive nilpotent Lie
algebra by height and truncating a full quantum group are different
operations.

\chapter{Hall states and protected superstates}
The representation varieties of $Q_{\Delta,H}$ carry the
perverse-sheaf complexes used in the composition-algebra construction.
The full equivariant derived category is a larger carrier.
Already for one real vertex and dimension one, its quotient is
$B\mathbb G_m$, with $H^*(B\mathbb G_m;\mathbb Q)=\mathbb Q[c_1]$,
where $|c_1|=2$. Its total equivariant cohomology is infinite-dimensional.
Finite root multiplicities therefore refer to the specified composition
algebra, not to total equivariant cohomology of this stack.

\begin{definition}[Igusa Hall category]
For every dimension vector $\mathbf d$, let
\[
 \mathfrak R_{H,\mathbf d}=
 [\operatorname{Rep}_{\mathbf d}(Q_{\Delta,H})/G_{\mathbf d}].
\]
The finite-height Hall category is the direct sum of the equivariant constructible derived categories on the $\mathfrak R_{H,\mathbf d}$, equipped with Hall convolution along the stack of short exact sequences.  Denote it by $\mathcal B_{\Delta,H}$ and set
\[
 \mathcal B_\Delta=\varinjlim_H\mathcal B_{\Delta,H}.
\]
\end{definition}

Independently, the protected one-particle superspace is the root-graded superspace
\[
 \mathcal H_{\mathrm{1p},H}
 =\bigoplus_{\alpha>0}\gIg^{(H)}_\alpha,
\]
with the parity of the Borcherds superalgebra.  Its multiparticle completion is the PBW/Fock superspace of $U(\nIg^{(H)})$.

\begin{theorem}[Hall and protected-state comparison]
The composition subalgebra in the finite-window Hall theorem
is the ordinary positive quantum generalized Kac--Moody algebra.
This statement does not identify the Grothendieck group of the full
category $\mathcal B_{\Delta,H}$ with that composition algebra.  The protected one-particle superspace has signed root multiplicities
\[
 \sdim\mathcal H_{\mathrm{1p},H,\alpha}
 =\sdim\gIg^{(H)}_\alpha.
\]
Its height-complete PBW character is
\[
 \sch\Fock(\nIg)
 =\prod_{\alpha>0}(1-e^{-\alpha})^{-\sdim\gIg_\alpha},
\]
and its Weyl-shifted inverse is $\Dmon(2Z)$.  The composition canonical basis and the protected root presentation use the specified simple-root indexing and Cartan grading; the parity refinement belongs to the Borcherds superdatum.
\end{theorem}
\begin{proof}
The composition-algebra assertion is the finite-window theorem. The protected assertion is PBW for the Borcherds Lie superalgebra.  Proper height passes both constructions to the limit.  The Gritsenko--Nikulin denominator theorem gives the final product identity.
\end{proof}

\chapter{Extension determinants and critical orientations}
Let \(\mathcal C\) be a three-Calabi--Yau dg category over
\(\mathbb C\), with its chosen trace. For families of objects
\(A,B\) over a base \(T\), let \(C(A,B)\) denote the relative
derived Hom complex on \(T\). Suppose these complexes are perfect,
and that the trace gives base-change-compatible duality
\[
 C(B,A)\simeq C(A,B)^\vee[-3].
\]
Put \(L(A,B)=\det C(A,B)\) and \(D(A)=L(A,A)\).
The determinant convention is
\(\det C=\bigotimes_i(\det C^i)^{(-1)^i}\), with symmetry sign
specified by \(\chi(C)\bmod2\).

\begin{theorem}[Extension determinant]
For an exact triangle \(A\to F\to B\to A[1]\), there is a
base-change-compatible isomorphism
\[
 D(F)\simeq D(A)\otimes D(B)\otimes L(A,B)^{\otimes2}.
\]
The determinant comparisons obtained by refining a flag agree, with
the stated determinant symmetry signs.
\end{theorem}
\begin{proof}
Apply derived Hom to the triangle in each variable. Determinant
additivity gives
\[
 D(F)\simeq L(B,A)\otimes D(A)\otimes D(B)\otimes L(A,B).
\]
Dualization and the odd shift each invert the determinant line.
Thus Calabi--Yau duality identifies \(L(B,A)\) with \(L(A,B)\).
Reorder by determinant symmetry to obtain the formula. For a full
flag, all iterated decompositions identify with the same ordered
tensor product of \(C(A_i,A_j)\), for its successive quotients
\(A_i\). Determinant additivity for refined filtrations proves
coherence. Each operation commutes with base change.
\end{proof}

A square-root orientation consists of lines \(o_A\) and
isomorphisms \(o_A^{\otimes2}\simeq D(A)\). Extension
compatibility requires maps
\[
 o_A\otimes o_B\otimes L(A,B)\longrightarrow o_F
\]
whose squares are the determinant comparisons. Choosing coordinates
on a smooth critical chart trivializes its ambient canonical line.
Hall convolution requires transition maps compatible with the cross
line and the zero object. On three-step flags, the two orientation
maps can differ by an element of \(\mu_2\), even when their
squares agree. That sign must be trivialized coherently.

\begin{example}[A nontrivial cross line]
Let \(X\) be a smooth projective Calabi--Yau threefold, with a
chosen volume form, and let \(x\in X\). On
\(X\times T\), where \(T=\mathbb P^1\), take
\[
 A=\mathcal O_X\boxtimes\mathcal O_T(1),\qquad
 B=\mathcal O_x\boxtimes\mathcal O_T.
\]
Then \(C(A,B)=\mathcal O_T(-1)[0]\), and relative Serre
duality gives \(C(B,A)=\mathcal O_T(1)[-3]\).
Thus \(L(A,B)=L(B,A)=\mathcal O_T(-1)\).
The self-Hom complexes of \(A\) and \(B\) have constant
determinant lines, since the factors \(\mathcal O_T(1)\)
cancel in the first self-Hom. Nevertheless
\[
 D(A\oplus B)=D(A)\otimes D(B)\otimes\mathcal O_T(-2).
\]
The orientation cross factor is \(\mathcal O_T(-1)\).
A sign change cannot replace this line of nonzero degree.
\end{example}

\begin{corollary}[Oriented critical coefficients]
Suppose the chosen critical charts have square-root orientation
maps on overlaps and extension flags, with unit and flag coherences.
Suppose their vanishing-cycle comparison maps preserve these data.
Then the local vanishing-cycle complexes descend to the specified
oriented coefficient object. Their Hall convolution is associative
whenever its pullback and pushforward maps are defined and satisfy
base change on the flag diagrams.
\end{corollary}
\begin{proof}
The overlap cocycle identifies the local coefficient complexes by
descent. On a three-step flag, orientation coherence identifies
the two Thom--Sebastiani coefficient maps. Proper or compact-support
base change, on the stipulated domain, identifies the two composites
of convolution. The zero flag gives the unit equation.
\end{proof}

For the quiver charts, the remaining construction is an equivariant
system of square roots and coefficient maps on the actual
representation and extension stacks, compatible with changes of
height. For compact \(S\times E\), the reduced complexes and
the common elliptic torsor must also be retained. A coordinate
volume on an affine chart specifies neither descent system by itself.

\chapter{The colored elliptic Ran space}
The elliptic degree is a global charge carried by the whole configuration.  Define
\[
 \Ran^{\mathrm{wr}}(E)=\coprod_{m\ge0}\Ran(E)_m
\]
with monoidal operation
\[
 (I,m)\sqcup(J,n)=(I\sqcup J,m+n)
\]
on disjoint supports.  The component $m=0$ is point-local.  The components $m>0$ are wrapped.

Let
\[
 \nIg=\bigoplus_{m\ge0}\nIg[m]
\]
be the elliptic-degree grading.  The root bracket obeys
\[
 [\nIg[m],\nIg[n]]\subset\nIg[m+n].
\]

\begin{construction}[Hybrid factorization algebra]
Form the chiral Chevalley factorization coalgebra
\[
 \mathcal F^{\mathrm{hyb}}_{\Delta,E}
 =C_*^{\mathrm{ch}}\!\left(\Cur_E(\nIg)\right).
\]
Retain the elliptic-degree grading on every Chevalley word and place its degree-$m$ summand on $\Ran(E)_m$.  Disjoint union adds the degree labels.
\end{construction}

\begin{theorem}[Local--wrapped factorization]
$\mathcal F^{\mathrm{hyb}}_{\Delta,E}$ is a graded factorization coalgebra on $\Ran^{\mathrm{wr}}(E)$.  Its degree-zero color is the current theory of $\nIg[0]$.  Every degree-$m$ color is a module for the degree-zero color.  The factorization product of degrees $m$ and $n$ lands in degree $m+n$.  Its chiral primitives are $\Cur_E(\nIg)$, and constant modes recover $\nIg$.
\end{theorem}
\begin{proof}
The chiral Chevalley construction is a factorization coalgebra for every chiral Lie algebra.  The bracket preserves elliptic-degree addition, so its differential and coproduct preserve the graded Ran placement.  The primitive theorem for the Chevalley coalgebra recovers the current Lie algebra.
\end{proof}

This construction retains the nonnegative root-degree label on a
point-supported current theory. It does not replace a sheaf projecting
onto the whole elliptic curve by a point-supported sheaf. The compact
wrapped comparison must map its common-torsor extension correspondences
to these graded current operations, preserving support and charge.
The bracket on the present carrier is the Borcherds bracket.

\chapter{Physical charge and Gram degree}
Let $\Gammphys$ be an additive physical charge lattice and let
\[
 \PiGram(c)=\left(\frac{Q^2}{2},Q\cdot P,\frac{P^2}{2}\right)
\]
be the Gram map.  Its polarization is
\[
 B(c,d)=\PiGram(c+d)-\PiGram(c)-\PiGram(d).
\]
Define the central extension
\[
 \widetilde\Gamma_\Pi=\Gammphys\times\Gammaind
\]
with multiplication
\[
 (c,\lambda)(d,\mu)
 =(c+d,\lambda+\mu+B(c,d)).
\]

\begin{theorem}[Graph homomorphism]
The map
\[
 j_\Pi:\Gammphys\longrightarrow\widetilde\Gamma_\Pi,
 \qquad
 c\longmapsto(c,\PiGram(c))
\]
is a group homomorphism.  For every $\Gammphys$-graded algebra $A$, the Rees lift
\[
 a_c\longmapsto a_cu^{\PiGram(c)}
\]
turns multiplication into an additive $\widetilde\Gamma_\Pi$-grading with cocycle $u^{-B(c,d)}$.
\end{theorem}
\begin{proof}
The defining multiplication gives
\[
 j_\Pi(c)j_\Pi(d)
 =(c+d,\PiGram(c)+\PiGram(d)+B(c,d))
 =j_\Pi(c+d).
\]
The algebra identity follows immediately.
\end{proof}

The universal Igusa category is graded directly by the Lorentzian root lattice.  A geometric $K3\times E$ theory enters through a charge functor to $\widetilde\Gamma_\Pi$.

\chapter{The universal protected field theory}
For each $H$ take cotangent BF theory with gauge Lie superalgebra $\nIg^{(H)}$ on $\R^2_{\mathrm{top}}\times E$.  Its fields are
\[
 \Omega^\bullet(\R^2)\widehattensor
 \Omega^{0,\bullet}(E)
 \widehattensor
 \left(\nIg^{(H)}[1]\oplus(\nIg^{(H)})^\vee[2]\right).
\]
The action is
\[
 S_H=\int_{\R^2\times E}
 \angles{B,(\dd+\dbar)A+\frac12[A,A]}.
\]

\begin{theorem}[Classical local theory and exact quantum harmonic sector]
For every finite height $H$ the shifted-cotangent action defines a classical mixed holomorphic--topological factorization algebra on $\R^2_{\mathrm{top}}\times E$.  Harmonic reduction on a compact product cell gives the finite-dimensional cotangent BV theory of $\nIg^{(H)}$, and this reduced theory satisfies the quantum master equation.  Complementary Lagrangian boundary polarizations give the boundary enveloping algebra $U(\nIg^{(H)})$, the two-sided bar complex $C_*(\nIg^{(H)})$, and the current envelope on $E$.  The proper-height inverse system is coefficientwise finite and defines a complete classical local theory together with an exact quantum harmonic sector.
\end{theorem}
\begin{proof}
The invariant cotangent pairing and the Jacobi identity give the classical master equation and factorization descent.  Harmonic reduction leaves a finite positive nilpotent Lie superalgebra.  Every even adjoint endomorphism raises proper height and is strictly upper triangular; every odd adjoint endomorphism has vanishing supertrace.  The reduced BV divergence therefore vanishes.  Canonical quantization in the selected polarization gives the enveloping algebra, and the augmentation boundary states give the two-sided bar, identified with Chevalley chains by PBW.  Each map preserves height, so the inverse system is defined coefficientwise.
\end{proof}


\chapter{Absolute convergence of the virtual root operator}
For each positive root $\alpha$, let $V_\alpha$ be the minimal parity realization of the signed multiplicity: it is purely even of dimension $\sdim\gIg_\alpha$ when this integer is nonnegative and purely odd of dimension $-\sdim\gIg_\alpha$ when it is negative.  Set
\[
 \Vcal_\Delta=\widehat\bigoplus_{\alpha>0}V_\alpha.
\]
This virtual root superspace is the canonical linear carrier of the denominator character; the Lie bracket remains carried by $\gIg$.

Let $h(\alpha)$ be the proper height.  Jacobi-form coefficient estimates give constants $C_0,C_1>0$ such that
\begin{equation}\label{eq:vol3-root-growth}
 \dim V_\alpha=\abs{\sdim\gIg_\alpha}\le C_0\exp(C_1h(\alpha))
\end{equation}
for every positive root.  The number of Gram triples of height $H$ is polynomial in $H$.

\begin{theorem}[Nonempty trace-class chamber]
Let $Z_0$ lie in the interior of the positive Weyl chamber.  For sufficiently large $t>0$,
\[
 \sum_{\alpha>0}\dim V_\alpha\,
 \abs{e^{-\alpha(tZ_0)}}<\infty.
\]
Hence the absolute-convergence chamber contains a conical tail in every interior direction.
\end{theorem}
\begin{proof}
Interior positivity gives a constant $c(Z_0)>0$ with
$\Re\alpha(Z_0)\ge c(Z_0)h(\alpha)$.
Combine this with \eqref{eq:vol3-root-growth}.  The contribution of roots of height $H$ is bounded by a polynomial in $H$ times
\[
 \exp\bigl((C_1-tc(Z_0))H\bigr).
\]
The series converges when $tc(Z_0)>C_1$.
\end{proof}

\begin{corollary}[Trace-class root Hamiltonian]
On the conical tail, the diagonal operator
\[
 K_Z\big|_{V_\alpha}=e^{-\alpha(Z)}\id
\]
is trace class on both parity summands.  The Fredholm product is holomorphic and equals the absolutely convergent Borcherds product.
\end{corollary}
\begin{proof}
The trace norm is the summable series of the preceding theorem, separated by parity.  Standard trace-class determinant theory gives holomorphy, and diagonalization gives the product over root eigenvalues.
\end{proof}


\chapter{The virtual root Fredholm operator}
Let
\[
 \mathcal H^{\mathrm{vir}}_\Delta
 =\widehat\bigoplus_{\alpha>0}V_\alpha
\]
with an orthonormal basis in each minimal parity realization.
This is the virtual root Hilbert space used in the preceding
absolute-convergence theorem.  In the absolute-convergence chamber define
\[
 K_Z\big|_{V_\alpha}=e^{-\alpha(Z)}\id,
 \qquad
 D_Z=1-K_Z.
\]

\begin{theorem}[Oriented Fredholm square root]
The operator $K_Z$ is trace class on the even and odd summands of the virtual root Hilbert space.  Its Fredholm Berezinian is
\[
 \sdet_F(D_Z)
 =\prod_{\alpha>0}(1-e^{-\alpha(Z)})^{\sdim\gIg_\alpha}.
\]
Choose an order of the positive roots, first by height and then lexicographically.  This choice orients every finite-height root space.  The skew doubling
\[
 \mathbb A_Z=
 \begin{pmatrix}0&D_Z\\-D_Z^{\mathsf t}&0\end{pmatrix}
\]
has an oriented Fredholm super-Pfaffian satisfying
\[
 \operatorname{sPf}_F(\mathbb A_Z)=\sdet_F(D_Z).
\]
Consequently
\[
 e^{-\rho}\operatorname{sPf}_F(\mathbb A_Z)=\Dmon(2Z).
\]
\end{theorem}
\begin{proof}
Absolute convergence of the Borcherds product is the trace-norm summability condition.  The determinant formula follows from the spectrum.  At finite height the Pfaffian of the oriented skew doubling is the determinant.  Trace-norm convergence passes the identity to the limit.  The denominator theorem supplies the Weyl shift.
\end{proof}

\section{Automorphic orientation}
The product defines a section of the automorphic line $\mathcal L^5\tensor\nu_5$.  On the absolute-convergence chamber, the oriented Fredholm Pfaffian is a trivialization of this line.  Analytic continuation uses the automorphy factor of $\Dmon$ and therefore glues the chamberwise Fredholm models with transition character $\nu_5$.

\begin{corollary}[Orientation character]
The automorphic orientation line of the root Dirac family is the sign local system $\nu_5$.  On every real simple reflection it has value $-1$.
\end{corollary}
\begin{proof}
The Fredholm product equals $D_5$ on the convergence chamber, and $D_5$ transforms with automorphy factor $\nu_5$.  This transformation law is the gluing rule for the automorphic orientation line.  A real simple root occurs with multiplicity one, so reflection reverses its oriented normal line.
\end{proof}

\chapter{The universal Igusa protected effective theory}
\begin{definition}[Universal Igusa completion]
The protected completion selected by the K3 half-index is the quintuple
\[
 \mathfrak K_{\Delta}=
 \left(
 \Gamma_\Delta,
 \mathcal B_\Delta,
 \mathcal H_{\mathrm{1p}},
 \mathcal F^{\mathrm{hyb}}_{\Delta,E},
 \mathbb A_Z
 \right).
\]
Its components are the pro-system of finite three-Calabi--Yau completions, the ordinary quiver Hall category, the Borcherds one-particle Lie superalgebra, the local--wrapped current factorization algebra, and the virtual root Dirac family.
\end{definition}

\begin{theorem}[Igusa protected completion]
The quintuple $\mathfrak K_\Delta$ satisfies:
\begin{enumerate}
\item its one-particle Lie superalgebra is $\gIg$;
\item its ordinary Hall presentation and Hong's quantum supergroup retain
the stated Cartan grading; comparison of their relation ideals and
parity-refined PBW spaces requires an additional map;
\item its chiral primitive algebra on $E$ is the local--wrapped current algebra of $\gIg$;
\item its raw multiparticle supercharacter is $P^{-1}$, where $e^{-\rho}P=\DenIg$; equivalently $e^{\rho}\sch\Fock(\nIg)=\DenIg^{-1}$;
\item the half-coordinate pullback of the doubled Weyl-normalized character is $-\ChiTen^{-1}$ in the Oberdieck--Pixton normalization;
\item its oriented first-order determinant is $\DenIg$ and its orientation monodromy is $\nu_5$.
\end{enumerate}
\end{theorem}
\begin{proof}
The first three statements are the Borcherds root, quiver Hall, Hong quantum-supergroup, Ginzburg, and colored-Ran theorems.  PBW gives the fourth.  The identity $\ChiTen=\Dmon^2$ and the established reduced scalar theorem give the fifth.  The Fredholm theorem gives the sixth.
\end{proof}

The tuple retains algebraic constructions associated with the Igusa
presentation and their scalar characters. A geometric type-IIA or
topological-string realization must construct compatible functors
from its oriented extension category to these targets. Primitive
states and extension amplitudes are inputs to those functors.
Their action on morphisms, convolution, orientations, and chiral
operations must also be specified and proved compatible.

\chapter{The primitive comparison and its extensions}
Let \(\mathfrak p\) be a positive-root-graded Lie superalgebra
with finite-dimensional homogeneous pieces and a proper positive
height. A primitive realization of the Igusa presentation is a
parity-preserving homogeneous map from its labelled simple-root
spaces to \(\mathfrak p\), satisfying the defining relations.

\begin{theorem}[Universal presentation map]
\label{igusa:primitive-presentation-map}
Let \(\mathfrak p\) be a positive-root-graded Lie superalgebra
over \(\mathbb C\), with finite-dimensional homogeneous pieces.
Suppose a homogeneous parity-preserving map from every simple-root
generating space of \(\nIg\) to \(\mathfrak p\) satisfies
the defining Borcherds--Serre relations. It extends uniquely to
a graded Lie-superalgebra map
\[
 f:\nIg\longrightarrow\mathfrak p.
\]
It induces continuous maps on height completions, enveloping
algebras, Chevalley chains, and level-zero currents. The map \(f\)
is an isomorphism if and only if it is surjective and the root-graded
PBW series agree in each parity.
\end{theorem}
\begin{proof}
Extend the chosen maps to the free Lie superalgebra on the
simple-root spaces. The relation hypothesis makes the defining
ideal lie in its kernel, giving \(f\). The generators prove
uniqueness. On enveloping algebras the map sends
\(x_1\cdots x_j\) to \(f(x_1)\cdots f(x_j)\).
On Chevalley chains it is the graded symmetric map on the shifted
generators. It commutes with the Chevalley differential because
\(f\) preserves brackets. The current map is
\(\partial^j x\mapsto\partial^j f(x)\), preserving
\([x_\lambda y]=[x,y]\). Root degree preservation gives
compatible maps at each height and hence on the stated limits.

If \(f\) is surjective, \(U(f)\) is surjective. Proper
positive height makes each root-graded PBW component
finite-dimensional. Equality of its even and odd dimensions makes
\(U(f)\) bijective in every component. PBW embeds each Lie
superalgebra into its enveloping algebra, so \(f\) is injective.
The converse follows by applying the constructions to its inverse.
\end{proof}


A geometric comparison also requires functors on the actual Hall
and critical categories, a compatible coproduct and Hopf pairing,
and maps of determinant lines and operator domains. These categories
and lines are not the domain of the primitive Lie map. The
additional maps must be constructed and checked against that map.
They are not consequences of its universal property.

\chapter{The eight diagonal-divisor forms}
Cl\'ery and Gritsenko classify eight genus-two modular forms whose divisor is the diagonal with order one \cite{GC}.  Gritsenko--Nikulin and the CHL denominator constructions give the sourced reflective presentations used for the corresponding rows \cite{GN97,GN98a,GN98b,GovindarajanKrishna2009,Govindarajan2011,GovindarajanSamanta2019}.  Every row has a minimal virtual product superspace obtained from its Fourier exponents, together with a Fock determinant and a diagonal Fredholm operator.  A sourced denominator presentation refines this product model by its Borcherds--Cartan bracket, Weyl group, and quantum supergroup.  The half-integral rows are formed on their multiplier covers.
\[
\begin{array}{c|c|c|c}
\text{form}&\text{group}&\text{weight}&\text{character order}\\ \hline
\Delta_5&\operatorname{Sp}_4(\Z)&5&2\\
\Delta_2&\Gamma_2&2&4\\
\Delta_1&\Gamma_3&1&6\\
\Delta_{1/2}&\Gamma_4&1/2&8\\
\nabla_3&\Gamma^{(2)}_0(2)&3&2\\
\nabla_2&\Gamma^{(2)}_0(3)&2&2\\
\nabla_{3/2}&\Gamma^{(2)}_0(4)&3/2&4\\
Q_1&\Gamma_2(2)&1&4
\end{array}
\]

\begin{theorem}[Eight-row algebraic realization]
Each diagonal-divisor row has a proper-height virtual root superspace, a super-Fock determinant, a colored abelian current theory, and a trace-class Fredholm denominator operator whose character is the corresponding automorphic product.  For every row equipped with its reflective denominator presentation, the finite-window Borcherds quiver, ordinary Hall category, three-Calabi--Yau completion, and parity-refined quantum supergroup refine this product model and recover the same denominator character.
\end{theorem}
\begin{proof}
Local finiteness of the product exponents gives the graded superspace and its proper-height completion.  The Fock and Fredholm constructions depend only on these signed multiplicities and reproduce the product coefficientwise.  A reflective denominator presentation supplies a symmetrizable Borcherds superdatum; the quiver, Hall, Ginzburg, current, and Hong constructions then apply to that datum and preserve its PBW denominator.
\end{proof}

\chapter{Algebraic and physical closure}
\begin{theorem}[Universal Igusa realization theorem]
The constructions have the following inputs and outputs. An arrow
labelled by a construction records its dependence on that input:
\[
\begin{gathered}
 \phiweak\xmapsto{\mathrm{Borch}}\Dmon,
 \qquad
 \phiweak\xmapsto{\mathrm{root\ datum}}\gIg,
 \qquad
 A_\Delta\xmapsto{\mathrm{quiver}}Q_{\Delta,H}
 \xmapsto{\mathrm{Hall}}\mathcal B_\Delta,\\
 kQ_{\Delta,H}\xmapsto{\Pi_3}\Gamma_{\Delta,H},
 \qquad
 \gIg\xmapsto{\Cur_E}\mathcal F^{\mathrm{hyb}}_{\Delta,E},
 \qquad
 \gIg\xmapsto{\mathrm{Hong}}\widehat U_q(\gIg),\\
 \gIg\xmapsto{\mathrm{root\ spectrum}}\mathbb A_Z,
 \qquad
 e^{-\rho}\operatorname{sPf}_F(\mathbb A_Z)=\DenIg(Z)=\Dmon(2Z),\\
 \mathcal Z^{\red}_{K3\times E}=-(\jmath^*\DenIg)^{-2}=-\ChiTen^{-1}.
\end{gathered}
\]
The system contains parity-refined one-particle spaces, an ordinary Hall extension product, the Borcherds root bracket, a local--wrapped collision law, a quantum supergroup, an automorphic orientation line, a first-order Fredholm determinant, and the reduced scalar square.
\end{theorem}
\begin{proof}
The Borcherds lift, Gritsenko--Nikulin presentation, and primitive
reduced curve-counting theorem give the scalar identities. The
enveloping and Chevalley maps are induced by the presentation.
The quiver Hall algebra and finite Calabi--Yau completion have their
respective quiver and path-algebra inputs. Their common grading
specifies the degrees a comparison must preserve. A map between
those constructions requires the additional functors, coefficient
maps, and pairings described above.
\end{proof}


\part{Technical appendices to the preceding book}
\chapter{Exact theta-series calculation}\label{app:theta-computation}
Set $t=q^{1/8}$ and $x=r^{1/2}$. Then
\[
\begin{aligned}
\vartheta_2(\tau,z)&=\sum_{n\in\Z}t^{4n^2+4n+1}x^{2n+1},\\
\vartheta_3(\tau,z)&=\sum_{n\in\Z}t^{4n^2}x^{2n},\\
\vartheta_4(\tau,z)&=\sum_{n\in\Z}(-1)^nt^{4n^2}x^{2n}.
\end{aligned}
\]
Squaring numerator and denominator, dividing in $\Q[x^{\pm1}][[t]]$, and summing the three quotients gives the coefficient table in Chapter~1. The calculation uses integer convolution and formal-series inversion; fractional powers cancel before conversion from $(t,x)$ to $(q,r)$.

\chapter{The root reflection matrices}
In the basis $(f_2,f_3,f_{-2})$, the form is
\[
 G=\begin{pmatrix}0&0&-1\\0&2&0\\-1&0&0\end{pmatrix}.
\]
The three reflection matrices are
\[
 S_1=\begin{pmatrix}1&4&4\\0&-1&-2\\0&0&1\end{pmatrix},\quad
 S_2=\begin{pmatrix}1&0&0\\-2&-1&0\\4&4&1\end{pmatrix},\quad
 S_3=\begin{pmatrix}1&0&0\\0&-1&0\\0&0&1\end{pmatrix}.
\]
Direct multiplication gives
\[
 S_i^2=I,
 \qquad S_i^{\mathsf T}GS_i=G.
\]

\chapter{Automorphic, homological, and Hall comparisons}
The half-index determinant, theta normalization, Borcherds denominator, reduced scalar theorem, charge Gram map, and geometric Hall comparison datum belong to the following constructions:
\begin{enumerate}
\item the automorphic square root is completed entirely within the Borcherds product theorem;
\item the BKM bracket comes from the Gritsenko--Nikulin algebra;
\item the bar--Chevalley comparison gives the denominator a homological source;
\item the current and BF constructions give an exact first-order algebraic model;
\item orientation data refine determinant lines, while an operator Pfaffian requires an actual Fredholm family;
\item Weyl reflections are attached to real roots, and imaginary roots enter through multiplicities;
\item the compact Hall comparison uses additive Mukai charge, the polarized central extension, and hybrid local--wrapped sectors.
\end{enumerate}

\cleardoublepage
\phantomsection
\addcontentsline{toc}{part}{Calabi--Yau Quantum Groups}
\chapter*{Book entrance: Calabi--Yau Quantum Groups}
\addcontentsline{toc}{chapter}{Book entrance: Calabi--Yau Quantum Groups}
\begin{summarybox}
An extension of geometric objects defines a correspondence. Its action on cohomology depends on the coefficient theory and the allowed pushforward. Tensor products and exchange require a coproduct and compatible intertwiners. A trace on a specified representation then produces an amplitude. The reconstruction problem is to retain these operations through the comparisons between Calabi--Yau categories, boundary fields, Hall algebras, and quantum groups.
\end{summarybox}
\part{Categorical and field-theoretic foundations}

\chapter{Calabi--Yau categories}
\section{Smoothness, properness, and the Calabi--Yau class}
Let $\Ccal$ be a small stable $\kfield$-linear category. Smoothness means that the diagonal bimodule is perfect. Properness means that the morphism complexes are perfect over $\kfield$. A $d$-Calabi--Yau structure is a nondegenerate negative-cyclic class
\[
 \theta_{\CY}\in HC^-_d(\Ccal)
\]
whose underlying Hochschild trace identifies the diagonal bimodule with its dual shifted by $d$.

For $E\in\Ccal$, the tangent complex of the derived moduli stack of objects is
\[
 T_E\M_{\Ccal}\simeq\RHom_{\Ccal}(E,E)[1].
\]
The Calabi--Yau trace pairs two tangent vectors by composition and trace.

\begin{theorem}[The two universal categorical outputs]
Let $\Ccal$ be a smooth proper $d$-Calabi--Yau category.
\begin{enumerate}
\item The Hochschild cochains $CC^*(\Ccal)$ carry a framed $\Etwo$-algebra structure.
\item Whenever the moduli stack of objects is representable in the PTVV setting, $\M_{\Ccal}$ carries a $(2-d)$-shifted symplectic form.
\end{enumerate}
\end{theorem}
\begin{proof}
The cyclic Deligne theorem upgrades the brace $\Etwo$ structure on Hochschild cochains by the circle action supplied by the cyclic trace \cite{BravRozenblyum}. The shifted form at $E$ is
\[
 \RHom(E,E)[1]^{\tensor2}
 \xrightarrow{\ \circ\ }
 \RHom(E,E)[2]
 \xrightarrow{\ \Tr\ }
 \kfield[2-d].
\]
Nondegeneracy is the Calabi--Yau duality isomorphism, and closedness is the negative-cyclic lift \cite{PTVV,BravDyckerhoff}.
\end{proof}

The integer $d$ controls the shift of the moduli form. The Hochschild operation remains framed $\Etwo$. These structures meet through deformation theory while retaining their separate operadic degrees.

\chapter{Local BV theories and factorization observables}
\section{Cyclic elliptic local Lie algebras}
Let $X$ be a complex manifold. A classical local field theory is specified by a sheaf $\Lcal$ of elliptic $L_\infty$ algebras, a local invariant pairing, and an action functional $S$ satisfying the classical master equation. For an open set $U\subset X$, the classical observables are
\[
 \Obs^{\mathrm{cl}}(U)=C^*_{\CE,\mathrm{cont}}(\Lcal(U)).
\]
Here $\Lcal(U)$ consists of unrestricted smooth fields with its
nuclear smooth-section topology. Continuous cochains use the completed
symmetric powers of its continuous dual, with the completed tensor
products of the elliptic-observable construction \cite{CG1,CG2}.
The dual consists of compactly supported distributions.
For $U\subset V$, restriction of fields $\Lcal(V)\to\Lcal(U)$
induces the covariant map on observables.
For pairwise disjoint $U_i\subset V$, restrict to
$\bigoplus_i\Lcal(U_i)$ and multiply the pulled-back cochains.
These are the prefactorization maps.

\begin{theorem}[Factorization observables]
An elliptic classical BV theory produces a factorization algebra $\Obs^{\mathrm{cl}}$. A renormalized quantization satisfying the quantum master equation produces a factorization algebra $\Obs^{\mathrm{q}}$ over $\kfield[[\hbar]]$.
\end{theorem}
\begin{proof}
Restriction is a morphism of local $L_\infty$ algebras, so its
contravariant continuous Chevalley map commutes with the differential.
For disjoint opens, the completed cochains of the finite direct sum
are the completed tensor product of the cochains. This constructs
the displayed maps and their composition identities.
The elliptic descent theorem in this completed coefficient category
supplies Weiss descent \cite{CG1,CG2}. In the quantum theory, renormalized time-ordered products define the structure maps, and the quantum master equation makes their differential square-zero and compatible with nested ultraviolet contractions \cite{CG1,CG2}.
\end{proof}

\section{Boundary and defect restriction}
Let $Y=C\times\R_{\ge0}$, with $C$ a complex curve, and let a boundary condition be given at $C\times\{0\}$. Boundary observables are holomorphic along $C$ and ordered along the inward normal line. A compact boundary object $b$ in the boundary factorization category gives
\[
 A_b=\RHom(b,b).
\]

\begin{theorem}[Ordered boundary algebra]
When boundary observables are holomorphic on $C$ and locally constant along the oriented normal direction, $A_b$ is an ordered chiral algebra
\[
 A_b\in\Alg_{\Eone}\bigl(\Fact_D(C),\tensor^!\bigr).
\]
If $b$ is a local compact generator and the Morita equivalences obey factorization descent, the boundary category is equivalent to $A_b$-modules in factorization coefficients.
\end{theorem}
\begin{proof}
Composition of boundary endomorphisms gives the internal $\Eone$ multiplication. Holomorphic collision gives the factorization coefficient on $C$. The componentwise locality identity makes multiplication compatible with every disjoint-support factorization map. Local compact-generator Morita theory descends over the Ran diagrams. These are the recognition and Morita theorems of Volume~I.
\end{proof}


\chapter{Three curve-level algebraic outputs}
Let $C$ be a smooth curve and let $\Dcal_C$ be the fixed Ran category.  The same geometric source can produce three algebraic outputs, and their distinction is functorial.

\begin{definition}[Curve-level faces]
A transverse face is
\[
 A^\perp\in\Alg_{\Eone}(\Fact_D(C),\tensor^!).
\]
An associative chiral face is
\[
 A^{\ch}_{\mathsf{Ass}}\in\Alg_{\Eone}(\Dcal_C,\starCh).
\]
A chiral Lie face is
\[
 \g^{\ch}\in\Alg_{\mathsf{Lie}}(\Dcal_C,\starCh).
\]
Their bars are $\BarPerp$, $\BarAssCh$, and $\BarLieCh$.
\end{definition}

\begin{theorem}[Source separation]
A boundary normal direction determines the transverse face.  A chiral-convolution multiplication determines the associative chiral face.  A Lie bracket for chiral convolution determines the Lie face.  The commutator functor sends the associative chiral face to the Lie face.  A chiral enveloping algebra $U^{\ch}(\g)$ satisfies
\[
 \BarAssCh(U^{\ch}(\g))\simeq\BarLieCh(\g)
\]
on the nilcomplete PBW locus.  A mixed source joins the transverse and associative chiral faces through an aligned distributive law.
\end{theorem}
\begin{proof}
Each statement is the corresponding operadic restriction or adjunction in the named monoidal category.  The PBW equivalence is the enveloping lane of chiral Koszul duality.  The mixed statement is the distributive-law theorem of Volume~I.
\end{proof}

This theorem replaces every dimension-based dispatch of an $\En$ level.  The geometry of the source selects the available directions; the operadic level records the operations actually constructed.

\chapter{Holomorphic Chern--Simons theory}
Let $X$ be a complex threefold with holomorphic volume form $\Omega_X$, and let $\mathfrak a$ be a quadratic Lie algebra. The BV fields are
\[
 \mathcal E_X=\Omega^{0,*}(X,\mathfrak a)[1],
\]
with degree $-1$ pairing
\[
 \omega(A,B)=\int_X\Omega_X\wedge\angles{A,B}.
\]
The action is
\[
 S_{\hCS}(A)=\int_X\Omega_X\wedge
 \left(\frac12\angles{A,\dbar A}+\frac16\angles{A,[A,A]}\right).
\]

\begin{theorem}[Classical hCS factorization algebra]
The action $S_{\hCS}$ satisfies the classical master equation. Its classical observables form a holomorphic factorization algebra on $X$.
\end{theorem}
\begin{proof}
The Hamiltonian vector field is the Maurer--Cartan differential $\dbar A+\frac12[A,A]$. Its square is the Dolbeault square, the derivation identity, and the Jacobi identity. The elliptic Dolbeault complex then gives factorization observables by the preceding theorem.
\end{proof}

A perturbative quantization is controlled by the local BV deformation complex. In complex dimension three, its one-loop gauge anomaly occupies the degree-four invariant-polynomial summand; mixed and gravitational components enter when the background fields carry the corresponding couplings. A quantization is a trivialization of the complete local anomaly class, compatible with restriction and factorization \cite{CostelloLiAnomaly}.

\chapter{Pushforward to a curve}
Let $p:X\to C$ be a holomorphic map to a complex curve. For a factorization algebra $\F$ on $X$, define
\[
 (p_*\F)(U)=\F(p^{-1}U).
\]

\begin{theorem}[Factorization pushforward]
The assignment $p_*\F$ is a factorization algebra on $C$. Suppose in addition that $p$ is proper and locally trivial in a topology compatible with the factorization theory, that $\F$ is locally constant in the fiber directions, and that product-chart excision and base change hold. Then on every trivializing open $U\subset C$ with fiber $F_p$ there is a natural equivalence
\[
 (p_*\F)(U)\simeq
 \int_{F_p}\F|_{F_p\times U},
\]
compatible on overlaps. If the resulting object is holomorphic and regular holonomic along $C$, Riemann--Hilbert algebraization gives a factorization $D$-module on $C$.
\end{theorem}
\begin{proof}
Disjoint opens in $C$ have disjoint inverse images, so the factorization products push forward. Weiss descent is inherited from inverse images of Weiss covers. On a product chart, local constancy in the fiber and factorization excision identify the value with fiberwise factorization homology. The base-change hypothesis glues these local equivalences. Regular-holonomic Riemann--Hilbert transports the holomorphic factorization object to the algebraic $D$-module setting.
\end{proof}

Thus a curve shadow is attached to an actual map and an actual field theory. Different fibrations and boundary conditions produce different ordered chiral algebras from the same categorical source.

\chapter{The derived centre and the physical bulk}
For an ordered boundary algebra $A$, define
\[
 Z_{\ord}(A)=\RHom_{A\tensor A^{\op}}(A,A).
\]
It is a canonical $\Etwo$ algebra. A bulk theory acting on the boundary gives
\[
 \beta:\Obs_{\mathrm{bulk}}|_{\partial}\longrightarrow Z_{\ord}(A).
\]

\begin{definition}[Open--closed completeness]
A boundary condition is open--closed complete on a protected sector when $\beta$ is an equivalence on that sector.
\end{definition}

The definition turns the phrase ``bulk equals centre'' into a comparison theorem. The centre exists algebraically for every ordered boundary algebra; the physical bulk is identified with it by the open--closed map.

\part{Hall and BPS algebras}

\chapter{Extension correspondences}
Let $\M=\coprod_{\gamma\in\Gamma}\M_\gamma$ be a moduli stack in an exact or derived category. Let $\M^{\mathrm{ext}}_{\alpha,\beta}$ classify extensions
\[
 0\longrightarrow E_\beta\longrightarrow E\longrightarrow E_\alpha\longrightarrow0.
\]
The correspondence is
\[
 \M_\alpha\times\M_\beta
 \xleftarrow{\ p\ }
 \M^{\mathrm{ext}}_{\alpha,\beta}
 \xrightarrow{\ q\ }
 \M_{\alpha+\beta}.
\]

\begin{theorem}[Hall associativity]
Assume a coefficient theory in which $p^*$ and $q_*$ are defined and satisfy base change for the two-step filtration stack. Then
\[
 a*b=q_*p^*(a\boxtimes b)
\]
defines an associative product.
\end{theorem}
\begin{proof}
The two parenthesizations are the two projections from the stack of filtrations $0\subset E_1\subset E_2\subset E_3$. Base change identifies both composites with pullback to this stack followed by pushforward to $E_3$.
\end{proof}

For critical CoHAs, the coefficient theory includes vanishing cycles and a coherent orientation. For symplectic CoHAs, it includes the selected equivariant Borel--Moore theory. Completion is chosen so that each output degree receives a convergent sum.

\chapter{Orientation data and the threefold Hall product}
For compact Calabi--Yau threefolds, natural orientation data on moduli of coherent sheaves and perfect complexes is constructed by Joyce and Upmeier \cite{JoyceUpmeier}.  A general cohomological Hall algebra for a $3$-Calabi--Yau category uses strong orientation data, which carries the coherent multiplicativity required by the attractor and extension correspondences \cite{KPS}.

\begin{theorem}[Threefold Hall algebra from strong orientation]
Let $\Ccal$ be a $3$-Calabi--Yau category satisfying the geometric hypotheses of Kinjo--Park--Safronov, and assume strong orientation data.  Then its cohomological Donaldson--Thomas complexes carry an associative cohomological Hall product.  The product is the parabolic-induction correspondence supplied by the attractor Lagrangian correspondence.
\end{theorem}
\begin{proof}
Kinjo--Park--Safronov prove the functoriality of the Donaldson--Thomas perverse sheaf under the attractor Lagrangian correspondence and use strong orientation to define the parabolic-induction morphism.  Their base-change and associativity theorem gives the CoHA \cite{KPS}.
\end{proof}

\begin{definition}[Reduced $K3\times E$ Hall datum]
A reduced $K3\times E$ Hall datum is a strong-oriented threefold Hall datum together with descent along the elliptic translation action, a reduced deformation theory, and a separated completion in physical charge.  A full quantum-group datum additionally includes a coproduct, Cartan extension, negative half, and nondegenerate continuous Hopf pairing.
\end{definition}

\chapter{The geometric \texorpdfstring{$E$}{E}-equivariant Hall source}
Let $S$ be a projective K3 surface, let $E$ be an elliptic curve, and put
\[
 X=S\times E.
\]
Write $\Coh_{\leq1}(X)$ for coherent sheaves whose support has dimension at most one.  This is an extension-closed abelian category.  Its numerical charge group is denoted $\Gamma_X$, and the degree of the projected one-cycle on $E$ is the additive function
\[
 m:\Gamma_X\longrightarrow\Z_{\geq0}.
\]
For $\gamma\in\Gamma_X$, let $\M_\gamma$ be the moduli stack of objects of charge $\gamma$.  Translation on the elliptic factor gives an action of $E$ on every $\M_\gamma$ and on every extension stack.

Oberdieck and Shen construct the $E$-equivariant Grothendieck ring of stacks and its Hall convolution.  In the present notation the completed algebra is
\[
 \Hcal_X^{E,\mathrm{mot}}
 =\widehat\bigoplus_{\gamma\in\Gamma_X}
 K_0^E(\operatorname{St}/\M_\gamma),
\]
where the completion is taken in effective curve class and Euler characteristic.

\begin{theorem}[Geometric reduced Hall algebra]
The extension correspondence in $\Coh_{\leq1}(X)$ defines an associative $E$-equivariant motivic Hall product on $\Hcal_X^{E,\mathrm{mot}}$.  The Oberdieck--Shen reduced integration morphism
\[
 I_X^E:\Hcal_X^{E,\mathrm{mot}}
 \longrightarrow
 \widehat{\Q[\Gamma_X]}[\epsilon]/(\epsilon^2)
\]
is multiplicative on the regular subalgebra on which their integration theorem is defined.  The coefficient of $\epsilon$ is the reduced contribution of the positive-dimensional $E$-orbits.
\end{theorem}
\begin{proof}
The stack of short exact sequences is $E$-equivariant because translation preserves exactness and numerical charge.  The stack of two-step filtrations identifies the two composites of three Hall products, so equivariant base change proves associativity.  The integration statement is the $E$-equivariant motivic Hall theorem of Oberdieck--Shen \cite{OberdieckShen}.  Their stabilizer stratification separates fixed contributions from moving elliptic orbits and packages the latter in the square-zero parameter $\epsilon$.
\end{proof}

\chapter{Local and wrapped geometry inside the Hall algebra}
Define the completed degree-zero and positive-degree pieces
\[
 \Hcal_X^{\mathrm{loc}}
 =\widehat\bigoplus_{m(\gamma)=0}
 K_0^E(\operatorname{St}/\M_\gamma),
 \qquad
 \Hcal_X^{\mathrm{wr}}
 =\widehat\bigoplus_{m(\gamma)>0}
 K_0^E(\operatorname{St}/\M_\gamma).
\]
A pure one-dimensional support of degree zero projects to a finite subset of $E$.  A support of positive degree projects onto $E$.  The first sector is point-local on the elliptic curve; the second is wrapped.

\begin{theorem}[Geometric local--wrapped decomposition]
Hall multiplication respects elliptic degree:
\[
 \Hcal_{m_1}*\Hcal_{m_2}\subseteq\Hcal_{m_1+m_2}.
\]
Consequently $\Hcal_X^{\mathrm{loc}}$ is a subalgebra, $\Hcal_X^{\mathrm{wr}}$ is a closed two-sided ideal, and
\[
 \Hcal_X^{E,\mathrm{mot}}
 \cong
 \Hcal_X^{\mathrm{loc}}\ltimes\Hcal_X^{\mathrm{wr}}
\]
as a complete vector space with multiplication determined by the Hall correspondences.  The reduced integration morphism preserves this decomposition.
\end{theorem}
\begin{proof}
The cycle of the middle term of a short exact sequence is the sum of the cycles of its subobject and quotient.  Projection to $E$ therefore adds degrees.  Degree zero is closed under addition, and positive degree absorbs addition by a nonnegative degree.  The same grading is present on the target group algebra of $I_X^E$.
\end{proof}

This theorem supplies the geometric version of the local--wrapped splitting constructed from the Borcherds root algebra in the companion volume.  The two constructions have the same additive elliptic degree and become comparable after a primitive root framing is chosen.

\chapter{Cohomological refinement and the Gram character}
A cohomological refinement uses the Donaldson--Thomas sheaf on the derived moduli stack.  The datum required for the convolution is the strong orientation coherence of Kinjo--Park--Safronov.  A strong orientation enhancement of the natural Joyce--Upmeier orientation therefore produces
\[
 \Hcal_X^{E,\mathrm{coh}}
 =\widehat\bigoplus_{\gamma\in\Gamma_X}
 H_*^{\mathrm{BM},E}(\M_\gamma,\Phi_\gamma),
\]
with its cohomological Hall product.

\begin{theorem}[Strong-oriented cohomological lift]
For every strong orientation enhancement on the $E$-equivariant derived moduli problem, the attractor correspondence constructs an associative cohomological Hall algebra $\Hcal_X^{E,\mathrm{coh}}$.  Elliptic degree again gives a subalgebra in degree zero and a closed ideal in positive degree.  Forgetting the cohomological coefficient system recovers the corresponding motivic Hall correspondences.
\end{theorem}
\begin{proof}
The parabolic-induction theorem of Kinjo--Park--Safronov gives the cohomological product from the strong-oriented attractor Lagrangian correspondence \cite{KPS}.  Numerical charge and elliptic degree are constant on connected components of every correspondence, so the proof of the preceding decomposition applies verbatim.
\end{proof}

For the dyonic charge lattice
\[
 \Gamma_{\mathrm{phys}}
 =\widetilde H(S,\Z)\oplus\widetilde H(S,\Z),
\]
write
\[
 \Pi(Q,P)=\left(\frac{Q^2}{2},Q\cdot P,\frac{P^2}{2}\right).
\]
To compare sheaf charges with this lattice, choose an additive map
\[
 c:\Gamma_X\longrightarrow\Gamma_{\mathrm{phys}}.
\]
Put $B_\Pi(a,b)=\Pi(a+b)-\Pi(a)-\Pi(b)$ and give
$\widetilde\Gamma_{\mathrm{phys}}=\Gamma_{\mathrm{phys}}\times\mathbb Z^3$
the group law
\[
 (a,u)(b,v)=(a+b,u+v+B_\Pi(a,b)).
\]
Then $j_\Pi(a)=(a,\Pi(a))$ is a homomorphism.
Assume $j_\Pi c$ preserves the selected charge completions.
On the regular subalgebra $\mathcal H_{\mathrm{reg}}$ of the reduced
integration theorem, define
\[
 I_{X,\Pi,c}^E=(j_\Pi c)_*I_X^E:
 \mathcal H_{\mathrm{reg}}\longrightarrow
 \widehat{\mathbb Q[\widetilde\Gamma_{\mathrm{phys}}]}
 [\epsilon]/(\epsilon^2).
\]
Write $I_{X,\Pi}^E$ for this map when $c$ has been fixed.

\begin{proposition}[Polarized reduced character]\label{cyiv:polarized-character}
The map $I_{X,\Pi,c}^E$ is multiplicative.
Its coefficients are the reduced invariants supplied by $I_X^E$,
regraded by $j_\Pi c$.
If $g\in1+F^1H$ is group-like in a complete filtered Hall bialgebra
and the logarithm converges, then $\log g$ is primitive.
\end{proposition}
\begin{proof}
Expanding $\Pi(a+b+d)$ gives
\[
 B_\Pi(a,b)+B_\Pi(a+b,d)=B_\Pi(b,d)+B_\Pi(a,b+d).
\]
This proves associativity of the group law.
The definition gives $j_\Pi(a+b)=j_\Pi(a)j_\Pi(b)$.
Composing its group-algebra map with reduced integration proves
multiplicativity on $\mathcal H_{\mathrm{reg}}$.
For the last assertion, the two factors in
$\Delta g=(g\widehat\otimes1)(1\widehat\otimes g)$ commute.
Taking the convergent logarithm gives
$\Delta\log g=\log g\widehat\otimes1+1\widehat\otimes\log g$.
\end{proof}

The primitive reduced Gromov--Witten theorem supplies the separate
character $-\chi_{10}^{-1}$ in classes primitive in the K3 factor
\cite{OP}.
Comparing this character with the Hall character requires the same
charge map, sector, and normalization.
A primitive root framing must also preserve the extension bracket
and parity of the chosen BPS states.

\chapter{BPS primitives}
A connected graded bialgebra has primitive Lie algebra
\[
 \Prim(H)=\set{x:\Delta x=x\tensor1+1\tensor x}.
\]
In cohomological Donaldson--Thomas theory, the BPS sheaf isolates the primitive part of the Hall algebra.

\begin{theorem}[Two-Calabi--Yau BPS algebra]
For the geometric two-Calabi--Yau categories covered by the BPS-algebra theorem, the BPS algebra is the positive enveloping algebra of a generalized Kac--Moody Lie algebra. This scope includes semistable sheaves on K3 and other symplectic surfaces under the stated good-moduli hypotheses.
\end{theorem}
\begin{proof}
The theorem is the Davison--Hennecart--Schlegel Mejia construction \cite{DHS}. The simple BPS sheaves give the simple-root spaces, extension correspondences give the bracket, and the integrality/PBW theorem identifies the Hall algebra with the positive enveloping algebra.
\end{proof}

This result supplies a global K3 generalized Kac--Moody algebra at the two-Calabi--Yau Hall level. Yangian structures arise from additional equivariant correspondences and coproducts in the surface and quiver models described below.

\chapter{The integral Hall--Borcherds amalgam}
Let $\mathfrak p_X^+$ be a complete charge-graded BPS Lie superalgebra supplied by a geometric Hall source, and let $\nDelta$ be the positive Igusa Borcherds Lie superalgebra.  The polarized charge extension places both gradings over one abelian Cartan algebra $\mathfrak h_\Pi$.  Form the two complete Borels
\[
 \mathfrak b_X=\mathfrak h_\Pi\ltimes\mathfrak p_X^+,
 \qquad
 \mathfrak b_\Delta=\mathfrak h_\Pi\ltimes\nDelta.
\]
The Cartan acts through the physical charge on the first Borel and through Gram degree on the second.

\begin{definition}[Integral protected Lie algebra]
The integral Hall--Borcherds algebra is the completed amalgamated coproduct
\[
 \mathfrak b_{X\boxtimes\Delta}
 =\mathfrak b_X
  \widehat\amalg_{\mathfrak h_\Pi}
  \mathfrak b_\Delta
\]
in complete charge-filtered dg Lie superalgebras.
\end{definition}

\begin{theorem}[Hall--Borcherds amalgamation]
The completed amalgam exists and represents pairs of continuous Lie morphisms from $\mathfrak b_X$ and $\mathfrak b_\Delta$ that agree on $\mathfrak h_\Pi$.  The Cartan projections of the two Borels induce canonical retractions
\[
 r_X:\mathfrak b_{X\boxtimes\Delta}\longrightarrow\mathfrak b_X,
 \qquad
 r_\Delta:\mathfrak b_{X\boxtimes\Delta}\longrightarrow\mathfrak b_\Delta.
\]
Its completed enveloping algebra is the amalgamated free product
\[
 \widehat U(\mathfrak b_{X\boxtimes\Delta})
 \cong
 \widehat U(\mathfrak b_X)
 \widehat *_{\widehat U(\mathfrak h_\Pi)}
 \widehat U(\mathfrak b_\Delta).
\]
The geometric Hall operations survive under $r_X$, while the Igusa bracket, PBW character, and automorphic determinant survive under $r_\Delta$.
\end{theorem}
\begin{proof}
Construct the ordinary Lie coproduct from the free product of the two Borels and divide by the ideal identifying their two Cartan copies.  Complete by total positive charge and proper root height.  Every fixed filtered degree uses finitely many Lie words, so the completion is separated and has the stated mapping property.  The projection $\mathfrak b_X\to\mathfrak h_\Pi\to\mathfrak b_\Delta$ together with the identity on $\mathfrak b_\Delta$ gives $r_\Delta$ by universality; the second retraction is symmetric.  The enveloping functor preserves coproducts of Lie algebras, and PBW identifies the completed algebraic filtration with the displayed completed free product.
\end{proof}

\begin{definition}[Comparison ideal]
A primitive comparison $\varphi:\nDelta\to\mathfrak p_X^+$ preserving Cartan degree defines the closed ideal
\[
 \mathcal J_\varphi
 =\overline{\angles{
   i_\Delta(x)-i_X(\varphi(x))\mid x\in\nDelta
 }}
 \subset\mathfrak b_{X\boxtimes\Delta}.
\]
The diagonal realization is
\[
 \mathfrak b_{X\simeq\Delta,\varphi}
 =\mathfrak b_{X\boxtimes\Delta}/\mathcal J_\varphi.
\]
\end{definition}

\begin{proposition}[Graph quotient of a primitive comparison]\label{cyiv:graph-quotient}
Let $\varphi:\mathfrak b_\Delta\to\mathfrak b_X$ be a continuous Lie
morphism restricting to the identity on $\mathfrak h_\Pi$.
Use a completion in which the amalgam and closed quotient have
their stated universal properties.
Then
\[
 \mathfrak b_{X\boxtimes\Delta}/\mathcal J_\varphi
 \cong\mathfrak b_X.
\]
The two factor maps become $\id$ and $\varphi$.
If $\varphi$ is an isomorphism, the quotient is also isomorphic
to $\mathfrak b_\Delta$.
\end{proposition}
\begin{proof}
The maps $\id$ and $\varphi$ agree on $\mathfrak h_\Pi$.
They give $q_\varphi:\mathfrak b_{X\boxtimes\Delta}\to\mathfrak b_X$,
which kills the generators of $\mathcal J_\varphi$.
The inclusion of $\mathfrak b_X$ gives a map $s$ in the other direction.
The composite $q_\varphi s$ is the identity.
The other composite is the identity on the geometric factor.
On the Borcherds factor, the defining relation identifies $x$ with
$s\varphi(x)$, so it is also the identity there.
The two factor images topologically generate the quotient.
Continuity proves that the maps are inverse on the completion.
\end{proof}

The graph quotient implements a specified primitive comparison.
Constructing that comparison requires geometric simple states whose
extension brackets satisfy the Borcherds relations.
The quotient formula alone produces neither those states nor their
orientation data.

\chapter{Coproducts, pairings, and doubles}
A positive Hall algebra becomes a quantum group after additional structure is supplied. Let $H^+$ be a complete graded bialgebra, let $H^0$ be a Cartan algebra, and let $H^-$ be the opposite or dual half. Let
\[
 \angles{-,-}:H^+\widehattensor H^-\longrightarrow\kfield
\]
be a nondegenerate Hopf pairing.

\begin{definition}[Completed Hall double]
The completed Drinfeld double is the topological algebra
\[
 \widehat D(H^+)=H^-\widehattensor H^0\widehattensor H^+
\]
with internal products in the three factors and cross-relations determined by the Hopf pairing. Its representation category is monoidal by the coproduct and braided by the universal $R$-matrix whenever the canonical tensor converges.
\end{definition}

\begin{theorem}[Double construction]
Let $H^+$ and $H^-$ be paired complete graded bialgebras, let $H^0$ be a Cartan bialgebra acting by the stated weight grading, and let
\[
 \angles{-,-}:H^+\widehattensor H^-\longrightarrow\kfield
\]
be a continuous nondegenerate Hopf pairing. Assume the triangular multiplication map from $H^-\widehattensor H^0\widehattensor H^+$ is separated and complete. Then the pairing cross-relations determine a unique completed double with this triangular decomposition. On every continuous module category on which the canonical tensors converge, the completed canonical element gives the universal $R$-matrix.
\end{theorem}
\begin{proof}
The Hopf pairing defines the cross-relations between the paired halves and makes the three canonical embeddings bialgebra morphisms. The Hopf identities identify the two reductions of every triple product, proving associativity. Separated triangular decomposition gives uniqueness. Finite-degree canonical tensors form a compatible inverse system; convergence on the selected module category yields the quasitriangular element \cite{DrinfeldDouble,EtingofSchiffmann}.
\end{proof}

\begin{theorem}[Quantum Hall--Borcherds amalgam]
Let $H_X$ and $H_\Delta$ be complete Hopf algebras containing the same complete Cartan Hopf algebra $H^0$, with continuous Hopf retractions onto $H^0$.  Their completed amalgamated free product
\[
 H_{X\boxtimes\Delta}
 =H_X\widehat *_{H^0}H_\Delta
\]
has a unique complete Hopf structure extending the two coproducts, counits, and antipodes.  It represents pairs of continuous Hopf morphisms agreeing on $H^0$ and has canonical Hopf retractions to both factors.  A mixed universal $R$-matrix is equivalent to a continuous cross-pairing satisfying the two hexagon identities; the intrinsic $R$-matrices of the two retracts remain defined independently.
\end{theorem}
\begin{proof}
The algebraic amalgamated free product exists because complete filtered algebras are obtained degreewise from finite quotients.  The two coproducts agree on $H^0$, so the universal property extends them to a coproduct on the free product.  Coassociativity is checked on the two generating Hopf subalgebras.  The same argument extends the counit and antipode.  The Hopf retractions are induced by the factor retractions onto $H^0$.  A quasitriangular structure extending both factors requires the additional commutation data between the two generating halves; writing this data as a cross-pairing turns the quasitriangular and Yang--Baxter equations into the two hexagon identities.
\end{proof}

The deconcatenation coproduct on a bar complex and the representation-theoretic coproduct on a quantum group occupy different structures. A geometric comparison between them is a theorem about a selected model.

\chapter{Ordered chiral bialgebras}
A tensor product of representations requires compatible actions on
threefold products and the unit.
Retain the factorization coefficient category
$(\Fact_D(C),\tensor^!)$ on the smooth complex curve $C$.
Fix a category $\widehat{\mathcal V}_C$ of complete filtered objects
with specified presentations $F\simeq\varprojlim_r F_r$,
where $F_r\in\Fact_D(C)$.
Use derived limits in the fixed Ran enhancement and define
\[
 F\widehat{\tensor}^{!}G
 =\varprojlim_{r,s}(F_r\tensor^!G_s).
\]
Assume these limits remain factorizing and give
$\widehat{\mathcal V}_C$ a symmetric monoidal structure with unit $\mathbf1$.
Its associativity, symmetry, and unit equivalences must preserve
the presentations and the Ran structure.
Morphisms are continuous maps compatible with these presentations.
These are hypotheses on the chosen completion.

\begin{definition}[Ordered chiral bialgebra]\label{def:vol4-ordered-chiral-bialgebra}
An ordered chiral bialgebra is a counital $\Eone$-coalgebra object
in complete ordered chiral algebras:
\[
 A\in\operatorname{CoAlg}_{\Eone}
 \left(\Alg_{\Eone}
 (\widehat{\mathcal V}_C,\widehat{\tensor}^{!})\right).
\]
It has continuous multiplication $m$ and unit $\eta$, and continuous
morphisms of unital ordered chiral algebras
\[
 \Delta:A\longrightarrow A\widehat{\tensor}^{!}A,
 \qquad \epsilon:A\longrightarrow\mathbf1.
\]
The counit is the chosen augmentation.
Write $\tau$ for the coefficient symmetry, including Koszul signs.
Suppressing ambient associativity and unit equivalences, compatibility reads
\[
\begin{aligned}
 \Delta m&=(m\widehat{\tensor}^{!}m)
  (\id\widehat{\tensor}^{!}\tau\widehat{\tensor}^{!}\id)
  (\Delta\widehat{\tensor}^{!}\Delta),\\
 \Delta\eta&=\eta\widehat{\tensor}^{!}\eta,
 \qquad \epsilon m=\epsilon\widehat{\tensor}^{!}\epsilon,
 \qquad \epsilon\eta=\id_{\mathbf1}.
\end{aligned}
\]
Coassociativity and the counit identities read
\[
\begin{aligned}
 (\Delta\widehat{\tensor}^{!}\id)\Delta
   &=(\id\widehat{\tensor}^{!}\Delta)\Delta,\\
 (\epsilon\widehat{\tensor}^{!}\id)\Delta
   &=\id_A=(\id\widehat{\tensor}^{!}\epsilon)\Delta.
\end{aligned}
\]
For a homotopy-coherent presentation, these identities include the
higher coherences of the displayed coalgebra object.
Every operation and coherence preserves diagonal transitions and
disjoint-locus factorization maps.
An antipode is a continuous convolution inverse $S:A\to A$ of $\id_A$:
\[
 m(S\widehat{\tensor}^{!}\id)\Delta
 =\eta\epsilon
 =m(\id\widehat{\tensor}^{!}S)\Delta.
\]
A bialgebra with an antipode is a Hopf object in this coefficient category.
\end{definition}

A continuous left $A$-module is an object $M\in\widehat{\mathcal V}_C$
with a unital $\Eone$ action
$\rho_M:A\widehat{\tensor}^{!}M\to M$.
Its action and coherences are morphisms in these factorization coefficients.
Write $\operatorname{Mod}_A(\widehat{\mathcal V}_C)$ for their category.

\begin{theorem}[Representation fusion]\label{thm:vol4-representation-fusion}
An ordered chiral bialgebra makes
$\operatorname{Mod}_A(\widehat{\mathcal V}_C)$ monoidal.
Its product is $M\widehat{\tensor}^{!}N$ with action
\[
\begin{aligned}
 A\widehat{\tensor}^{!}(M\widehat{\tensor}^{!}N)
 &\xrightarrow{\Delta\widehat{\tensor}^{!}\id}
 (A\widehat{\tensor}^{!}A)\widehat{\tensor}^{!}
 (M\widehat{\tensor}^{!}N)\\
 &\xrightarrow{\id\widehat{\tensor}^{!}\tau_{A,M}
                       \widehat{\tensor}^{!}\id}
 (A\widehat{\tensor}^{!}M)\widehat{\tensor}^{!}
 (A\widehat{\tensor}^{!}N)\\
 &\xrightarrow{\rho_M\widehat{\tensor}^{!}\rho_N}
 M\widehat{\tensor}^{!}N.
\end{aligned}
\]
The unit is $\mathbf1$ with action induced by $\epsilon$.
The ambient associator and unit equivalences are module equivalences.

Suppose a continuous universal $R$-matrix is also given.
It is an invertible morphism
$R:\mathbf1\to A\widehat{\tensor}^{!}A$ for the tensor-product algebra,
satisfying
\[
\begin{gathered}
 R\Delta(a)=\Delta^{\op}(a)R,\qquad
 \Delta^{\op}=\tau_{A,A}\Delta,\\
 (\Delta\widehat{\tensor}^{!}\id)(R)=R_{13}R_{23},
 \qquad
 (\id\widehat{\tensor}^{!}\Delta)(R)=R_{13}R_{12},\\
 (\epsilon\widehat{\tensor}^{!}\id)(R)=\eta
 =(\id\widehat{\tensor}^{!}\epsilon)(R).
\end{gathered}
\]
The first equation is an identity of morphisms from $A$ to its completed tensor square.
Leg notation uses the unit and symmetry of $\widehat{\mathcal V}_C$.
In a homotopy-coherent presentation, require compatible lifts of these identities.
Then $c_{M,N}=\tau_{M,N}\circ R_{M,N}$ is a braiding,
where $R_{M,N}$ acts through the external tensor-product module action.
\end{theorem}
\begin{proof}
The symmetry combines the independent module actions into an
$A\widehat{\tensor}^{!}A$-action.
Restriction along the unital algebra map $\Delta$ gives the displayed action.
For three modules, the two actions use
$(\Delta\widehat{\tensor}^{!}\id)\Delta$ and
$(\id\widehat{\tensor}^{!}\Delta)\Delta$.
Coassociativity makes the ambient associator $A$-linear.
The counit identities make the ambient unit equivalences $A$-linear.
The coalgebra coherences give the pentagon and triangle coherences.
The completion hypotheses keep these maps continuous and factorizing.
Quasitriangularity makes $c_{M,N}$ an intertwiner and invertibility makes it an equivalence.
The two coproduct identities for $R$ give the hexagons.
Its counit identities give compatibility with the unit.
\end{proof}

The tensor in this construction is the completed transverse tensor
$\tensor^!$ on factorization coefficients.
A chiral-convolution or Hall coproduct requires its own construction
and a comparison in the specified category.
A spectral version additionally requires a difference coordinate,
a Laurent-series completion, and compatible domains for iterated expansions.
Its identities must hold in those domains before defining module actions.

\begin{example}[A counital algebra map without coassociativity]
In ordinary complex vector spaces, let $A=\C[x]$ and put
\[
 \Delta(x)=x\tensor1+1\tensor x+x\tensor x^2,
 \qquad \epsilon(x)=0.
\]
Polynomial substitution extends these assignments to unital algebra maps.
Writing $F(u,v)=u+v+uv^2$, the equations
$F(0,v)=v$ and $F(u,0)=u$ prove both counit identities.
But
\[
\begin{split}
 F(F(u,v),w)-F(u,F(v,w))
 &=-uv^2w^4-uv^2w^2-2uvw^3-2uvw\\
 &=-uvw(1+w^2)(vw+2)\ne0.
\end{split}
\]
On three regular modules, the two actions differ already on
$1\tensor1\tensor1$.
Thus the canonical tensor associator is not $A$-linear.
This ordinary example shows why multiplicativity and the counit identities
must be supplemented by coassociativity.
It makes no assertion about a geometric realization in factorization coefficients.
\end{example}

\part{Exact model cases}

\chapter{The \texorpdfstring{$A_2$}{A2} Hall algebra}
For the quiver $1\to2$ over $\mathbb F_q$, the twisted Hall algebra has generators $E_1,E_2$ and $v^2=q$. Counting subrepresentations gives
\[
 E_1^2E_2-(v+v^{-1})E_1E_2E_1+E_2E_1^2=0.
\]
For a representation of dimension vector $(2,1)$, the arrow has rank zero or one. The relevant flag counts are
\[
\begin{array}{c|ccc}
\rank&112&121&211\\ \hline
0&q+1&q+1&q+1\\
1&q+1&1&0.
\end{array}
\]
After the Hall twists $v^{-1},1,v$, both rows give zero. This is the elementary geometric origin of the quantum Serre relation \cite{Ringel,Green}.

\chapter{The tripled Jordan quiver and \texorpdfstring{$\mathbb C^3$}{C3}}
Let $Q$ have one vertex and loops $x,y,z$, with potential
\[
 W=\Tr x[y,z].
\]
The Jacobi algebra is $\C[x,y,z]$. The critical CoHA uses the vanishing cycles of $\Tr W$ on the representation spaces. Its positive currents are generated by the one-vertex dimension sectors.

\begin{theorem}[Affine-Yangian positive half]
After the standard equivariant localization and completion, the critical CoHA of the tripled Jordan quiver is identified with the positive half of the affine Yangian of $\mathfrak{gl}_1$. The independently defined affine Yangian carries its standard Cartan and negative current halves; the positive-half comparison embeds the CoHA into that triangular presentation.
\end{theorem}
\begin{proof}
The CoHA product becomes the shuffle product under localization. Dimensional reduction identifies the tripled-quiver critical theory with the equivariant surface CoHA, and Davison's affine-BPS theorem identifies its completed positive algebra with the strictly positive affine Yangian. The shuffle kernel satisfies the affine-Yangian current relations with equivariant parameters $\epsilon_1,\epsilon_2,\epsilon_3$ and Calabi--Yau relation $\epsilon_1+\epsilon_2+\epsilon_3=0$. PBW and current comparison identify the positive algebra. The independently presented full affine Yangian is the triangular algebra generated by its negative currents, Cartan currents, and this positive half \cite{DavisonAffineBPS,DavisonMeinhardt,SVYangians,Tsymbaliuk}.
\end{proof}

The vacuum representation produces a $W_{1+\infty}$-type vertex algebra through an evaluation homomorphism. The CoHA, its double, and the vacuum algebra occupy successive levels: positive algebra, paired quantum group, and distinguished representation.

\chapter{Surface CoHAs and ADE Yangians}
Let $S_\Gamma\to\C^2/\Gamma$ be the minimal resolution for a finite subgroup $\Gamma\subset\operatorname{SL}_2(\C)$. The exceptional curves have ADE intersection matrix. Moduli of sheaves supported on these curves carry convolution correspondences.

\begin{theorem}[Curve-supported surface CoHA]
The completed positive CoHA of one-dimensional sheaves supported on the exceptional curves is identified with a nonstandard positive half of the affine Yangian of the corresponding ADE Lie algebra.
\end{theorem}
\begin{proof}
The theorem is the curve-supported surface construction of Diaconescu--Porta--Sala--Schiffmann--Vasserot \cite{DPPSSV}. The simple curve sheaves give the simple generators, Hecke correspondences give the current operators, and the completed PBW comparison gives the affine Yangian presentation.
\end{proof}

This theorem is a local K3 model whenever an ADE configuration occurs on a K3 surface. The global K3 surface adds compact curve classes and monodromy beyond the local root lattice.

\chapter{Coulomb branches and shifted Yangians}
For a three-dimensional $\mathcal N=4$ gauge theory, the BFN Coulomb branch is defined by equivariant Borel--Moore homology of an affine-Grassmannian correspondence. Its convolution algebra quantizes the Coulomb branch \cite{BFN}. In quiver gauge theories, shifted Yangians and their truncations provide presentations of these quantizations. Boundary conditions select the shift and truncation data.

The BFN construction supplies a second geometric source of Yangians. Its grading and coproduct arise from loop rotation and convolution, while the surface CoHA source arises from sheaf extensions. A comparison between the two is a model-specific equivalence.

\part{K3 symmetries}

\chapter{The Mukai lattice}
For a K3 surface $S$, the Mukai lattice is
\[
 \Mukai(S,\Z)=H^0(S,\Z)\oplus H^2(S,\Z)\oplus H^4(S,\Z)
\]
with pairing
\[
 ((r,c,s),(r',c',s'))=c\cdot c'-rs'-r's.
\]
It is even unimodular of signature $(4,20)$ and rank $24$.

\chapter{The Nakajima Heisenberg algebra}
Let
\[
 \mathbb H_S=\bigoplus_{n\ge0}H^*(S^{[n]}).
\]
For $m\in\Z\setminus\{0\}$ and $\alpha\in H^*(S)$, the Nakajima correspondence gives an operator $p_m(\alpha)$.

\begin{theorem}[K3 Heisenberg action]
The operators satisfy
\[
 [p_m(\alpha),p_n(\beta)]
 =m\,\delta_{m+n,0}\int_S\alpha\cup\beta\;\id.
\]
They generate the rank-twenty-four Heisenberg algebra acting irreducibly on $\mathbb H_S$ with vacuum $1\in H^*(S^{[0]})$.
\end{theorem}
\begin{proof}
The commutator is the excess intersection of the two incidence correspondences. Opposite creation and annihilation correspondences meet along the diagonal with multiplicity $m$ and intersection pairing $\int_S\alpha\beta$. Correspondences of the same sign commute. The vacuum construction gives the Fock representation \cite{Nakajima97}.
\end{proof}

This action is global, compact, and abelian at the Lie level. It is an exact K3 symmetry and supplies the rank-twenty-four Fock sector of any broader K3 algebra.

\chapter{The orthogonal LLV algebra}
For a compact hyperkähler manifold, every Lefschetz class gives an $\mathfrak{sl}_2$ triple on cohomology. The Lie algebra generated by all such triples is the Looijenga--Lunts--Verbitsky algebra.

\begin{theorem}[Orthogonal LLV symmetries]
For a K3 surface $S$, the LLV algebra acting on $H^*(S)$ is
\[
 \mathfrak g_{\mathrm{LLV}}(S)
 \cong\mathfrak{so}\bigl(\widetilde H(S),(-,-)_{\mathrm{Muk}}\bigr)
 \cong\mathfrak{so}(4,20).
\]
For an irreducible hyperk\"ahler manifold $M$ of $K3^{[n]}$ type with $n\ge2$, one has $b_2(M)=23$ and
\[
 \mathfrak g_{\mathrm{LLV}}(M)\cong\mathfrak{so}(4,21).
\]
Moduli spaces of stable objects on $S$ carry their own LLV actions; Mukai and monodromy correspondences relate these actions to the K3 lattice.
\end{theorem}
\begin{proof}
The Looijenga--Lunts--Verbitsky theorem identifies the Lie algebra generated by all Lefschetz triples with the orthogonal algebra of the Mukai extension $H^2(M)\oplus U$. For $S$, the extension has signature $(4,20)$. For $K3^{[n]}$ type with $n\ge2$, the Beauville--Bogomolov lattice has signature $(3,20)$, so adjoining the hyperbolic plane gives signature $(4,21)$ \cite{LooijengaLunts,Verbitsky96}.
\end{proof}

The LLV algebra is a classical cohomological symmetry. A Yangian deformation of the full orthogonal action consists of two further structures: a Drinfeld current presentation and a compatible spectral coproduct.

\chapter{Local ADE Yangians}
An ADE configuration of $(-2)$-curves on $S$ gives a root sublattice of $H^2(S,\Z)$. Sheaves supported on the configuration produce the surface CoHA of the preceding part and hence a completed positive affine ADE Yangian. Nakajima quiver varieties supply compatible Yangian actions on their equivariant cohomology.

Thus K3 geometry contains local nonabelian Yangian sectors indexed by ADE enhancements of the Picard lattice. Picard monodromy transports the ADE root sublattices over the corresponding Noether--Lefschetz loci. A family of sheaf moduli and Hecke correspondences compatible with that monodromy transports the Yangian actions themselves.

\chapter{The K3 symmetry trichotomy}
The three exact structures are
\begin{center}
\small
\begin{tabularx}{0.96\textwidth}{@{}>{\raggedright\arraybackslash}p{0.23\textwidth}>{\raggedright\arraybackslash}X>{\raggedright\arraybackslash}p{0.28\textwidth}@{}}
\toprule
\textbf{Structure}&\textbf{Carrier}&\textbf{Operation}\\
\midrule
Mukai Heisenberg&$\bigoplus_nH^*(S^{[n]})$&creation and annihilation\\
LLV $\mathfrak{so}(4,20)$&$H^*(S)$; separately $\mathfrak{so}(4,21)$ on $K3^{[n]}$ type&Lefschetz Lie action\\
Local ADE Yangian&curve-supported equivariant moduli&Hecke and current correspondences\\
\bottomrule
\end{tabularx}
\end{center}
A global K3 quantum algebra is a system that contains these three actions with explicit comparison homomorphisms. The two-Calabi--Yau BPS theorem supplies a global generalized Kac--Moody positive algebra. A full Yangian or double is specified by a compatible coproduct, spectral parameter, negative half, Cartan part, pairing, and completion.

\chapter{K3 correspondences on fixed coefficients}
Fix rational cohomology and a collection $\mathscr M_S$ of smooth
projective moduli spaces on a projective K3 surface $S$.
Include the Hilbert schemes $S^{[n]}$ and close the collection under
finite products.
For $M,N\in\mathscr M_S$, define the graded kernel space
\[
 K(M,N)=H^*(M\times N,\mathbb Q)[2\dim_{\mathbb C}M].
\]
A kernel of degree $r$ acts from $H^j(M)$ to $H^{j+r}(N)$ by
\[
 T_\alpha(v)=(p_N)_*(p_M^*v\cup\alpha).
\]
The shift in $K(M,N)$ accounts for integration over $M$.
All pushforwards in this construction are proper.

\begin{proposition}[Kernel composition]\label{cyiv:k3-kernels}
The product
\[
 \beta\circ\alpha=(p_{13})_*
 (p_{12}^*\alpha\cup p_{23}^*\beta)
\]
defines an associative graded category with diagonal identities.
The kernel action identifies $K(M,N)$ with
$\operatorname{Hom}^*(H^*(M),H^*(N))$.
\end{proposition}
\begin{proof}
The K\"unneth isomorphism and Poincar\'e duality identify
$H^*(M\times N)[2\dim M]$ with $H^*(M)^\vee\otimes H^*(N)$.
Under this identification, $T_\alpha$ is the associated linear map.
Integration over the middle factor contracts its dual pair, so
$T_{\beta\circ\alpha}=T_\beta T_\alpha$.
Faithfulness then proves associativity and the diagonal identity.
External products give tensor products with their Koszul signs.
\end{proof}

A geometric symmetry is specified by selected kernels in this category.
Nakajima incidence kernels connect different Hilbert schemes.
Lefschetz multiplication and its adjoint act on each individual
cohomology space.
Using every cohomology kernel would instead give all linear maps;
the selected geometric kernels carry the representation-theoretic content.

\chapter{The algebra generated by geometric actions}
Put $V_S=\bigoplus_{M\in\mathscr M_S}H^*(M,\mathbb Q)$.
Choose a set of homogeneous operators $T_\lambda$ on $V_S$, each sending
a vector with finite support to a vector with finite support.
The Nakajima operators on the Hilbert-scheme summand satisfy this
condition: each changes the point number by a fixed integer.
Individual Lefschetz operators, extended by zero on other summands,
also satisfy it.
Let $F$ be the free unital graded algebra on symbols $t_\lambda$.

\begin{proposition}[Geometric action algebra]\label{cyiv:k3-action}
The map
\[
 \rho:F\longrightarrow\operatorname{End}_{\mathbb Q}(V_S),
 \qquad t_\lambda\longmapsto T_\lambda
\]
defines a faithful action of $A_S=F/\ker\rho$.
For every algebra $B$, maps $A_S\to B$ are exactly maps $F\to B$
that kill $\ker\rho$.
Products and commutators of the selected kernels are their convolution
products and commutators on the relevant moduli spaces.
\end{proposition}
\begin{proof}
Finite words in the selected operators preserve finite support, so
$\rho$ is defined by composition.
The quotient acts faithfully by definition of its kernel.
The factorization property is the universal property of an algebra
quotient.
The convolution assertion follows from Proposition~\ref{cyiv:k3-kernels}.
\end{proof}

An equivariant ADE construction uses a separate coefficient field,
for example $\operatorname{Frac}H_T^*(\mathrm{pt},\mathbb Q)$,
and its prescribed localization and proper pushforwards.
Joining it to $A_S$ requires maps between the actual state spaces that
intertwine those operators and their coefficients.
A completion of $A_S$ requires a specified topology and continuity of
every generator.
The indefinite Mukai norm alone supplies neither condition.
The desired common K3 quantum algebra must construct these local-to-global
actions and their mixed convolution relations.

\chapter{Tensor reconstruction and geometric kernels}
Let $\mathscr C$ be an essentially small rigid linear monoidal category
and let $\omega:\mathscr C\to\operatorname{Vect}_{\mathrm{fd}}(\mathbb Q)$
be a strong monoidal functor.
A choice of geometric kernels specifies the morphisms of $\mathscr C$;
it is part of the input.
Define
\[
 H=\left(\bigoplus_{M\in\mathscr C}\omega(M)^*\otimes\omega(M)\right)
 \big/\left\langle
 [\lambda\circ\omega(f)\otimes v]_M
 -[\lambda\otimes\omega(f)v]_N\right\rangle,
\]
where $f:M\to N$, $v\in\omega(M)$, and $\lambda\in\omega(N)^*$.

\begin{proposition}[Matrix-coefficient coalgebra]\label{cyiv:k3-coend}
For a basis $(e_i)$ with dual basis $(e^i)$, the formulas
\[
 \Delta[\lambda\otimes v]_M
 =\sum_i[\lambda\otimes e_i]_M\otimes[e^i\otimes v]_M,
 \qquad \epsilon[\lambda\otimes v]_M=\lambda(v)
\]
are basis-independent and make $H$ a coalgebra.
Tensor product gives its multiplication, and duals give its antipode.
Each $\omega(M)$ is a comodule by
\[
 v\longmapsto\sum_i e_i\otimes[e^i\otimes v]_M.
\]
Every selected kernel is a comodule morphism.
\end{proposition}
\begin{proof}
The identity tensor $\sum_i e_i\otimes e^i$ is basis-independent.
Inserting it twice proves coassociativity; evaluation proves both
counit identities.
The defining relations express naturality under $f$ and are preserved
by these formulas.
The strong monoidal constraints identify the coefficient of $M\otimes N$
with the product of coefficients of $M$ and $N$.
Their pentagon and unit identities give the bialgebra axioms.
The evaluation and coevaluation identities for duals give the antipode
identities.
The same relations prove that $\omega(f)$ intertwines the coactions.
\end{proof}

The linear dual of $H$ represents linear natural endomorphisms of
$\omega$.
Tensor-compatible natural automorphisms are the corresponding
multiplicative, unit-preserving families, not all elements of this dual.
An equivalence with a comodule category requires the additional exactness
and reconstruction hypotheses stated in the Tannaka theorem below.
The action algebra $A_S$ acts by selected kernels; the coend records
symmetries commuting with those kernels.

\part{The exact \texorpdfstring{$K3\times E$}{K3 x E} protected completion}

\chapter{The protected input}
The K3 elliptic genus is $2\phi_{0,1}$.
The Borcherds half determines the weight-five denominator $D_5$ and its
Borcherds Lie superalgebra $\mathfrak g_\Delta$.
Fix its positive root spaces with their parity and proper height.
The positive algebra and its root Fock space are defined from this Lie
superalgebra, independently of a compact Hall realization.

For each finite zero-potential quiver window $Q_{\Delta,H}$, form the
Ginzburg algebra $\Gamma_{\Delta,H}$.
Restriction to full subquivers gives the projections of
Proposition~\ref{cyiv:ginzburg-projection}.
Write $\Gamma_\Delta$ for this specified inverse system.
A nondegenerate continuous Calabi--Yau structure on a completed limit
is additional data.
An ordinary Hall category $\mathcal B_\Delta$ must carry its chosen
quiver coefficients and extension correspondences.
Its identification with a root presentation requires the corresponding
Hall theorem on that precise datum.

\begin{definition}[Root-datum realizations]
A protected root-datum realization records
\[
 \mathfrak P(K3\times E)=
 (\Gamma_\Delta,\mathcal B_\Delta,\mathcal H_{\mathrm{1p}},
 \mathcal F^{\mathrm{hyb}}_{\Delta,E},\widehat U_q(\mathfrak g_\Delta),
 \mathbb A_Z).
\]
Here $\mathcal H_{\mathrm{1p}}=\bigoplus_{\alpha>0}(\mathfrak g_\Delta)_\alpha$
with its parity.
The current-factorization, paired quantum-group, and Fredholm entries
retain their own coefficient categories, completions, and defining
hypotheses.
Maps among these entries are specified realization maps, not consequences
of using the same root labels.
\end{definition}

The enveloping and Chevalley constructions are functorial under the
primitive maps of Theorem~\ref{cyiv:primitive-comparison}.
The root Fock trace is computed from the actual graded root spaces.
A Fredholm realization adds a family of operators on a specified domain
and the convergence and orientation data identifying its Pfaffian with
the denominator.
These are the analytic inputs to the square-root comparison below.

\begin{definition}[Integral protected realization]
Let $\mathfrak p_X^+$ be the BPS primitive Lie algebra of a selected strong-oriented cohomological lift of the geometric $E$-equivariant Hall source.  The integral protected realization is
\[
 \mathfrak P_{\mathrm{int}}(X,\Delta)
 =\left(
   \Hcal_X^{E,\mathrm{mot}},
   \mathfrak P(K3\times E),
   \mathfrak b_{X\boxtimes\Delta},
   r_X,r_\Delta,
   I_{X,\Pi}^E
 \right).
\]
At the motivic level the first, second, and character components are defined.  A strong orientation enhancement supplies $\mathfrak p_X^+$ and hence the amalgamated primitive component.
\end{definition}

\begin{proposition}[Primitive graph comparison]
Suppose the two primitive Borels, their common Cartan action, and their
completed amalgam are supplied as above.
Their Cartan retractions induce retractions from the amalgam to each
factor.
For each specified primitive map $\varphi$ preserving Cartan action,
its graph quotient is canonically the geometric Borel.
\end{proposition}
\begin{proof}
The identity on one Borel and the Cartan projection of the other agree
on the common Cartan, so the amalgam gives the retraction.
The quotient statement is Proposition~\ref{cyiv:graph-quotient}.
\end{proof}

The motivic integration map and primitive comparison are different maps.
Strong orientation gives the coefficient coherence of a cohomological
Hall construction under its stated geometric hypotheses.
Extracting a primitive Lie algebra additionally requires the appropriate
BPS or coproduct construction.

\chapter{Cyclic open-string operations}
Let $A$ be a finite-dimensional, uncurved cyclic $A_\infty$ algebra
over $\mathbb C$, with a nondegenerate pairing of degree $-3$.
Use the cohomological suspension $s:A\to A[1]$ of degree $-1$.
The suspended operations $b_n:(A[1])^{\otimes n}\to A[1]$ have degree one.
Our sign convention is the coderivation convention $b^2=0$ on the
tensor coalgebra, with the Koszul rule for moving homogeneous inputs.
Write $\omega$ for the resulting degree-$-1$ symplectic form on $A[1]$.
Fix the Hamiltonian convention $\iota_{X_F}\omega=dF$ and
$[X_F,X_G]=X_{\{F,G\}}$.

For the universal degree-zero coordinate $x$ on $A[1]$, put
\[
 S(x)=\sum_{n\ge1}\frac1{n+1}\omega(x,b_n(x^{\otimes n})).
\]
Cyclic invariance is read with the suspended Koszul signs.
All functions are formal power series at the origin.
An equivalent unsuspended convention writes this as
$\frac12\langle a,m_1a\rangle+
\sum_{n\ge2}\frac1{n+1}\langle a,m_n(a^{\otimes n})\rangle$.

\begin{proposition}[Cyclic master equation]\label{cyiv:cyclic-master}
The action satisfies $\{S,S\}=0$.
\end{proposition}
\begin{proof}
Cyclicity makes the $n+1$ derivatives of the arity-$n+1$ term equal,
with the suspended Koszul signs.
Thus $X_S(x)=\sum_{n\ge1}b_n(x^{\otimes n})$.
The coderivation identity implies $X_S^2=0$.
The Hamiltonian identity gives $X_{\{S,S\}}=2X_S^2=0$.
Nondegeneracy of $\omega$ makes $\{S,S\}$ constant.
Uncurvedness removes $b_0$, so $S$ starts in order two and its bracket
has no constant term.
The constant is therefore zero.
\end{proof}

A finite-height Ginzburg realization requires such a cyclic model on
the selected proper objects before this scalar action is applied.
A Hall realization additionally requires a moduli problem with
critical coefficient complexes.
For an extension $A_0\to F\to B_0$, put
$L(A_0,B_0)=\det R\operatorname{Hom}(A_0,B_0)$.
Three-Calabi--Yau duality gives
\[
 R\operatorname{Hom}(B_0,A_0)
 \simeq R\operatorname{Hom}(A_0,B_0)^\vee[-3].
\]
Applying the determinant functor to the two derived Hom triangles gives
\[
 \det R\operatorname{Hom}(F,F)
 \cong \det R\operatorname{Hom}(A_0,A_0)
 \otimes\det R\operatorname{Hom}(B_0,B_0)
 \otimes L(A_0,B_0)^{\otimes2}.
\]
Indeed, the four subquotients are the two diagonal Hom complexes and
the two cross complexes.
Dualization and the odd shift cancel on their determinant lines,
so the two cross determinants coincide.
Orientation lines require square-root isomorphisms along this formula,
compatible on triple flags and with any quotient descent.
The cyclic master equation does not select these isomorphisms.

\chapter{An algebraic open--closed comparison}
A finite-dimensional ordinary Lie algebra $\mathfrak l$ gives a
precise coefficient model for an enveloping-algebra centre.
Let $U=U(\mathfrak l)$ over $\mathbb C$.
The adjoint module $U^{\mathrm{ad}}$ has action
$x\cdot u=xu-ux$.
Its Chevalley cochains are
\[
 C^n(\mathfrak l,U^{\mathrm{ad}})
 =\operatorname{Hom}_{\mathbb C}(\Lambda^n\mathfrak l,U),
\]
with differential
\[
\begin{split}
 (df)(x_1,\ldots,x_{n+1})
 &=\sum_i(-1)^{i+1}[x_i,f(x_1,\widehat x_i,\ldots,x_{n+1})]\\
 &\quad+\sum_{i<j}(-1)^{i+j}
 f([x_i,x_j],x_1,\widehat x_i,\widehat x_j,\ldots,x_{n+1}).
\end{split}
\]
The second line places the bracket first and retains the order of the
remaining arguments.

\begin{theorem}[Hochschild--Chevalley comparison]\label{cyiv:hh-ce}
Restriction followed by antisymmetrization defines a quasi-isomorphism
\[
 C^\bullet_{\mathrm{HH}}(U,U)
 \longrightarrow C^\bullet(\mathfrak l,U^{\mathrm{ad}}),
 \qquad
 \phi\longmapsto\left[(x_1,\ldots,x_n)\mapsto
 \sum_{\sigma\in S_n}\operatorname{sgn}(\sigma)
 \phi(x_{\sigma(1)},\ldots,x_{\sigma(n)})\right].
\]
This is a comparison of ordinary cochain complexes.
\end{theorem}
\begin{proof}
The Chevalley bimodule resolution has terms
$K_n=U\otimes\Lambda^n\mathfrak l\otimes U$.
Its differential has the left-minus-right multiplication terms and
the bracket terms with the signs dual to the displayed differential.
Filter $K_n$ by total PBW degree, counting each exterior generator
as degree one.
The multiplication terms preserve this degree and bracket terms lower it.
The associated graded is the Koszul resolution of the diagonal in
$\operatorname{Sym}(\mathfrak l)\otimes\operatorname{Sym}(\mathfrak l)$,
with regular sequence $x_i\otimes1-1\otimes x_i$.
That Koszul complex is exact away from degree zero.
The filtration is bounded below and exhaustive, and each element has
finite filtration degree, so lifting a leading term proves exactness
of $K_\bullet$ by induction.
Its augmentation has image $U$.
Thus $K_\bullet$ is a free bimodule resolution.

Antisymmetrization gives a map from $K_\bullet$ to the ordinary bar
resolution, lifting the identity of $U$.
Adjacent products in its bar differential pair to commutators;
$xy-yx=[x,y]$ gives exactly the bracket terms of $K_\bullet$.
The outer products give its left-minus-right terms.
A comparison between projective resolutions lifting the identity is
a homotopy equivalence.
Applying $\operatorname{Hom}_{U\otimes U^{\mathrm{op}}}(-,U)$ gives
the stated formula and quasi-isomorphism.
\end{proof}

This is the classical Cartan--Eilenberg comparison.
For a finite Lie superalgebra, the same construction uses the
super-exterior powers and permutation Koszul signs.
Odd generators produce symmetric exterior directions, so the resolution
can be unbounded; its algebraic direct sums and cochain products remain
part of the convention.
The PBW and super-Koszul proof applies with those conventions.

A physical mixed holomorphic--topological realization requires a map
from its independently defined bulk observables to this cochain object.
The theorem specifies the adjoint coefficients and the actual algebraic
comparison.
An $E_2$ refinement, a current-factorization realization, and compatibility
with root-height transitions require maps preserving their additional
operations.
A quotient of Lie algebras alone does not provide a covariant map of
Hochschild centres.

\chapter{The exact scalar and first-order observables}
Set
\[
 \DenIg(Z)=D_5(2Z),
 \qquad \jmath(Z)=Z/2,
 \qquad \jmath^*\DenIg=D_5.
\]
The one-particle root spaces give
\[
 \sch\Fock(\gDelta_+)
 =\prod_{\alpha>0}(1-e^{-\alpha})^{-\sdim\gDelta_\alpha}.
\]
The oriented root operator gives
\[
 e^{-\rho}\operatorname{sPf}_F(\mathbb A_Z)=D_5(2Z).
\]
The primitive reduced curve-counting theorem gives
\[
 \mathcal Z^{\red}_{K3\times E}=-(\jmath^*\DenIg)^{-2}=-\chi_{10}^{-1}.
\]

\begin{theorem}[Protected square-root theorem]
Assume the stated root-space Fredholm comparison with its orientation,
Weyl shift, and convergence domain.
Within that realization, the denominator section $\DenIg$ is simultaneously:
\begin{enumerate}
\item the Weyl denominator of $\gDelta$;
\item the Euler character of the bar--Chevalley complex;
\item the oriented Fredholm Pfaffian of $\mathbb A_Z$;
\item the PBW determinant of the parity-refined root superspace $\mathcal H_{\mathrm{1p}}$;
\item the automorphic denominator whose half-coordinate pullback squares to the reduced scalar denominator.
\end{enumerate}
All five identifications preserve root degree and Weyl transport.
\end{theorem}
\begin{proof}
The first two are the denominator and bar--Chevalley theorems.  The third is the Fredholm theorem.  PBW for the Borcherds Lie superalgebra identifies the root superspace determinant, giving the fourth.  The half-coordinate identity $\jmath^*\DenIg=D_5$, the theta normalization, and the reduced curve-counting theorem give the fifth.
\end{proof}

\chapter{Physical meaning}
The root-datum realizations retain several distinct constructions.
The finite Ginzburg algebras and their projections form an inverse system.
A completed three-Calabi--Yau sector requires a continuous nondegenerate
Calabi--Yau structure and compatible cyclic models on the selected proper objects.
The quiver Hall category retains its chosen coefficients and extension
correspondences. A Hall theorem on that datum must identify its root presentation.
The Borcherds root superspace defines the root-graded Fock space and its
supercharacter. Its identification with physical one-particle states
requires a parity-preserving graded comparison intertwining the primitive actions.

A realization on the elliptic curve must construct the current object
with its local and wrapped supports, collision operations, and comparison maps.
The quantum Borcherds algebra, with its specified pairing and completion,
acts as a line-operator symmetry only through a representation preserving
fusion and exchange. The root operator gives the Fredholm product on its
specified convergence domain. A physical determinant interpretation requires
an operator comparison and an orientation-preserving identification of the
determinant lines and sections. The Fock character gives the scalar square.
An operator trace additionally requires its state space and trace domain.

Simple-root states satisfying the defining relations give a primitive Lie map.
The explicit enveloping and Chevalley extensions preserve products and brackets.
A compact string realization must additionally construct its brane functor,
Hall orientation maps, paired quantum coproduct, and trace-line comparison.

\chapter{Primitive comparisons and their algebraic extensions}
Let $\mathfrak n_\Delta$ be the positive Borcherds Lie superalgebra.
Choose homogeneous parity-preserving images of its simple generators
in a charge-graded Lie superalgebra $\mathfrak p$.
Assume these images satisfy every defining relation.

\begin{theorem}[Extension of a primitive comparison]\label{cyiv:primitive-comparison}
The chosen images determine a unique Lie-superalgebra map
$f:\mathfrak n_\Delta\to\mathfrak p$.
It induces
\[
 U(f)(x_1\cdots x_r)=f(x_1)\cdots f(x_r),
 \qquad C_*(f)=\operatorname{Sym}^c(f[1]).
\]
These are respectively an enveloping-algebra map and a Chevalley
chain-coalgebra map.
Filtration-preserving maps extend to the specified inverse-limit
completions.
If $f$ is surjective and both algebras have finite charge components
and proper positive height, equality of their parity-refined PBW
characters implies that $f$ and $U(f)$ are isomorphisms.
\end{theorem}
\begin{proof}
The defining presentation gives $f$.
Its tensor-algebra extension carries
$xy-(-1)^{|x||y|}yx-[x,y]$ to the corresponding relation.
It therefore descends to $U(f)$.
The shifted symmetric-coalgebra map commutes with the Chevalley
differential because $f$ preserves brackets.
Compatible maps on filtration quotients give maps on their inverse limits.
Surjectivity of $f$ gives surjectivity of $U(f)$.
Character equality makes each finite charge component an isomorphism,
separately in both parities.
The PBW embeddings then imply injectivity of $f$.
\end{proof}

At level zero the current Lie conformal algebra is
$\mathbb C[\partial]\otimes\mathfrak n_\Delta$ with
$[x_\lambda y]=[x,y]$ on constant generators.
The map $\partial^j x\mapsto\partial^j f(x)$ preserves this bracket
and sesquilinearity.
A nonzero central level additionally requires preservation of the
specified invariant pairing.
Hall convolution, Calabi--Yau structures, quantum coproducts, and
determinant lines require their own comparison data.

\chapter{Finite quivers and Calabi--Yau completions}
For a finite quiver $Q$ with zero potential, let $\Gamma(Q,0)$ be its
path-length-completed Ginzburg algebra.
Path products are read from left to right: if $a:i\to j$ and
$b:j\to k$, then $ab:i\to k$.
For each arrow $a:i\to j$, add $a^*:j\to i$ of degree $-1$.
At each vertex add a loop $t_i$ of degree $-2$.
Original arrows have degree zero, and
\[
 da=da^*=0,\qquad
 dt_i=e_i\left(\sum_a[a,a^*]\right)e_i.
\]

\begin{proposition}[Restriction to a full subquiver]\label{cyiv:ginzburg-projection}
If $Q_0$ is a full subquiver, retaining its generators and killing all
other vertices, arrows, reverse arrows, and loops defines a continuous
dg-algebra projection
\[
 \Gamma(Q,0)\longrightarrow\Gamma(Q_0,0).
\]
The inclusion of retained graded generators need not commute with
the differential.
\end{proposition}
\begin{proof}
The projection respects source and target idempotents and extends
multiplicatively.
At a retained vertex, every term involving a discarded arrow vanishes.
The other terms are precisely $dt_i$ in the smaller algebra.
All arrow differentials vanish, so the projection is a chain map.
It preserves the path-length filtration and is continuous.
For the final assertion, start with one vertex $i$ and no arrows,
where $dt_i=0$.
Add a vertex $j$ and an arrow $a:i\to j$.
In the larger algebra $dt_i=aa^*$, a nonzero length-two path.
Thus the graded inclusion is not a dg map.
\end{proof}

Ginzburg completion supplies a three-Calabi--Yau algebra for each finite
quiver in its stated setting.
The displayed projections specify an inverse system for compatible
zero-potential windows.
A Calabi--Yau structure on a completed limit requires its own
nondegeneracy and continuity argument.
Comparisons with Hall categories and quantized Borcherds algebras must
also preserve orientations, root relations, and the paired Hopf data.

\part{Quantum groups, exchange, and representation theory}

\chapter{From positive half to full quantum group}
Suppose the compact or local positive algebra $H^+$ is paired continuously with a negative half $H^-$, carries a completed coproduct, and admits Cartan data and a separated triangular completion. The completed double
\[
 \qgroup=\widehat D(H^+)
\]
has triangular decomposition $H^-\widehattensor H^0\widehattensor H^+$. Its universal $R$-matrix acts on continuous tensor products of modules. The classical limit is a Lie bialgebra whose cobracket is the first-order term of the coproduct asymmetry.

For $\C^3$, $H^+$ is the positive affine Yangian. For local ADE surfaces, it is the completed positive affine ADE Yangian. For the Igusa algebra, the paired Borel construction of the preceding part supplies a height-complete formal quantum group; a compact Hall realization maps to it through primitive, coproduct, and pairing comparisons.

\chapter{Tannaka reconstruction}
Let $\mathcal L$ be an essentially small rigid $\kfield$-linear abelian monoidal category with finite-dimensional morphism spaces, finite-length objects, and exact tensor product. Let
\[
 \omega:\mathcal L\to\Vect_{\mathrm{fd}}
\]
be a faithful exact strong monoidal functor. Assume the coend
\[
 H=\int^{M\in\mathcal L}\omega(M)^*\tensor\omega(M)
\]
exists and that the canonical comparison functor to finite-dimensional $H$-comodules is comonadic. Tensor product makes $H$ a bialgebra, and rigidity supplies its antipode. A braiding gives a coquasitriangular form; a meromorphic braiding gives its spectral version.

\begin{theorem}[Rigid Tannaka reconstruction]
Under the preceding finite and comonadic hypotheses, the comparison functor is a monoidal equivalence
\[
 \mathcal L\simeq\operatorname{Comod}_{\mathrm{fd}}(H).
\]
Braiding transports to the universal coquasitriangular form on $H$.
\end{theorem}
\begin{proof}
Matrix coefficients define the universal $H$-coaction on every $\omega(M)$. Tensor product and duality give the bialgebra and antipode structures. The Barr--Beck comonadic theorem identifies $\mathcal L$ with coalgebras for the induced comonad, and the coend identifies that comonad with $H\tensor-$. The hexagon identities transport to the coquasitriangular identities.
\end{proof}

This reconstruction explains why a line-operator category, a fiber functor, and exchange together produce a quantum group. The ordered boundary algebra provides ordered composition.  A compatible bialgebra coproduct supplies tensor products of its modules, and exchange supplies additional braiding data.

\chapter{KZ and rational exchange}
For a simple Lie algebra $\mathfrak g$, the Knizhnik--Zamolodchikov connection is
\[
 \nabla=d-\frac{\hbar}{2\pi i}
 \sum_{i<j}\Omega_{ij}\,d\log(z_i-z_j).
\]
Flatness follows from the infinitesimal braid relations. Drinfeld--Kohno identifies its monodromy tensor category with the corresponding Drinfeld--Jimbo category after a choice of associator.

Rational difference exchange instead gives the Yangian $R$-matrix
\[
 R(u)=1-\frac{\hbar P}{u}
\]
in the fundamental representation. The Yang--Baxter equation is the confluence identity for three exchanges. The KZ and Yangian constructions share the geometry of configuration spaces while using different flat connections.

\chapter{The centre and the double}
The Hochschild centre $Z_{\ord}(A)$ controls central operations on the
ordered algebra in its specified coefficient category.
To form a Drinfeld centre of modules, first specify a monoidal module category.
For an ordered chiral bialgebra satisfying
Definition~\ref{def:vol4-ordered-chiral-bialgebra}, use
\[
 \mathcal M_A=\operatorname{Mod}_A(\widehat{\mathcal V}_C)
\]
with the monoidal product of Theorem~\ref{thm:vol4-representation-fusion}.
An object of $Z(\mathcal M_A)$ is a module $M$ together with natural equivalences
\[
 c_{M,N}:M\widehat{\tensor}^{!}N
 \xrightarrow{\sim}N\widehat{\tensor}^{!}M
 \qquad(N\in\mathcal M_A)
\]
satisfying the half-braiding hexagon and unit identities, with their coherences.
This categorical centre is braided.
A comparison with the Hochschild centre must specify a realization functor,
the source and target of the comparison, and the operations it preserves.
Dualizability alone does not supply such a comparison.
Derived loop-space descriptions apply in the geometric categories and
under the hypotheses of their respective centre theorems
\cite{BenZviFrancisNadler}.

The centre and the double answer complementary questions.
The centre classifies central actions in the chosen category.
The double reconstructs a quantum group from paired positive and negative halves.
A field-theory comparison requires explicit realizations of the relevant
line and bulk operations and a map compatible with both constructions.

\part{A compact six-dimensional protected source}

\chapter{Holomorphic cotangent BF theory on \texorpdfstring{$K3\times E$}{K3 x E}}
Let $S$ be a K3 surface with holomorphic symplectic form $\sigma_S$, let $E$ be an elliptic curve with holomorphic form $\dd z$, and put
\[
 X=S\times E,
 \qquad
 \Omega_X=\sigma_S\wedge\dd z.
\]
For a finite proper-height quotient $\nIg^{(H)}$ of the positive Igusa Borcherds algebra, take fields
\[
 A\in\Omega^{0,\bullet}(X,\nIg^{(H)})[1],
 \qquad
 B\in\Omega^{0,\bullet}(X,(\nIg^{(H)})^\vee)[1].
\]
The invariant evaluation pairing and $\Omega_X$ define the BV action
\begin{equation}\label{eq:vol4-compact-hbf}
 S_H(A,B)=
 \int_X\Omega_X\wedge
 \angles{B,\dbar A+\frac12[A,A]}.
\end{equation}
This is the holomorphic cotangent theory of the dg Lie algebra of $\nIg^{(H)}$-valued Dolbeault forms.

\begin{theorem}[Classical compact-background source]
The action \eqref{eq:vol4-compact-hbf} satisfies the classical master equation and defines a holomorphic factorization algebra on the compact Calabi--Yau threefold $X$.
\end{theorem}
\begin{proof}
The shifted cotangent pairing is invariant.  The classical master equation is the Jacobi identity of $\nIg^{(H)}$ together with $\dbar^2=0$ and the derivation property of $\dbar$.  Local observables of an elliptic cyclic $L_\infty$ algebra form a factorization algebra by the standard BV construction.
\end{proof}

\chapter{Tree exactness and the quantum master equation}
Every interaction vertex in \eqref{eq:vol4-compact-hbf} has one $B$-leg and two $A$-legs.  The propagator pairs an $A$-leg with a $B$-leg.

\begin{lemma}[Loop bound]
A connected Feynman graph of the cotangent theory has loop number at most one.
\end{lemma}
\begin{proof}
A graph with $V$ interaction vertices has exactly $V$ internal or external $B$-half-edges.  Every internal edge consumes one of them, hence $E\le V$.  The loop number is $E-V+1\le1$.
\end{proof}

The positive algebra is graded by strictly positive root height.  For homogeneous $x\in\nIg^{(H)}$, the operator $\ad_x$ raises height.  Its coadjoint operator lowers dual height.  A product of adjoint operators labelled by positive roots has total nonzero height shift and therefore has zero trace.

\begin{proposition}[Vanishing of one-loop Lie weights]\label{cyiv:loop-weights}
For every finite positive-height quotient, the Lie weight of each
one-loop graph in the cotangent theory vanishes.
\end{proposition}
\begin{proof}
A one-loop graph has $E=V$.
Every $B$-leg is consumed by an internal edge, so its external legs are
$A$-legs.
Its Lie weight is a cyclic trace of adjoint operators, or the dual
trace with the conventionally equivalent signs.
Each external label has positive root height.
The product therefore strictly raises height on the finite-dimensional
adjoint representation and has zero trace.
The same argument applies to the dual representation.
\end{proof}

This vanishing removes the one-loop Lie coefficients in a perturbative
construction with the stated propagator and vertex types.
A renormalized quantum factorization algebra still requires admissible
analytic kernels, local extensions, the quantum master equation, and
compatible height transitions on the chosen observable spaces.
These are the remaining quantum realization data.

\chapter{Classical reduction along K3}
Choose a K\"ahler metric on $S$ and normalize a closed harmonic
$(0,2)$-form $\eta$ by $\int_S\sigma_S\wedge\eta=1$.
Let $K=\mathbb C[\eta]/(\eta^2)$ with $|\eta|=2$ and zero differential.
There is a unital dg-algebra map
\[
 i:K\longrightarrow\Omega^{0,*}(S),\qquad \eta\longmapsto\eta.
\]
Its image is closed under multiplication because $(0,4)$-forms vanish.
The Dolbeault cohomology of a K3 surface is one-dimensional in degrees
zero and two and zero in degree one, so $i$ is a quasi-isomorphism.
It preserves the trace defined by integration against $\sigma_S$.

Let $\mathfrak l$ be a finite-dimensional Lie algebra, and write
$\mathfrak d=\mathfrak l\ltimes\mathfrak l^*$ for its cotangent Lie
algebra, with invariant evaluation pairing.
For an open $U\subset E$, tensor with $\Omega^{0,*}(U)$, using the
completed smooth tensor product in the $S$ direction.
The field comparison is
\[
 K\otimes\Omega^{0,*}(U)\otimes\mathfrak d
 \xrightarrow{i\otimes\id}
 \Omega^{0,*}(S\times U)\otimes\mathfrak d.
\]
The bracket is multiplication of forms followed by the Lie bracket,
with Koszul signs.

\begin{proposition}[Classical field comparison]\label{cyiv:k3-classical-reduction}
The displayed map is a trace-preserving dg Lie quasi-isomorphism,
compatible with restriction in $U$.
Its induced minimal operations in the K3 factor have no higher products.
\end{proposition}
\begin{proof}
Multiplicativity of $i$ proves bracket compatibility.
Hodge projection $p$ and $h=\bar\partial^*G$ on the compact surface satisfy
$pi=\id$ and $\bar\partial h+h\bar\partial=\id-ip$.
Tensoring these continuous operators with forms on $U$ gives the same
contraction, with the tensor-product signs.
Thus $i\otimes\id$ is a quasi-isomorphism.
Restriction acts only on the $U$ factor and commutes with it.
The normalized trace on $K$ is the restriction of the Dolbeault trace.
Choose the contraction with $hi=0$.
In every higher transferred tree, the first product of inputs from
$i(K)$ remains in $i(K)$ and the next internal homotopy kills it.
Thus every higher transferred product vanishes.
\end{proof}

The coefficient algebra $K$ retains both the degree-zero and degree-two
K3 modes.
The proposition compares classical field complexes and their brackets.
Constructing an equivalence of quantum observable factorization algebras
requires compatible quantization, support conventions, and the induced
maps on observable completions.
The ordinary Hochschild comparison of Theorem~\ref{cyiv:hh-ce} is a
separate algebraic calculation, not that quantum field map.

\chapter{Protected states and the Igusa trace}
Let
\[
 \mathcal H_{\mathrm{1p}}
 =\widehat\bigoplus_{\alpha>0}\gIg_\alpha
\]
with its Borcherds parity.  The multiparticle state space is the completed super-Fock space
\[
 \mathcal H_{\mathrm{Fock}}=\widehat\Sym(\mathcal H_{\mathrm{1p},\bar0})
 \tensor\widehat\Lambda(\mathcal H_{\mathrm{1p},\bar1}).
\]
For $Z$ in the absolute-convergence chamber, let $e^{-H_Z}$ act by $e^{-\alpha(Z)}$ on the root space of degree $\alpha$.

\begin{theorem}[Protected Hamiltonian trace]
The supertrace is
\[
 \str_{\mathcal H_{\mathrm{Fock}}}(e^{-H_Z})
 =\prod_{\alpha>0}
 (1-e^{-\alpha(Z)})^{-\sdim\gIg_\alpha}.
\]
After multiplication by the inverse Weyl line, its reciprocal is $\DenIg(Z)=D_5(2Z)$.  After the half-coordinate pullback and doubling, the disconnected protected scalar is $-(\jmath^*\DenIg)^{-2}=-\chi_{10}^{-1}$ in the Oberdieck--Pixton normalization.
\end{theorem}
\begin{proof}
An even oscillator contributes $(1-e^{-\alpha})^{-1}$ and an odd oscillator contributes $1-e^{-\alpha}$.  Multiplication over the root basis gives the product.  The denominator identity and the reduced curve-counting theorem supply the final two equalities.
\end{proof}

\chapter{The compact-source comparison}
The compact threefold provides the classical cotangent field theory
and the dg Lie comparison of Proposition~\ref{cyiv:k3-classical-reduction}.
The root datum separately gives the Borcherds algebra, its enveloping
algebra, and the root-graded Fock space.
The latter has the trace computed above.
Proposition~\ref{cyiv:loop-weights} supplies the finite one-loop Lie
cancellation, while Theorem~\ref{cyiv:hh-ce} computes the ordinary
algebraic centre model with adjoint coefficients.

\begin{problem}[Geometric compact comparison]\label{cyiv:compact-comparison}
Construct objects $E_i$ representing the simple branes and maps on their
derived morphism complexes that preserve units, composition, and higher
composition homotopies.
Realize these maps on extension stacks with the reduced critical
coefficients and orientation isomorphisms on triple flags.
Construct compatible tensor and unit constraints on the resulting
boundary functor.
Compare the chosen quantum bulk action with the independently defined
boundary centre, including support, collisions, and completion.
Finally construct a trace and a map of its value line to the automorphic
line in the same charge normalization.
\end{problem}

A primitive map satisfying the root relations gives the enveloping and
Chevalley maps of Theorem~\ref{cyiv:primitive-comparison}.
A strong monoidal equivalence transports categorical centres.
Neither construction determines the brane functor or its cyclic trace.
These are the additional maps required for compact geometric
reconstruction.

\part[Physics]{Physics of the comparison maps}

\chapter{Holomorphic--topological boundary theories}
A three-dimensional holomorphic--topological theory on $C\times\R$ has local observables that are holomorphic along $C$ and locally constant along $\R$. A boundary produces an ordered chiral algebra. A second topological direction permits exchange and gives an $\Etwo$ enhancement or a meromorphic braided line category. The bulk acts through the derived centre.

The finite cotangent BF model gives the exact local calculation:
\[
 \Barc(U(\mathfrak l))\simeq C_*(\mathfrak l).
\]
For $\mathfrak l=\nDelta^{(H)}$, this calculation realizes every finite-height Igusa denominator sector.

\chapter{Six-dimensional hCS on \texorpdfstring{$K3\times E$}{K3 x E}}
Choose a holomorphic volume form $\Omega_S\wedge dz_E$. The hCS theory has classical fields
\[
 \Omega^{0,*}(S\times E,\mathfrak a)[1]
\]
and action $S_{\hCS}$. Its factorization observables are native to the threefold. Pushforward along $p_E$ gives a factorization algebra on $E$ whose coefficient complex contains the Dolbeault modes of $S$.

A protected reduction to a chiral algebra consists of a deformation retract of the K3 Dolbeault complex, compatibility of the transferred $L_\infty$ operations with the trace, a boundary or defect selecting ordered fusion, and a quantization trivializing the local anomaly. These data produce the ordered chiral algebra to which the Hall/BPS comparison may be applied.

The Hodge equality $\chi(\Ocal_{S\times E})=0$ controls one compact supertrace. The gauge and gravitational anomalies are classes in the local BV deformation complex and are evaluated by their own characteristic forms.

\chapter{Omega backgrounds}
Let a torus $T$ act holomorphically on the local geometry and preserve the selected supercharge. For a generator $V$ of the action, the equivariant differential is
\[
 Q_\Omega=Q+\iota_V,
 \qquad Q_\Omega^2=\mathcal L_V.
\]
After localization, fixed loci produce deformation parameters and often a spectral $R$-matrix. Toric threefolds and ADE surface resolutions carry such actions. A deformation family equipped with monodromy-compatible moduli correspondences and wall-crossing equivalences transports the localized algebra from algebraic or equivariant K3 fibers to the selected generic fiber.

\chapter{Anomalies as typed classes}
The principal anomaly classes are
\[
\begin{array}{c|c}
\text{construction}&\text{deformation complex}\\ \hline
\text{hCS quantization}&\text{local BV cochains}\\
\text{critical CoHA orientation}&\text{Picard/gerbe theory of the critical stack}\\
\text{Hall coproduct}&\text{correspondence and localization complex}\\
\text{projective conformal blocks}&\text{Atiyah algebra of the determinant line}\\
\text{automorphic multiplier}&\text{character group of the modular cover}.
\end{array}
\]
A comparison theorem is a morphism between the corresponding complexes together with equality of the transported classes. This formulation turns numerical coincidences into geometric identities precisely when an intertwining map exists.

\part[Automorphic forms]{Automorphic characters and scalar shadows}

\chapter{Borcherds products as characters}
Let $\mathfrak n_+$ be a locally finite positive Lie superalgebra. The Chevalley Euler character is
\[
 \chi C_*(\mathfrak n_+)
 =\prod_{\alpha>0}(1-e^{-\alpha})^{\sdim\mathfrak g_\alpha}.
\]
A Borcherds product becomes the Weyl-shifted character of $C_*(\mathfrak n_+)$. This construction produces a character of an existing algebra. Hall geometry supplies a second algebra, and a primitive comparison theorem identifies the two.

\chapter{The universal Borcherds weight formula}
For a weight-zero weak Jacobi form
\[
 \varphi(\tau,z)=\sum c(n,\ell)q^nr^\ell,
\]
the Borcherds lift has weight
\[
 \wt(\operatorname{Borch}(\varphi))=\frac{c(0,0)}2.
\]
For the K3 half-index, $c(0,0)=10$ and the weight is $5$. This formula lives entirely on the automorphic layer. Hodge Euler characteristics, Heisenberg ranks, central charges, and BV anomalies are evaluated by their own traces.

\chapter{The Gritsenko--Cl\'ery atlas}
The eight diagonal-divisor forms are
\[
 \Delta_5,\Delta_2,\Delta_1,\Delta_{1/2},
 \nabla_3,\nabla_2,\nabla_{3/2},Q_1
\]
with weights
\[
 5,2,1,\frac12,3,2,\frac32,1.
\]
Each form is an automorphic product on its own paramodular or Hecke congruence group with its own character or multiplier \cite{GC}. A geometric realization is specified row by row by a charge lattice, Hall problem, trace, and primitive algebra comparison.

\chapter{Scalarization hierarchy}
The constructions descend through four levels:
\[
\begin{aligned}
 \text{factorization or Hall object}
 &\longrightarrow \text{primitive Lie algebra}\\
 &\longrightarrow \text{PBW character}
 \longrightarrow \text{scalar partition function}.
\end{aligned}
\]
Each arrow retains the quotient of structure displayed at its target. The first object contains supports, extension maps, and brackets. The second retains the infinitesimal algebra. The third retains signed graded dimensions. The fourth adds a trace and normalization. Upward reconstruction is supplied by the corresponding geometric operations and comparison morphisms.

\part{The unified theorem system}

\chapter{The foundational theorem}
\begin{definition}[Complete Calabi--Yau quantum-group datum]
A complete comparison datum specifies independent categorical, field,
boundary, Hall, Hopf, and trace objects together with:
\begin{enumerate}
\item PTVV-representable moduli for the smooth proper Calabi--Yau category;
\item a renormalized BV quantization satisfying the quantum master equation and a local compact factorization generator $b$;
\item a boundary reduction or proper locally trivial map to a curve satisfying factorization pushforward, together with the supplied curve-level faces and aligned distributive law;
\item a strong-oriented Hall coefficient theory whose pull--push maps satisfy base change and whose positive algebra is separated and complete;
\item paired negative and Cartan halves, a completed coproduct, a continuous nondegenerate Hopf pairing, and a separated triangular completion on which the canonical tensors converge;
\item compatible open--closed and Hall--boundary morphisms;
\item a parity-refined primitive comparison for every automorphic identification.
\end{enumerate}
\end{definition}

\begin{theorem}[Calabi--Yau quantum-group assembly]
Let a complete datum consist of:
\begin{enumerate}
\item a Calabi--Yau category and a representable moduli-of-objects problem;
\item a quantized cyclic elliptic BV theory and a boundary or defect generator;
\item a Hall extension correspondence with coefficient functor, base-change theorem, strong orientation data, and separated charge completion;
\item paired positive and negative bialgebras, Cartan data, a nondegenerate continuous Hopf pairing, and a completed double;
\item every curve-level distributive, PBW, and analytic comparison used below.
\end{enumerate}
Then the datum determines functorially:
\begin{enumerate}
\item framed $\Etwo$ Hochschild cochains and the shifted-symplectic moduli structure;
\item the three curve-level algebraic faces and their three bars;
\item the positive Hall/BPS algebra and its primitive Lie algebra;
\item the completed quantum double and braided representation category;
\item every supplied comparison among categorical, field-theoretic, Hall, current, quantum, and automorphic realizations;
\item the Weyl-shifted PBW/Chevalley character of any parity-refined primitive Lie-superalgebra identification.
\end{enumerate}
\end{theorem}
\begin{proof}
Cyclic Deligne and shifted symplectic geometry give the first output \cite{BravRozenblyum,PTVV,BravDyckerhoff}.  Boundary factorization gives the curve faces, and Volume~I gives their bars and centres.  The Hall correspondence and its base-change theorem give the associative product; strong orientation supplies the coefficient coherence required in the $3$-Calabi--Yau setting \cite{KPS}.  The paired bialgebra data give the double and universal $R$-matrix.  The comparison maps commute by hypothesis.  PBW and Chevalley antisymmetrization give the final character.  Each construction is natural under a morphism preserving the displayed data.
\end{proof}

The local quiver, equivariant surface, K3 correspondence, and Igusa Borcherds constructions have the carriers stated above.
A geometric $K3\times E$ comparison requires the primitive map and determinant-line map of Problem~\ref{cyiv:compact-comparison}.

\chapter{Landscape of established realizations}
\begin{center}
\footnotesize
\begin{tabularx}{0.99\textwidth}{@{}>{\raggedright\arraybackslash}p{0.15\textwidth}>{\raggedright\arraybackslash}p{0.25\textwidth}>{\raggedright\arraybackslash}X>{\raggedright\arraybackslash}p{0.20\textwidth}@{}}
\toprule
\textbf{Geometry}&\textbf{Positive algebra}&\textbf{Full or centred object}&\textbf{Character}\\
\midrule
$A_2$ quiver&$U_v^+(\mathfrak{sl}_3)$&$U_v(\mathfrak{sl}_3)$&PBW root product\\
$\C^3$&$Y^+(\widehat{\mathfrak{gl}}_1)$&$Y(\widehat{\mathfrak{gl}}_1)$&plane-partition and Fock traces\\
$\widetilde{\C^2/\Gamma}$&$Y^+(\widehat{\mathfrak g}_\Gamma)$&paired completion when the negative half and coproduct are supplied&surface Fock traces\\
$K3$&$U(\mathfrak g_{\BPS}^+)$&Heisenberg and LLV actions&K3 Hall characters\\
$K3\times E$&Borcherds root envelope; geometric reduced character&$U_q(\gDelta)$ and the root-current model and classical BF theory&$-\ChiTen^{-1}$\\
\bottomrule
\end{tabularx}
\end{center}
The $A_2$ and $\C^3$ rows provide complete standard quantum-group models. The ADE row provides an established completed positive half together with its geometric current actions; its paired full double is formed from the negative half, Cartan data, coproduct, continuous Hopf pairing, and triangular completion stated in the table. The K3 row provides global two-Calabi--Yau BPS and cohomological symmetries. The final row records independently constructed root, character, and classical field objects.
Their geometric Hall realization requires the comparison in Problem~\ref{cyiv:compact-comparison}.

\chapter{Constructed comparisons and their geometric extensions}
The Hall flag argument constructs associativity once its coefficient
and orientation maps are defined.
The tensor-action formula constructs representation fusion from a
coassociative counital coproduct.
The matrix-coefficient formulas construct the coend of a supplied
fibre functor.
The primitive presentation gives enveloping and Chevalley maps.
The classical cyclic and K3 field calculations retain their actual
differentials, traces, and brackets.

The quantum field realization requires admissible quantization and
observable completions.
The compact Hall realization requires the brane and orientation maps
of Problem~\ref{cyiv:compact-comparison}.
The Fredholm and automorphic comparisons require their analytic domains
and value-line isomorphisms.
These constructions preserve the full geometric reconstruction problem
beyond its finite algebraic and classical calculations.

\chapter{Actions and comparison maps}
For a specified boundary algebra $A_b$, a coherent bulk action is a map
\[
 \beta:\Obs_{\mathrm{bulk}}|_\partial\longrightarrow Z_{\ord}(A_b).
\]
The target is the derived endomorphism object of the diagonal bimodule
in the chosen factorization coefficient category.
A Hall construction instead starts with
\[
 \M_\alpha\times\M_\beta\xleftarrow{p}\mathcal E_{\alpha,\beta}
 \xrightarrow{q}\M_{\alpha+\beta}.
\]
Critical coefficients, orientations, and admissible pull--push maps
give its product.
Framed moduli supply a representation after their action
correspondences and triple-flag identities have been constructed.
A paired Hopf structure gives $D(H^+,H^0,H^-)$.
Its coproduct defines tensor products of continuous modules, and a
convergent universal $R$-matrix gives exchange.

A strong monoidal equivalence between this module category and a
boundary category transports categorical centres.
It does not identify that categorical centre with the object
$Z_{\ord}(A_b)$ without a realization theorem.
A trace on a specified representation then gives an amplitude.
Its comparison with an automorphic section requires a map of value
lines with the same normalization.

\chapter{Comparison data for Calabi--Yau constructions}
Comparisons among Calabi--Yau categories, holomorphic field theory, chiral algebras, Hall algebras, quantum groups, and automorphic products require the following structures:
\begin{enumerate}
\item for a smooth proper $d$-Calabi--Yau category with its nondegenerate negative-cyclic class, the categorical outputs are framed $\Etwo$ Hochschild cochains and, when the moduli stack of objects is representable in the PTVV setting, $(2-d)$-shifted symplectic moduli; the shift $d$ does not change the Hochschild operadic level to $\En$;
\item a geometric BV theory and boundary category produce factorization observables and transverse ordered chiral algebras; associative chiral convolution is recorded as a separate supplied structure;
\item pushforward to a curve uses an actual map $p:X\to C$; formal cycle-specialization notation alone specifies no pushforward;
\item extension stacks, coefficient sheaves, orientations, and completions define the Hall algebra;
\item coproduct, Cartan, Hopf pairing, and topology define the double;
\item K3 Heisenberg, LLV orthogonal symmetry, local ADE Yangians, and global BPS generalized Kac--Moody algebras occupy distinct strata joined by explicit comparison maps where constructed;
\item the graph of the $K3\times E$ Gram map becomes multiplicative in the polarized central extension;
\item point-local and wrapped sectors form a hybrid factorization problem;
\item the Igusa form is the Weyl-shifted PBW/Chevalley character of $\gDelta$, and compact Hall geometry maps to it through the primitive comparison;
\item anomaly identities are transported through explicit maps between deformation complexes.
\end{enumerate}
The categorical, geometric, Hall, and automorphic operations thus have separate carriers and comparison maps.

\part{Technical appendices to the preceding book}
\chapter{Typed data table}
\begin{center}
\small
\begin{tabularx}{0.98\textwidth}{@{}>{\raggedright\arraybackslash}p{0.20\textwidth}>{\raggedright\arraybackslash}X>{\raggedright\arraybackslash}p{0.28\textwidth}@{}}
\toprule
\textbf{Object}&\textbf{Minimal data}&\textbf{Primary operation}\\
\midrule
$\CY$ category&nondegenerate negative-cyclic trace&Hochschild and shifted-symplectic structures\\
Factorization algebra&local observables and Weiss descent&disjoint-open product\\
Transversely ordered chiral algebra&factorization coefficient with an internal $\Eone$ structure for $\tensor^!$&ordered fusion\\
Associative chiral algebra&an $\Eone$ structure for chiral convolution $\starCh$&nonlocal chiral composition\\
$\CoHA$&extension correspondence and coefficient theory&pull--push product\\
Bialgebra&compatible product and coassociative counital coproduct&fusion of representations\\
Double&positive and negative halves, Cartan data, Hopf pairing&universal $R$-matrix\\
Borcherds product&Jacobi or vector-valued modular input&graded character\\
\bottomrule
\end{tabularx}
\end{center}

\chapter{The three K3 brackets}
The Heisenberg bracket is the excess intersection of two Nakajima incidence kernels.  The LLV bracket is the commutator of diagonal Lefschetz kernels and their adjoints.  The local Yangian bracket is the convolution commutator of ADE Hecke kernels with spectral equivariance.  In $\operatorname{Corr}_{K3}(S)$ all three are compositions of geometric kernels.  Their mixed commutators are the corresponding triple-intersection classes.  The action algebra and Tannaka coend retain these classes while preserving the three geometric constructions.

\chapter{Exact arithmetic checks}
The companion verification calculation establishes:
\begin{enumerate}
\item the low Fourier rows of $\phi_{0,1}$ through $q^3$ by exact theta-series division;
\item the vector-valued cocycle identity for the Gram polarization;
\item the rational Yang--Baxter equation on $(\C^2)^{\tensor3}$;
\item the $A_2$ Hall flag counts at ranks zero and one;
\item the type-II Gram matrix and reflection identities.
\end{enumerate}
Each calculation uses integer or rational arithmetic and appears as a worked formula in the relevant chapter.

\cleardoublepage
\phantomsection
\addcontentsline{toc}{part}{Universal Chiral BV and Einstein Completion}
\chapter*{Coherent boundary actions and BV fields}
\addcontentsline{toc}{chapter}{Coherent boundary actions and BV fields}
\begin{summarybox}
Boundary operations determine a deformation problem once their interface and coherence maps are supplied. A local realization gives a shifted cotangent BV theory. Comparing it with a physical bulk requires a map of fields, actions, and coherent boundary operations.
\end{summarybox}
\part{The full mixed boundary datum}

\chapter{Three products and one interface}
Holomorphic collisions and transverse composition act on different coefficient objects. A common boundary theory must specify their interaction on an interface. This interaction is part of the input to the following constructions.

Let $X$ be a smooth complex curve and let $\Dcal_X=D\text{-}\Mod(\Ran X)$ in a fixed six-functor enhancement.  Write $\tensor^!$ for the pointwise product and $\starCh$ for chiral convolution.

\begin{definition}[Complete mixed boundary datum]
A complete mixed boundary datum on $X$ is a tuple
\[
 \mathbb A=(L,A^{\ch},A^\perp,K,\iota,\lambda)
\]
with the following components.
\begin{enumerate}
\item $L$ is a complete chiral Lie algebra.
\item $A^{\ch}$ is a complete associative algebra for $\starCh$.
\item $A^\perp$ is a complete $\Eone$ algebra in $(\Fact_D(X),\tensor^!)$.
\item $K$ is an $(A^{\ch},A^\perp)$ interface bimodule.
\item $\iota:L\to(A^{\ch})_{\Lie}$ is a chiral Lie morphism.
\item $\lambda$ is the aligned mixed distributive law relating the two module actions on $K$.
\end{enumerate}
Every structure map is continuous for the chosen weight filtrations. The tuple includes the module actions, the mixed coherence homotopies, and their identities in the chosen enhancement. When a reduced bar is used, the relevant augmentations are supplied as additional data. Existence of this tuple for a physical boundary theory is a separate construction.
\end{definition}

The three native bars are
\[
 \BarLieCh(L),\qquad
 \BarAssCh(A^{\ch}),\qquad
 \BarPerp(A^\perp).
\]
The interface has a fourth bar.  Its degree-$r$ terms are alternating words in the two associative colours with one marked interface letter.  Its differential contracts adjacent compatible letters.

\begin{construction}[Interface bar]
Let $\overline A^{\ch}$ and $\overline A^\perp$ be the augmentation ideals.  Set
\[
 \BarMix(\mathbb A;K)
 =\bigoplus_{p,q\ge0}
 (s\overline A^{\ch})^{\tensor p}\tensor sK\tensor
 (s\overline A^\perp)^{\tensor q}.
\]
The differential is the sum of the internal differential, the two bar differentials, the left action on $K$, the right action on $K$, and the mixed terms supplied by $\lambda$.
\end{construction}

\begin{theorem}[Four-colour differential]
The interface differential squares to zero.  Its codimension-two terms are exactly the associative relations in the two colours, the two module relations, and the mixed distributive relation.
\end{theorem}
\begin{proof}
Filter by total word length.  Every component of the square contracts two adjacent edges.  Disjoint contractions cancel by the determinant-line sign.  Nested contractions in one colour cancel by associativity.  A contraction meeting the interface cancels by the corresponding module axiom.  The remaining square has one edge of each colour and is the aligned distributive square.  The datum $\lambda$ identifies its two composites.
\end{proof}

\chapter{Compatible modules and the mixed centre}
An operation on a boundary theory must act on boundary-changing morphisms as well as on objects. Natural transformations retain this compatibility, and derived natural transformations retain its higher homotopies.

A compatible module is a triple $(M^{\ch},M^\perp,\kappa)$ with an $A^{\ch}$ action, an $A^\perp$ action, and an interface comparison over $K$.  Let $\Mod_{\mathbb A}^{\mathrm{mix}}$ be their stable category.

\begin{definition}[Mixed derived centre]
The mixed derived centre is
\[
 Z_{\mathrm{mix}}(\mathbb A)
 =\End\!\left(\id_{\Mod_{\mathbb A}^{\mathrm{mix}}}\right).
\]
It is computed by the endomorphisms of the diagonal compatible bimodule.
\end{definition}

\begin{theorem}[Universal interior action]
$Z_{\mathrm{mix}}(\mathbb A)$ carries an internal $\Etwo$ algebra structure.  Every action already defined as a coherent natural endomorphism of the identity determines a point of this centre. The space of its factorizations through the representing endomorphism object is contractible.  Morita-equivalent complete mixed data have equivalent centres.
\end{theorem}
\begin{proof}
Endomorphisms of the identity functor form the higher centre of a module category.  Composition and insertion of natural transformations give the brace operations and hence the $\Etwo$ structure.  A supplied coherent natural action defines an endomorphism of the identity. A physical bulk insertion determines such an action only after its action on modules, morphisms, and higher compositions has been constructed.  The factorization is the universal property of the endomorphism object.  Morita equivalence identifies the module categories and their identity functors.
\end{proof}

\chapter{Relative induction from a chiral algebra}
Let $V$ be a chiral algebra with conformal vacuum.  Its algebraic collar $\Coll(V)$ is the free transverse $\Eone$ extension generated by one boundary copy of $V$, modulo the chiral relations already present in $V$.  Weight completion gives $\widehat\Coll(V)$.

\begin{theorem}[Universal collar]
For every complete transversely ordered chiral algebra $B$ and every chiral morphism $V\to B$, there is a unique continuous $\Eone$ morphism
\[
 \widehat\Coll(V)\longrightarrow B
\]
extending the chiral map.  The construction is left adjoint to restriction from transverse ordered algebras to chiral coefficients.
\end{theorem}
\begin{proof}
The algebraic collar is the relative operadic induction from the chiral operad to the mixed operad.  Relative induction is left adjoint to restriction.  Completion preserves the adjunction on positive complete objects because every map is determined weight by weight.
\end{proof}

When relative induction supplies a regular interface and a continuous aligned distributive law, these data define a mixed envelope with $L=V_{\Lie}$, $A^{\ch}=U^{\ch}(L)$, and $A^\perp=\widehat\Coll(V)$. Constructing that interface law in the selected geometric coefficient category is part of the envelope construction.

\part{The deformation algebra and its BV envelope}


\chapter{The centre-extended deformation algebra}
Varying a boundary product changes its compatibility with the interface and the central action. The deformation complex must therefore include all these operations. A local field realization is further data needed to turn their deformation equations into a BV action.

Fix a cofibrant coloured operad $\mathcal P_{\mathrm{mix}}$ presenting the structures of a complete mixed datum and its central action.  Let $\mathcal P_{\mathrm{mix}}^{\ash}$ be its bar cooperad.  The convolution complex
\[
 \Def_{\mathbb A}^{\mathrm{full}}
 =\Hom_{\mathbb S}
 \left(\mathcal P_{\mathrm{mix}}^{\ash},\End_{\mathbb A}\right)
\]
contains the chiral operation, both associative products, the interface, and the centre action.

\begin{theorem}[Complete deformation $L_\infty$ algebra]
The convolution complex has a canonical complete $L_\infty$ structure.  Its Maurer--Cartan points are complete mixed boundary structures together with a coherent central action.  Twisting by the point defined by $\mathbb A$ gives the tangent algebra of deformations of the full datum.
\end{theorem}
\begin{proof}
The infinitesimal decomposition maps of the bar cooperad insert collections of operations into one another.  Composition in the endomorphism operad turns these decompositions into the convolution brackets.  The cooperad identities imply the $L_\infty$ identities.  A Maurer--Cartan element is precisely an operad morphism from a cofibrant resolution of $\mathcal P_{\mathrm{mix}}$ to $\End_{\mathbb A}$.  Twisting linearizes this mapping space at the selected structure.
\end{proof}

\section{Field equations and Bianchi identities}
Let $\alpha$ be a degree-one field in the localized deformation algebra.  Its curvature is
\[
 \mathcal F(\alpha)
 =\sum_{n\ge1}\frac1{n!}\ell_n(\alpha^{\tensor n}).
\]
The linearization at $\alpha$ is
\[
 d_\alpha x
 =\sum_{n\ge0}\frac1{n!}
   \ell_{n+1}(\alpha^{\tensor n},x).
\]

\begin{theorem}[Universal Bianchi identity]
For every $\alpha$,
\[
 d_\alpha\mathcal F(\alpha)=0.
\]
At a Maurer--Cartan point, $d_\alpha^2=0$.
\end{theorem}
\begin{proof}
Both identities are the exponential form of the $L_\infty$ Jacobi relations.  Expand the coderivation square on the cofree cocommutative coalgebra and evaluate it on $e^\alpha$ and $e^\alpha x$.
\end{proof}

\chapter{Shifted cotangent transgression}
Let $\flie$ be a complete elliptic $L_\infty$ algebra on a three-manifold $M$.  Its shifted cotangent theory has fields
\[
 \mathcal E=T^*[-1]\flie[1]
 =\flie[1]\oplus\flie^\vee[-2].
\]
The canonical action is
\[
 S(\alpha,\beta)
 =\angles{\beta,\mathcal F(\alpha)}.
\]

\begin{theorem}[Universal BV envelope]
The action satisfies the classical master equation.  Its Euler--Lagrange equations are
\[
 \mathcal F(\alpha)=0,
 \qquad
 d_\alpha^\vee\beta=0.
\]
The tangent complex at a solution is the cotangent complex of the Maurer--Cartan moduli problem.
\end{theorem}
\begin{proof}
The Hamiltonian vector field of $S$ is the cotangent lift of the homological vector field defining the $L_\infty$ structure.  The square of that vector field is zero, hence $\{S,S\}=0$.  Variation in $\beta$ gives the Maurer--Cartan equation.  Variation in $\alpha$ gives the dual linearized equation.
\end{proof}

\begin{definition}[Local realization and chiral BV envelope]
A local realization of the complete mixed datum consists of an elliptic local $L_\infty$ algebra $\flie_{\mathbb A}^{\mathrm{full}}$ on $M=X^{\mathrm{an}}\times\R$, a specified continuous filtered $L_\infty$ quasi-isomorphism
\[
 \Gamma(M,\flie_{\mathbb A}^{\mathrm{full}})
 \longrightarrow \Def_{\mathbb A}^{\mathrm{full}},
\]
and the local pairing needed for cotangent transgression. Here $\Gamma$ denotes the chosen complex of global sections with its declared topology. The comparison includes compatibility with the specified boundary operations. The dual fields use the density-valued continuous dual, with compact support or boundary conditions that make the integrated pairing and integration by parts defined.

For a supplied local realization, define
\[
 \BVEnv_3^{\mathrm{full}}(\mathbb A)
 =T^*[-1]\MC(\flie_{\mathbb A}^{\mathrm{full}}).
\]
Its action is the cotangent action displayed above. Constructing a local realization from an abstract mixed datum remains an additional existence problem.
\end{definition}

\chapter{Intrinsic algebraic quantization}
Let $\mathfrak X$ be a finite perfect derived formal moduli problem with a homological orientation $\mu$: the corresponding Berezinian is invariant under the homological vector field.  Functions on its shifted cotangent are completed polyvectors:
\[
 \Ocal(T^*[-1]\mathfrak X)
 =\widehat{\PV}(\mathfrak X)
 =\prod_{p\ge0}\Gamma(\mathfrak X,\wedge^pT_{\mathfrak X})[p].
\]
The homological orientation defines the divergence operator
\[
 \Delta_\mu:\widehat{\PV}(\mathfrak X)
 \longrightarrow\widehat{\PV}(\mathfrak X)[1].
\]

\begin{theorem}[Oriented BV complex]
The operator $\Delta_\mu$ has order two and satisfies
\[
 \Delta_\mu^2=0,
 \qquad
 [d_{\mathfrak X},\Delta_\mu]=0,
\]
\[
 \Delta_\mu(fg)=\Delta_\mu(f)g+(-1)^{|f|}f\Delta_\mu(g)
 +(-1)^{|f|}\{f,g\}.
\]
Consequently
\[
 \QHdg(\mathfrak X,\mu)
 =\left(\widehat\PV(\mathfrak X)\llbracket\hbar\rrbracket,
 d_{\mathfrak X}+\hbar\Delta_\mu\right)
\]
is a square-zero quantum BV complex.
\end{theorem}
\begin{proof}
In oriented Darboux coordinates, $\Delta_\mu$ is the sum of paired second derivatives.  Their graded commutation gives $\Delta_\mu^2=0$.  Invariance of the Berezinian under the homological vector field gives $[d_{\mathfrak X},\Delta_\mu]=0$.  Expanding the operator on a product gives the bracket as the cross term.
\end{proof}

The orientation identifies polyvectors with differential forms by contraction.  Under this identification,
\[
 \iota_\mu\Delta_\mu=d_{\mathrm{dR}}\iota_\mu.
\]
Hence the BV complex is the Hodge-completed derived de Rham complex with differential
\[
 d_{\mathfrak X}+\hbar d_{\mathrm{dR}}.
\]
The Rees parameter records Hodge degree.

\begin{definition}[Intrinsic quantum complex of a mixed datum]
For a complete mixed datum $\mathbb A$ equipped with compatible homological orientations on its finite perfect charts, define $\QHdg_3^{\mathrm{full}}(\mathbb A)$ by applying the preceding BV complex to finite perfect charts of
\[
 \MC(\flie_{\mathbb A}^{\mathrm{full}})
\]
and taking the separated inverse limit over the chart and weight filtrations.
\end{definition}

\begin{theorem}[Quantum master equation]
A degree-zero action $S=S_0+\hbar S_1+\cdots$ defines a quantum state $e^{S/\hbar}\mu$ precisely when
\[
 dS+\frac12\{S,S\}+\hbar\Delta_\mu S=0.
\]
The equation is equivalent to
\[
 (d+\hbar\Delta_\mu)e^{S/\hbar}=0.
\]
\end{theorem}
\begin{proof}
Apply the second-order product formula to the exponential and collect powers of $\hbar$.
\end{proof}

\chapter{Perturbative factorization realization}
The algebraic quantum differential does not specify a propagator or a space of singular kernels. Fix a local realization, elliptic gauge-fixing data, support and boundary conditions, and a renormalization prescription. A heat-kernel propagator is defined away from the diagonal. Extending its graph integrals to collisions requires counterterms and compatibility with the boundary conditions.

\begin{definition}[Factorization realization data]
A factorization realization consists of renormalized local observable complexes, continuous disjoint-insertion maps, and coherent homotopies for their compositions. It includes extensions of graph weights to the required collision strata, their Stokes identities, and a quantum master equation in the permitted local functional complex. Weiss descent is a further required property of these actual complexes. A boundary realization also includes a comparison with the supplied mixed datum, preserving its two products, Lie operation, interface, and mixed law.
\end{definition}

Under these hypotheses the disjoint-insertion maps and composition homotopies define the stated factorization object. The boundary comparison identifies its boundary operations with those of $\mathbb A$. The construction problem is to produce these data from the physical fields and action. A graphwise identity supplies the corresponding composition equation, whereas descent and boundary identification require their own maps and proofs.

\part{Boundary duality, exchange, and reduction}


\chapter{The Koszul--Morita square}
A change from modules to comodules must transport the interface as well as each algebra. The mixed bar supplies the comparison object. Exchange then asks for motion in an additional topological direction, while BRST reduction asks whether the differential respects every transported operation.

Let $P$ be a dualizable boundary state for $A^\perp$.  Its boundary dual is
\[
 A_P^!=\RHom_{A^\perp}(P,P).
\]
The ordinary bar construction gives explicit twisted tensor functors.
Their extension to the mixed datum must also transport the interface and
its geometric operations.

Let $k$ be a field and let $A_{\mathrm{alg}}=\bigoplus_{w\ge0}A_w$
be an augmented associative algebra in cohomological degree zero.
Assume $A_0=k$, augmentation projects onto $A_0$, and multiplication adds weights.
Set $A=\prod_{w\ge0}A_w$ and $C=\prod_{w\ge0}B(A_{\mathrm{alg}})_w$,
with degreewise products and the weight topology.
The bar suspension lowers cohomological degree by one, with $b_2[a|b]=[ab]$.
The degree-one twisting map $\tau:C\to A$ sends $[a]$ to $a$
and vanishes on the empty word and longer words.

A module or comodule has the form $\prod_{w\ge w_0}M_w$,
with differential preserving weight, for some integer $w_0$.
Actions and coactions preserve weight.
Comodule coactions are counital and coassociative, and use the tensor product
\[
 (C\widehat\otimes M)_w
 =\bigoplus_{u+v=w}B(A_{\mathrm{alg}})_u\otimes M_v.
\]
The same convention defines $A\widehat\otimes M$.
These sums are finite in each weight.
Repeated reduced coactions vanish in weight $w$ after more than $w-w_0$ steps.
Invert quasi-isomorphisms in each weight to obtain $D_{\mathrm{wt}}(A)$
and $D_{\mathrm{wt}}(C)$.

\begin{proposition}[Ordinary twisted tensor equivalence]
The functors
\[
 F(M)=C\widehat\otimes_\tau M,\qquad
 G(N)=A\widehat\otimes_\tau N
\]
induce inverse equivalences $D_{\mathrm{wt}}(A)\simeq D_{\mathrm{wt}}(C)$.
Writing $\Delta c=\sum c_{(1)}\otimes c_{(2)}$
and $\rho n=\sum n_{(-1)}\otimes n_{(0)}$, their differentials are
\[
\begin{aligned}
 d_F(c\otimes m)={}&d_Cc\otimes m+(-1)^{|c|}c\otimes d_Mm
 +\sum(-1)^{|c_{(1)}|}c_{(1)}\otimes\tau(c_{(2)})m,\\
 d_G(a\otimes n)={}&a\otimes d_Nn
 -\sum a\tau(n_{(-1)})\otimes n_{(0)}.
\end{aligned}
\]
The counit and unit are
\[
 \epsilon_M(a\otimes c\otimes m)=a\epsilon_C(c)m,\qquad
 \eta_N(n)=\sum n_{(-1)}\otimes1\otimes n_{(0)}.
\]
\end{proposition}
\begin{proof}
The twisting equation makes both differentials square to zero.
The action, coaction, and counit identities show that the displayed maps
are chain maps and satisfy the triangle identities.
For $GF(M)$, the extra degeneracy of the two-sided bar inserts
$1[\bar a_0|a_1|\cdots|a_r]m$ into $a_0[a_1|\cdots|a_r]m$,
where $\bar a_0=a_0-\epsilon_A(a_0)1$.
With simplicial boundary $\sum_i(-1)^id_i$, it contracts the augmented bar.
Multiplying length-$r$ words by $(-1)^r$ gives the displayed differential convention.
This contraction preserves weight, so the counit is a weightwise quasi-isomorphism.
For $FG(N)$, filter each weight by the weight of the final $N$-factor.
The reduced coaction lowers that weight.
The associated graded unit is the contracted right-sided bar
$k\otimes N_v\to B(k,A_{\mathrm{alg}},A_{\mathrm{alg}})\otimes N_v$.
The filtration is finite in each total weight, so the unit is also a quasi-isomorphism.
The same finite filtrations show that $F,G$ preserve weightwise quasi-isomorphisms,
since tensoring complexes over $k$ preserves them.
Thus the unit and counit give the asserted inverse equivalences.
\end{proof}

For a right module $R$ and a left module $L$ with the same lower-weight
bounds, the two-sided bar has underlying weight-completed complex
\[
 \prod_n\ \bigoplus_{u+v_1+\cdots+v_r+w=n}
 R_u\otimes s\bar A_{v_1}\otimes\cdots\otimes s\bar A_{v_r}\otimes L_w,
 \qquad v_i>0,
\]
where the sum also runs over $r\ge0$ and $\bar A=\ker(A_{\mathrm{alg}}\to k)$.
Internal differentials and the two endpoint actions and adjacent products
form its bar differential. Its realization computes $R\otimes_A^{\mathbb L}L$
in these weight-graded categories.
Comparing it with a derived cotensor product requires a two-sided twisting
comparison on the chosen comodule category and completion.
Weight completion here is not identified with derived augmentation-adic completion.

\begin{proposition}[Centre transport under a mixed equivalence]
Suppose an equivalence between the compatible mixed module categories is supplied,
with coherent compatibility with all four bar objects and the interface operations.
It identifies the derived endomorphism objects of their identity functors.
\end{proposition}
\begin{proof}
Conjugate a natural endomorphism by the equivalence and a quasi-inverse.
The unit and counit homotopies give inverse maps on derived endomorphism objects
and preserve composition and its coherent operations.
\end{proof}

Constructing such a mixed equivalence requires the restriction, collision,
interface, and distributive-law comparisons for the actual geometric coefficients.
Its unit and counit must respect those maps and the chosen inverse limits.
A diagonal bimodule comparison must produce an actual compatible bicomodule
before tensor products can be identified with geometric cotensor products.
Transport of mixed Hochschild chains and traces further requires their
cyclic comparison and the trace domains used in their definition.
These are additional mixed Koszul--Morita constructions.
The ordinary weight-graded equivalence supplies its stated one-colour functors,
without identifying the other mixed carriers or their physical traces.

\chapter{Planar exchange and Tannaka reconstruction}
A second topological direction gives paths exchanging parallel defects.  Their homotopy classes form braid groups.  A rigid meromorphic line category $\mathscr L$ with a faithful fibre functor $\omega$ has coend
\[
 \Ocal(\mathbb G)=\int^{M\in\mathscr L}\omega(M)^*\tensor\omega(M).
\]

\begin{theorem}[Quantum symmetry from exchange]
The coend is a coquasitriangular bialgebra.  Its continuous dual acts on every line object.  The meromorphic braiding gives the universal $R$-matrix on completed matrix coefficients.  A spectral rational braiding yields an RTT presentation.
\end{theorem}
\begin{proof}
Tensor product in $\mathscr L$ gives the coproduct on matrix coefficients.  Evaluation gives the counit.  Rigidity gives the antipode on the Hopf envelope.  The braiding defines the coquasitriangular pairing, and the hexagon identities give the Yang--Baxter equation.  Writing this naturality on a generating fibre produces the RTT relation.
\end{proof}

\chapter{BRST reduction}
Let $Q$ be a square-zero derivation of a mixed datum.  Extend $Q$ factorwise to every bar word.  Since $Q$ is a derivation for each structure map, it anticommutes with every bar differential.

\begin{theorem}[BRST--bar bicomplex]
The total differential
\[
 D=Q+b_{\ch}+b_\perp+b_{\mathrm{mix}}
\]
satisfies $D^2=0$.  A bounded-below complete filtration gives a convergent spectral sequence
\[
 E_1=H_Q\bigl(\BarMix(\mathbb A;K)\bigr)
 \Longrightarrow H_D\bigl(\BarMix(\mathbb A;K)\bigr).
\]
When $Q$-cohomology is concentrated and tensor-exact, the target is the mixed bar of the reduced datum.
\end{theorem}
\begin{proof}
The square vanishes by the derivation identity and the four-colour bar theorem.  The filtration by BRST degree gives the spectral sequence.  Concentration and tensor exactness identify the $E_1$ differential with the bar differential of the reduced operations.
\end{proof}

\part{Conformal covariance and Einstein completion}


\chapter{Projective connections and the Virasoro tangent complex}
A conformal boundary algebra transforms under changes of coordinates. Projective connections record its central term. Three-dimensional gravity supplies independently defined connection and coframe fields whose boundary transformation law can be compared with this algebraic structure.

Let $\Pcal_c(X)$ be the formal moduli problem of projective connections with central parameter $c$.  At a projective connection $p$, its tangent complex is locally
\[
 \Omega^{0,\bullet}(X,T_X)
 \xrightarrow{\mathcal D_{p,c}}
 \Omega^{0,\bullet}(X,K_X^{\tensor2}),
\]
where
\[
 \mathcal D_{p,c}(\epsilon)
 =\epsilon\partial p+2(\partial\epsilon)p+\frac{c}{12}\partial^3\epsilon.
\]

\begin{theorem}[Projective covariance]
A conformal map $\Vir_c\to A$ defines a formal family of mixed structures over $\Pcal_c(X)$.  Its tangent algebra is the semidirect product of the Virasoro projective complex with the deformation algebra relative to a fixed projective connection.
\end{theorem}
\begin{proof}
The conformal vector acts on fields by infinitesimal coordinate change.  The central term is the affine cocycle in the transformation law of a projective connection.  The action therefore gives a map from the projective tangent algebra to derivations of the mixed datum.  The tangent algebra of the associated family is the semidirect product.
\end{proof}

\chapter{Three-dimensional gravity as a flat Cartan theory}
Let $M$ be an oriented three-manifold, let $\Lambda=-\ell^{-2}$, and let $\mathfrak g_\Lambda$ have basis $J_a,P_a$ with
\[
 [J_a,J_b]=\epsilon_{ab}{}^cJ_c,
 \qquad [J_a,P_b]=\epsilon_{ab}{}^cP_c,
 \qquad [P_a,P_b]=\ell^{-2}\epsilon_{ab}{}^cJ_c.
\]
The invariant pairing is $\angles{J_a,P_b}=\eta_{ab}$ and vanishes on the two diagonal summands.  A Cartan connection is
\[
 \mathcal A=\omega^aJ_a+e^aP_a.
\]
Its curvature is
\[
 F_{\mathcal A}
 =\left(R^a+\frac1{2\ell^2}\epsilon^a{}_{bc}e^b\wedge e^c\right)J_a
 +T^aP_a.
\]

\begin{theorem}[Cartan--Einstein equivalence]
On the locus where $e$ is invertible, $F_{\mathcal A}=0$ is equivalent to
\[
 T^a=0,
 \qquad
 G_{\mu\nu}+\Lambda g_{\mu\nu}=0,
 \qquad
 g_{\mu\nu}=\eta_{ab}e^a_\mu e^b_\nu.
\]
The Chern--Simons action for $\mathcal A$ equals the Einstein--Hilbert action with cosmological constant, up to its boundary term and normalization.
\end{theorem}
\begin{proof}
The $P$ component of the flatness equation is the torsion equation.  Invertibility of $e$ gives the Levi--Civita spin connection.  The $J$ component states that the curvature two-form has constant sectional curvature $-\ell^{-2}$.  In three dimensions the Riemann tensor is determined by the Ricci tensor, so this is the Einstein equation.  Expanding the Chern--Simons three-form with the off-diagonal invariant pairing gives $e\wedge R+\ell^{-2}e^3$ plus an exact form 
\cite{AchucarroTownsend,Witten3DGravity}.
\end{proof}

\chapter{Four equivalent linear complexes}
Fix a constant-curvature solution $(e,\omega)$.  There are four deformation complexes.
\begin{enumerate}
\item The flat adjoint complex $\Omega^\bullet(M,\mathfrak g_\Lambda)$ with differential $d_{\mathcal A}$.
\item The Einstein--Cartan complex on $(\delta e,\delta\omega)$ with linearized torsion and curvature.
\item The metric Einstein complex on $h_{\mu\nu}$ with linearized Einstein operator.
\item The Calabi complex resolving local Killing fields by the Killing, linearized curvature, and Bianchi operators.
\end{enumerate}

\begin{theorem}[Chain equivalence of gravitational presentations]
An invertible torsion-free dreibein gives chain maps among the four complexes.  These maps are quasi-isomorphisms on every geodesically convex open set.  Their cohomology describes infinitesimal isometries, physical linearized fields, and Bianchi identities in the four languages.
\end{theorem}
\begin{proof}
The dreibein identifies internal translations with tangent vectors and internal Lorentz rotations with skew endomorphisms.  Splitting $d_{\mathcal A}$ in this basis gives the linearized torsion and curvature operators.  Eliminating $\delta\omega$ by the linearized torsion equation gives the metric Einstein operator.  The symbol sequence is the Calabi symbol sequence and is exact on constant-curvature space.  The Poincar\'e lemma for the flat adjoint connection and exactness of the symbol homotopies give local quasi-isomorphisms.
\end{proof}

\chapter{Brown--Henneaux and the projective boundary}
Write the AdS connection as
\[
 A^\pm=\omega\pm\ell^{-1}e.
\]
The Einstein action is the difference of two $\mathfrak{sl}_2$ Chern--Simons actions with level
\[
 k=\frac{\ell}{4G}.
\]

\begin{theorem}[Brown--Henneaux central charge]
The classical asymptotic Poisson algebra of each chiral Chern--Simons sector has Virasoro central charge
\[
 c=6k=\frac{3\ell}{2G}.
\]
The projective connection carried by the boundary stress tensor is the projective variable of the preceding conformal family.
\end{theorem}
\begin{proof}
Hamiltonian reduction of the boundary $\mathfrak{sl}_2$ current Poisson algebra gives the classical Virasoro bracket.  Its central term is $6k$.  Substitution of the gravitational level gives the second equality.  The coadjoint transformation of the boundary stress tensor is the projective-connection transformation law 
\cite{BrownHenneaux}.
\end{proof}

\chapter{Einstein completion of a conformal pair}
Let $A_L$ and $A_R$ be conformal mixed boundary data with equal Brown--Henneaux central charge.  Their projective families map to the same $\Pcal_c(X)$.  Define their relative field algebra by the homotopy fibre product over the projective tangent algebra.

\begin{definition}[Candidate Einstein BV completion]
Suppose the relative algebra and a continuous $L_\infty$ action of the Cartan algebra have been supplied, together with their local realization. The associated cotangent candidate is
\[
 \EinBV_3(A_L,A_R)
 =T^*[-1]\MC\!\left(
   \mathfrak g_{\mathrm{grav}}
   \ltimes_{\mathrm{Ham}}
   \mathfrak m_{A_L,A_R}
  \right).
\]
Here $\mathfrak g_{\mathrm{grav}}$ is the Cartan deformation algebra and $\mathfrak m_{A_L,A_R}$ is the relative mixed deformation algebra after projective matching.
\end{definition}

\begin{theorem}[Cotangent action for a supplied coupled algebra]
Suppose projective matching and a continuous $L_\infty$ action have constructed the displayed complete semidirect product. Suppose also that its local realization and integrated dual pairing are defined. Its shifted cotangent has action
\[
 S(\alpha,\beta)=\langle\beta,\mathcal F(\alpha)\rangle
\]
and satisfies the classical master equation. Its Bianchi identity is the $L_\infty$ identity of the coupled algebra.
\end{theorem}
\begin{proof}
The supplied semidirect product has a square-zero homological vector field. Its cotangent lift is Hamiltonian for the displayed action and remains square-zero. This proves the classical master equation and the Bianchi identity by the preceding cotangent construction.
\end{proof}

The gravitational summand of this cotangent action is the cotangent theory of Cartan deformations. It introduces independent dual fields. To identify it with Einstein BV theory requires a comparison preserving the BV pairing and the action, including boundary terms, on the invertible-coframe sector. The flatness equations alone do not provide that comparison. A physical coupled completion must also construct the matter action and its equivariant moment map. The intended Einstein reconstruction therefore retains these field and action comparisons as existence problems.

\chapter{Higher-spin completion}
Replace $\mathfrak{sl}_2$ by $\mathfrak{sl}_N$ and choose its principal $\mathfrak{sl}_2$ embedding.  The adjoint representation decomposes as
\[
 \mathfrak{sl}_N
 \cong\bigoplus_{s=2}^{N}V_{2s-2},
 \qquad
 \sum_{s=2}^{N}(2s-1)=N^2-1.
\]

\begin{theorem}[Principal higher-spin spectrum]
The principal embedding gives one massless gauge field of each integer spin $2,3,\ldots,N$.  Classical Drinfeld--Sokolov reduction of the boundary affine Poisson algebra gives the principal $W_N$ algebra.  A pair of principal $W_N$ boundary data therefore has an $\mathfrak{sl}_N\oplus\mathfrak{sl}_N$ Einstein--Cartan completion with the same spin spectrum.
\end{theorem}
\begin{proof}
The irreducible $\mathfrak{sl}_2$ module $V_{2s-2}$ corresponds to a frame-like field of spin $s$.  The dimension identity exhausts the adjoint.  Principal Drinfeld--Sokolov reduction produces one generating current in each conformal weight $2,\ldots,N$.  The boundary currents and bulk gauge fields therefore match representation by representation 
\cite{CampoleoniHigherSpin,HenneauxRey}.
\end{proof}


\part{Exact gravitational complexes and quantum tests}


\chapter{The four complexes in coordinates}
The Cartan flatness equations and the metric Einstein equations describe the same classical solutions on the invertible-coframe locus. Their infinitesimal complexes make the comparison explicit. Quantum determinants additionally depend on gauge fixing, boundary conditions, and determinant-line data.

\section{Flat adjoint complex}
At a flat Cartan connection $\mathcal A$, the first complex is
\[
 \Omega^0(M,\mathfrak g_\Lambda)
 \xrightarrow{d_{\mathcal A}}
 \Omega^1(M,\mathfrak g_\Lambda)
 \xrightarrow{d_{\mathcal A}}
 \Omega^2(M,\mathfrak g_\Lambda)
 \xrightarrow{d_{\mathcal A}}
 \Omega^3(M,\mathfrak g_\Lambda).
\]
Degree zero is the gauge parameter, degree one is the fluctuation, degree two is the linearized curvature, and degree three is the Bianchi identity.

\section{Einstein--Cartan complex}
Write a gauge parameter as $(\xi^a,\lambda^a)$ and a fluctuation as $(u^a,v^a)=(\delta e^a,\delta\omega^a)$.  The first differential is
\[
 \begin{aligned}
 u^a&=D_\omega\xi^a+\epsilon^a{}_{bc}e^b\lambda^c,\\
 v^a&=D_\omega\lambda^a+\ell^{-2}\epsilon^a{}_{bc}e^b\xi^c.
 \end{aligned}
\]
The second differential is
\[
 \begin{aligned}
 \delta T^a&=D_\omega u^a+\epsilon^a{}_{bc}v^b\wedge e^c,\\
 \delta F^a&=D_\omega v^a+\ell^{-2}\epsilon^a{}_{bc}e^b\wedge u^c.
 \end{aligned}
\]
The third differential is the linearized pair of Cartan Bianchi identities.

\begin{proposition}[Chain identity]
The composite of the first two Einstein--Cartan differentials vanishes on a constant-curvature torsion-free background.
\end{proposition}
\begin{proof}
Apply $D_\omega^2$ to $\xi$ and $\lambda$.  Replace the background curvature by $-\ell^{-2}e\wedge e$ and use $D_\omega e=0$.  The curvature terms cancel the explicit $\ell^{-2}$ terms.  The remaining expressions are Jacobi identities for $\epsilon_{ab}{}^c$.
\end{proof}

\section{Metric complex}
The dreibein identifies the translation parameter with a vector field.  Modulo Lorentz rotations,
\[
 h_{\mu\nu}=2\eta_{ab}e^a_{(\mu}u^b_{\nu)}.
\]
The metric differential begins with the Killing operator
\[
 K(\zeta)_{\mu\nu}=\nabla_\mu\zeta_\nu+\nabla_\nu\zeta_\mu.
\]
The linearized cosmological Einstein operator is
\[
 \begin{aligned}
 (E_\Lambda h)_{\mu\nu}
 =\frac12\bigl(&-\nabla^2h_{\mu\nu}-\nabla_\mu\nabla_\nu h
 +\nabla_\mu\nabla^\rho h_{\rho\nu}
 +\nabla_\nu\nabla^\rho h_{\rho\mu}\\
 &-g_{\mu\nu}(\nabla^\rho\nabla^\sigma h_{\rho\sigma}-\nabla^2h)
 \bigr)-\Lambda(h_{\mu\nu}-g_{\mu\nu}h).
 \end{aligned}
\]
It satisfies $E_\Lambda K=0$ and $\nabla^\mu(E_\Lambda h)_{\mu\nu}=0$.

\begin{theorem}[Explicit metric reduction]
Solving the linearized torsion equation for $v=\delta\omega$ and substituting into the linearized curvature equation gives $E_\Lambda h=0$.  The kernel of the map from $(u,v)$ to $h$ is the infinitesimal Lorentz action.
\end{theorem}
\begin{proof}
The torsion equation is algebraic in the antisymmetric part of the covariant derivative of $u$ and determines $v$.  The symmetric part of $u$ is $h/2$.  Substitution into $\delta F=0$ yields the standard variation of the Ricci tensor and the displayed Einstein operator.  An antisymmetric $u_{\mu\nu}$ is exactly a Lorentz rotation of the dreibein.
\end{proof}

\section{Calabi complex}
On a constant-curvature manifold, the Calabi complex begins
\[
 \Gamma(TM)
 \xrightarrow{K}
 \Gamma(S^2T^*M)
 \xrightarrow{R'}
 \Gamma(\mathcal R)
 \xrightarrow{B'}
 \Gamma(\mathcal B),
\]
where $R'$ is the linearized Riemann tensor and $B'$ is the linearized algebraic and differential Bianchi operator.  In three dimensions the Riemann tensor is equivalent to the Einstein tensor, so the metric and Calabi complexes are locally quasi-isomorphic.

\chapter{Boundary BFV theory and the Virasoro moment map}
The variation of Chern--Simons theory on a manifold with boundary gives the boundary symplectic form
\[
 \Omega_{\partial}=\frac{k}{4\pi}\int_{\partial M}\angles{\delta A\wedge\delta A}.
\]
A chiral polarization selects $A_{\bar z}$ as coordinate.  The residual moment map is the flatness constraint.

\begin{theorem}[Boundary current algebra]
Canonical quantization of the boundary symplectic form gives the affine current algebra at level $k$.  Drinfeld--Sokolov reduction of the principal $\mathfrak{sl}_2$ constraint gives the Virasoro moment map
\[
 T(z)=\frac1{2(k+h^\vee)}:J^aJ_a:(z)+\text{improvement},
\]
with the classical Brown--Henneaux limit $c=6k$.
\end{theorem}
\begin{proof}
The inverse of $\Omega_\partial$ is the current Poisson bracket.  Quantization gives the affine central extension.  The BRST complex of the principal constraint eliminates the gauge directions and retains the improved quadratic current.  Its classical central term is $6k$ in the gravitational normalization.
\end{proof}

\chapter{Matter stress tensor and the Einstein moment map}
Let $\mathfrak m$ be the relative field algebra of a conformal pair.  Projective covariance gives a Hamiltonian action of vector fields and therefore a moment map
\[
 \mu_{\mathrm{matt}}:\mathfrak m\longrightarrow
 \Omega^{0,\bullet}(X,K_X^{\tensor2})[-1].
\]
Its classical component is the stress tensor.

\begin{theorem}[Conservation from the master equation]
The coupled BV action
\[
 S_{\mathrm{tot}}=S_{\mathrm{grav}}+S_{\mathrm{matt}}
 +\angles{h,\mu_{\mathrm{matt}}}
\]
satisfies the classical master equation precisely when the matter moment map is equivariant.  The component of that equation linear in the diffeomorphism ghost is
\[
 \nabla^\mu T_{\mu\nu}=0.
\]
\end{theorem}
\begin{proof}
The master bracket of the coupling with itself is the defect of equivariance.  The bracket with the gravitational action is the infinitesimal diffeomorphism of the moment map.  Their cancellation is the Hamiltonian equivariance identity.  Pairing the identity with the diffeomorphism ghost and integrating by parts gives stress-tensor conservation.
\end{proof}

\chapter{One-loop gravity and analytic torsion}
Expand Chern--Simons theory around an acyclic flat connection $\mathcal A$.  Gauge fixing gives the Laplacians of the flat adjoint de Rham complex.

\begin{theorem}[One-loop determinant]
The absolute value of the one-loop partition function is the square root of the Ray--Singer torsion of the flat adjoint bundle:
\[
 |Z_{1\text{-loop}}(\mathcal A)|
 =\tau_{\mathrm{RS}}(M,\ad\mathcal A)^{1/2}.
\]
The phase is the eta-invariant contribution fixed by the framing convention.
\end{theorem}
\begin{proof}
Bosonic and ghost Gaussian integrals give alternating determinants of the Laplacians in the flat de Rham complex.  Their product is the analytic torsion.  The odd signature operator supplies the determinant phase.  This is the standard stationary-phase calculation for Chern--Simons theory \cite{RaySinger1971,Witten3DGravity}.
\end{proof}

The same determinant complex is obtained from the Einstein--Cartan, metric, and Calabi presentations through their chain equivalences.  One-loop equality is therefore a consequence of quasi-isomorphism together with compatible determinant-line orientations.

\chapter{BTZ, Cardy, and the area law}
For a unitary modular boundary theory with left and right central charges $c_L,c_R$, the asymptotic entropy is
\[
 S_{\mathrm{Cardy}}
 =2\pi\sqrt{\frac{c_L}{6}\left(L_0-\frac{c_L}{24}\right)}
 +2\pi\sqrt{\frac{c_R}{6}\left(\bar L_0-\frac{c_R}{24}\right)}.
\]
For the BTZ black hole,
\[
 L_0-\frac c{24}=\frac{(r_++r_-)^2}{16G\ell},
 \qquad
 \bar L_0-\frac c{24}=\frac{(r_+-r_-)^2}{16G\ell},
 \qquad
 c=\frac{3\ell}{2G}.
\]

\begin{theorem}[Cardy--Bekenstein equality]
Substitution gives
\[
 S_{\mathrm{Cardy}}=\frac{2\pi r_+}{4G}.
\]
This is the Bekenstein--Hawking entropy of the BTZ horizon.
\end{theorem}
\begin{proof}
Each square root becomes $\pi(r_+\pm r_-)/(4G)$.  Their sum is $2\pi r_+/(4G)$.  The modular asymptotic is Cardy's theorem, and the geometric charges are the Chern--Simons holonomies of the BTZ solution \cite{Cardy1986,BTZ1992}.
\end{proof}

\chapter{Higher-spin holonomy and boundary charges}
In $\mathfrak{sl}_N$ Chern--Simons gravity, a smooth Euclidean saddle is selected by the conjugacy class of the thermal holonomy.  The invariant polynomials
\[
 \Tr(\operatorname{Hol}^s),
 \qquad 2\le s\le N,
\]
fix the spin-$s$ charges.  Drinfeld--Sokolov reduction identifies the same charges with zero modes of the $W_N$ currents.

\begin{theorem}[Bulk--boundary charge equality]
For a principal $\mathfrak{sl}_N$ connection in Drinfeld--Sokolov gauge, the invariant polynomials of the spatial holonomy equal the classical $W_N$ zero-mode Hamiltonians.  Their Poisson brackets are the classical $W_N$ brackets.
\end{theorem}
\begin{proof}
Drinfeld--Sokolov gauge writes the connection as the principal raising operator plus lowest-weight components.  The coefficients of the characteristic polynomial are the invariant polynomials and are polynomial functions of those components.  Hamiltonian reduction gives the $W_N$ Poisson structure on the same functions.
\end{proof}


\part{Anomalies, Hall sources, and automorphic completion}


\chapter{Four native anomaly classes}
A local anomaly, a determinant-line curvature, and an arity-zero bar term belong to different complexes. A cancellation in one complex becomes a geometric cancellation only through a comparison preserving the relevant class and its primitive. Hall and automorphic constructions require their coefficient and orientation maps for the same reason.

The construction carries four curvature classes.
\begin{enumerate}
\item A local BV anomaly lies in the degree-one local deformation complex.
\item A projective anomaly is the curvature of the determinant line over parameter space.
\item A framing anomaly is the transgression of a characteristic class to a three-manifold.
\item A curved Koszul term is the arity-zero component of a curved coalgebra.
\end{enumerate}

\begin{theorem}[Curvature-killing extension]
Let $\kappa$ be a closed central degree-two class in the deformation complex of a theory.  Adjoining a degree-one generator $u$ with $du=\kappa$ produces a central extension in which the lifted curvature vanishes.  Tensoring with the inverse invertible field theory gives the geometric realization of the same cancellation.
\end{theorem}
\begin{proof}
The extended differential squares to zero because $d\kappa=0$ and $\kappa$ is central.  The new Maurer--Cartan equation contains the compensating field $u$.  An invertible theory with curvature $-\kappa$ adds the opposite class under tensor product.
\end{proof}

\chapter{Hall and Calabi--Yau sources}
Let $\Ccal$ be a Calabi--Yau category with a Hall moduli problem.  An oriented extension correspondence gives a Hall algebra $H^+$.  A factorization pushforward to the boundary curve gives a chiral-convolution algebra.  Relative induction supplies the remaining colours.

\begin{proposition}[Transport under a supplied mixed Morita equivalence]
Suppose a Hall construction supplies a complete mixed boundary datum and a Morita equivalence of its compatible module category with that of $\mathbb A$. The equivalence identifies their derived centres. A comparison of local deformation algebras, BV pairings, and actions is additional data for a BV comparison. Quantum line and amplitude comparisons additionally require the monoidal and cyclic structures used to define them.
\end{proposition}
\begin{proof}
Conjugating natural transformations by the equivalence and its inverse identifies the endomorphisms of the two identity functors. The unit and counit homotopies provide the inverse comparison and preserve composition. This proves the centre comparison. It does not construct the additional local, symplectic, or cyclic maps.
\end{proof}

An algebra map alone does not induce the required forward centre map. For example, the homomorphism $\C[x]\to M_2(\C)$ with $x\mapsto\operatorname{diag}(0,1)$ sends the central element $x$ to a noncentral matrix. Its commutator with the matrix unit $E_{12}$ is $-E_{12}$. Thus even the degree-zero condition fails for a centre map extending the given algebra map. Hall multiplication maps must be supplemented by an actual coherent module comparison before centre transport can be applied.

For $X=K3\times E$, suppose a geometric Hall source supplies the complete charge-graded BPS Lie superalgebra and Cartan action required in Volume~IV. Its semidirect product with the common Cartan defines $\mathfrak b_X$. Volume~IV then constructs the completed Lie amalgam with the Igusa Borcherds Borel $\mathfrak b_\Delta$. Applying the completed enveloping functor gives the cocommutative Hopf algebra
\[
 H_{X\boxtimes\Delta}
 =\widehat U(\mathfrak b_X)
  \widehat *_{\widehat U(\mathfrak h_\Pi)}
  \widehat U(\mathfrak b_\Delta).
\]
The two Hopf retractions preserve the geometric and automorphic sectors separately.  A $q$-deformed mixed Hopf algebra is supplied by the quantum amalgam theorem when a geometric Hall coproduct and pairing are present.

\section{Mixed realization of the amalgam}
The Hopf retractions above are maps of their stated algebraic objects. A mixed realization must construct chiral and transverse operations, an interface, and a compatible distributive law. To realize the retractions, it must then supply maps preserving those operations and the chosen completions.

A BV retraction requires further maps between the local deformation algebras, together with compatible pairings and actions. The counterexample to forward centre functoriality prevents deducing these maps from the Hopf retractions alone. The intended geometric--automorphic theory asks for this mixed realization and for a primitive comparison quotient on its actual boundary operators. Its interaction ideal is the kernel of that constructed quotient, rather than an ideal determined by equality of scalar characters.

\chapter{Igusa root states and the mixed-boundary comparison}

The Gritsenko--Nikulin presentation defines the root Lie superalgebra
$\mathfrak g_{\Delta_5}$ and its positive subalgebra $\mathfrak n_+$.
Let $d_\alpha=\operatorname{sdim}(\mathfrak g_{\Delta_5})_\alpha$,
and fix the positive cone, its proper height, and Weyl vector $\rho$.
Put
\[
 P=\prod_{\alpha>0}(1-e^{-\alpha})^{d_\alpha},\qquad
 \DenIg=e^{-\rho}P=D_5(2Z),\qquad \jmath(Z)=Z/2.
\]
The products are first interpreted in the positive-height completion.

\begin{theorem}[Root character and reduced scalar]
The root Fock space $\operatorname{Sym}(\mathfrak n_+)$, formed with
super symmetry, has supercharacter $P^{-1}$. Consequently
\[
 -\jmath^*\!\left(e^{2\rho}
   \operatorname{sch}(\operatorname{Sym}(\mathfrak n_+)^{\otimes2})\right)
 =-(\jmath^*\DenIg)^{-2}=-\chi_{10}^{-1}.
\]
For a nonsingular projective K3 surface $S$ and an elliptic curve $E$, this scalar also has an enumerative interpretation. Let $N_{g,h,d}$ be the reduced disconnected Gromov--Witten invariant in class $(\beta_h,d)$, where $\beta_h$ is nonzero and primitive with square $2h-2$. Use the incidence insertion $\beta_h^\vee\otimes[\mathrm{point}]$, with $\langle\beta_h,\beta_h^\vee\rangle=1$. Then the series
\[
 \mathcal Z(u,q,s)=\sum_{g,h,d\ge0}N_{g,h,d}u^{2g-2}q^{h-1}s^{d-1}
\]
is the Laurent expansion of $-\chi_{10}(r,q,s)^{-1}$ under $r=e^{iu}$, by \cite[Theorem~1]{OP}.
\end{theorem}
\begin{proof}
An even root vector contributes $(1-e^{-\alpha})^{-1}$ to the
supercharacter of its symmetric algebra. An odd root vector contributes
$1-e^{-\alpha}$, because its symmetric algebra in super vector spaces
is an exterior algebra. Multiplying the two contributions gives
$P^{-1}$. Proper height makes every coefficient calculation finite.
The Weyl factor gives $e^{2\rho}P^{-2}=\DenIg^{-2}$.
Pullback by $\jmath$ gives $D_5(Z)^{-2}=\chi_{10}(Z)^{-1}$.
Oberdieck--Pixton, Theorem~1, identifies the separately defined reduced numerical generating series with its negative \cite{OP}. This last theorem
is an enumerative input; the character calculation supplies the root side.
\end{proof}

A primitive realization starts with parity-preserving root classes in a charge-graded Lie superalgebra $\mathfrak p$ satisfying the defining relations. Require the classes to preserve charge. The free Lie presentation gives an even charge-preserving map $f:\mathfrak n_+\to\mathfrak p$.

Applying $f$ to tensor words constructs $U(f)$. With trivial coefficients, use Chevalley chains $C_*^{\mathrm{CE}}(\mathfrak a)=\operatorname{Sym}(\mathfrak a[1])$, where the homological suspension raises degree by one and the differential contracts pairs by the bracket. Applying the suspension of $f$ to every factor gives a chain map because $f$ preserves brackets.

For the level-zero Lie conformal current, define $\operatorname{Cur}(\mathfrak a)=\C[\partial]\otimes\mathfrak a$, with $\partial$ even. Its generator bracket is $[x_\lambda y]=[x,y]$. Extend it by $[\partial x_\lambda y]=-\lambda[x_\lambda y]$ and $[x_\lambda\partial y]=(\partial+\lambda)[x_\lambda y]$. The formula $\partial^j x\mapsto\partial^j f(x)$ defines $\operatorname{Cur}(f)$. Charge preservation makes these maps compatible with the specified proper-height quotients and their inverse limits.

A complete mixed boundary realization requires further geometric maps.
The compact source consists of sheaves and extension flags on a common
elliptic torsor. Its critical realization must supply reduced coefficient
complexes, orientation cross-line maps, unit and flag coherences, and the
admissible pull--push operations. A chiral comparison must preserve actual
support and moving-insertion operations. A quantum-double comparison must
preserve coproducts, pairings and the domain of exchange operators.
The finite Calabi--Yau quiver completions require their specified dg maps;
ordinary full-subquiver inclusions do not in general give those maps.

The mixed datum to which a BV envelope is applied must therefore include
these operations and compatible maps between them. Identifying the
primitive algebra of that envelope requires an actual map of its
deformation complex to the root complex. A determinant comparison requires
a skew perfect family, a square-root line and section, and an isomorphism
to the automorphic line preserving transport and normalization.
An operator-trace interpretation further requires its state space and
trace domain. Constructing this datum and these comparisons is the Igusa
mixed-boundary problem. The root character identity specifies the scalar
that their final evaluation must recover.

\chapter{The carrier lattice}
The complete theory uses the carrier set
\[
 \mathfrak C=
 \{\mathrm{diag},\mathrm{ch},\perp,\mathrm{int},
   \mathrm{ex},\mathrm{cyc},\mathrm{mod},\mathrm{BV},\mathrm{Ein}\}.
\]
For a subset $S\subset\mathfrak C$, let $\operatorname{Str}_S(\mathbb A)$ be the homotopy limit of all mapping spaces attached to incidence cells supported in $S$.

\begin{theorem}[Carrier descent]
When $S$ and $T$ form a relation-excisive pair, the square
\[
\begin{tikzcd}
\operatorname{Str}_{S\cup T}(\mathbb A)\arrow[r]\arrow[d]&
\operatorname{Str}_S(\mathbb A)\arrow[d]\\
\operatorname{Str}_T(\mathbb A)\arrow[r]&
\operatorname{Str}_{S\cap T}(\mathbb A)
\end{tikzcd}
\]
is a homotopy pullback.  Thus a global structure is assembled from its carrier restrictions together with their common interface.
\end{theorem}
\begin{proof}
Relation-excision identifies the incidence category of $S\cup T$ with the pushout of the incidence categories for $S$ and $T$ over their intersection.  Homotopy limits send this pushout to the displayed pullback.
\end{proof}

\begin{definition}[Relation-complete carrier cover]
Let $U\subseteq\mathfrak C$ and let $\mathcal U=\{S_i\}_{i\in I}$ be a finite family with $U=\bigcup_iS_i$.  The family is \emph{relation-complete} when the canonical functor from the Cech homotopy colimit
\[
 \underset{\varnothing\ne J\subseteq I}{\operatorname{hocolim}}
 \mathcal I_{\cap_{j\in J}S_j}
 \longrightarrow \mathcal I_U
\]
is an equivalence.  Equivalently, every incidence cell and every relation among its faces is contained in one member of the cover, with all overlaps recorded by the intersections.
\end{definition}

\begin{theorem}[Power-set coherence]
The assignment
\[
 S\longmapsto\operatorname{Str}_S(\mathbb A)
\]
is a contravariant functor on the full power set of $\mathfrak C$.  For every relation-complete carrier cover $\mathcal U=\{S_i\}_{i\in I}$ of $U$, the restriction map is an equivalence
\[
 \operatorname{Str}_U(\mathbb A)
 \xrightarrow{\sim}
 \underset{\varnothing\ne J\subseteq I}{\operatorname{holim}}
 \operatorname{Str}_{\cap_{j\in J}S_j}(\mathbb A).
\]
Thus every subtheory, every intersection of subtheories, and the complete theory are values of one descent object.
\end{theorem}
\begin{proof}
An inclusion $S\subseteq T$ induces the full inclusion $\mathcal I_S\subseteq\mathcal I_T$ and therefore the restriction map on homotopy limits.  Composition of inclusions gives functoriality.  A relation-complete cover presents $\mathcal I_U$ as the displayed homotopy colimit of its intersection diagram.  Mapping that colimit into the diagram of structure spaces converts it into the stated homotopy limit.  The resulting equivalence is natural under refinement of the carrier cover.
\end{proof}

\begin{definition}[Exact-support deformation complex]
For $S\subseteq\mathfrak C$, let $\Def_{\leq S}$ be the subcomplex generated by incidence cells supported in $S$, and set
\[
 \Def_{=S}
 =\operatorname{cofib}\!\left(
   \underset{T\subsetneq S}{\operatorname{hocolim}}
   \Def_{\leq T}
   \longrightarrow
   \Def_{\leq S}
 \right).
\]
Thus $\Def_{=S}$ retains the deformation cells whose minimal carrier support is exactly $S$.
\end{definition}

\begin{theorem}[Carrier-support spectral sequence]
Let $\Def_{\mathbb A}^{\mathrm{full}}$ be the cellular convolution complex of the complete incidence diagram.  Filter it by the cardinality of minimal carrier support.  The finite filtration gives a convergent spectral sequence
\[
 E_1^{p,q}
 =\bigoplus_{\substack{S\subseteq\mathfrak C\\|S|=p}}
 H^{p+q}(\Def_{=S})
 \Longrightarrow
 H^{p+q}\!\left(\Def_{\mathbb A}^{\mathrm{full}}\right).
\]
A class on the $S$-summand is a deformation or curvature class whose first complete geometric carrier is $S$.  The first differential is the signed incidence map between exact-support quotients.
\end{theorem}
\begin{proof}
The support-cardinality filtration is finite because $\mathfrak C$ is finite.  Its $p$th associated-graded quotient is the direct sum of the exact-support complexes $\Def_{=S}$ with $|S|=p$.  The spectral sequence of a finite filtered complex gives the displayed first page and converges strongly to the cohomology of the total convolution complex.  The cellular determinant orientation gives the incidence signs.
\end{proof}

\chapter{Sixteen structural laws}
The theory is governed by the following laws.
\begin{enumerate}
\item A geometric carrier determines the type of its operation.
\item The diagonal stores the complete principal-part filtration.
\item The interval stores ordered fusion.
\item The plane stores exchange.
\item The circle stores trace.
\item The stable curve stores modular sewing.
\item The transverse, associative-chiral, and chiral-Lie bars are distinct.
\item Chiral PBW compares the last two on enveloping objects.
\item The interface bar records the mixed distributive law.
\item The derived centre is the universal interior action.
\item The deformation algebra is the tangent algebra of the centre-extended datum.
\item Maurer--Cartan curvature is the field equation.
\item The $L_\infty$ identity is the Bianchi identity.
\item Shifted cotangent is the universal classical BV envelope.
\item Hodge-completed derived de Rham theory is the intrinsic algebraic quantum complex of the oriented formal moduli problem.
\item A gravitational realization requires projective matching, a coupled local field algebra, and a comparison preserving the physical BV action and pairing.
\end{enumerate}

\chapter{Coherent actions and gravitational reconstruction}
A supplied mixed boundary datum defines compatible modules and their derived natural endomorphisms. A specified coloured operad organizes deformations of that datum together with its central action. These are constructions relative to the full operations and their coherence maps. The centre alone does not recover the datum from which the deformation algebra is formed.

A supplied local realization of the deformation algebra admits shifted cotangent transgression. With compatible invariant orientations, the finite perfect charts also carry the algebraic BV differential. The physical bulk is defined independently by its fields, action, gauge symmetry, boundary conditions, and quantization. Reconstruction requires a coherent action map from its observables to an independently defined geometric centre, followed by an equivalence theorem on the stated sector.

For gravity, the invertible coframe gives the classical Cartan--Einstein correspondence. The cotangent construction supplies a formal comparison theory with extra dual fields. Its identification with gravitational BV fields requires an action-preserving and pairing-preserving map. Quantum boundary currents, a continuum measure, global holonomy sectors, allowed fillings, and compatible states remain necessary for the quantum and global comparison.

Hall and Calabi--Yau geometry can supply boundary operations when their correspondences, coefficients, orientations, and comparison maps are constructed. Exchange additionally uses a monoidal category and a fibre functor. Cyclic pairings and sewing determine amplitudes only on their stated domains. These comparisons are the mathematical steps by which local boundary operations could reconstruct physical and geometric theories.

\part{Technical appendices to the preceding book}

\chapter{The Chern--Simons normalization}
The gravitational charge calculation uses the action normalization. The mixed bar calculation uses the determinant orientation of its edges. The following formulas specify these conventions.

With invariant pairing $\angles{J_a,P_b}=\eta_{ab}$, the Chern--Simons action is
\[
 S_{\mathrm{CS}}[\mathcal A]
 =\frac{k}{4\pi}\int_M
 \angles{\mathcal A\wedge\dd\mathcal A
 +\frac23\mathcal A\wedge\mathcal A\wedge\mathcal A}.
\]
For $A^\pm=\omega\pm\ell^{-1}e$ one finds
\[
 S_{\mathrm{CS}}[A^+]-S_{\mathrm{CS}}[A^-]
 =\frac{k}{2\pi\ell}\int_M
 \epsilon_{abc}e^a\wedge
 \left(R^{bc}+\frac1{3\ell^2}e^b\wedge e^c\right)
 +\int_{\partial M}(\cdots).
\]
The equality with the Einstein--Hilbert normalization gives $k=\ell/(4G)$.

\chapter{The principal embedding count}
The exponents of $\mathfrak{sl}_N$ are $1,2,\ldots,N-1$.  Under the principal $\mathfrak{sl}_2$ embedding,
\[
 \mathfrak{sl}_N=\bigoplus_{m=1}^{N-1}V_{2m}.
\]
The dimensions satisfy
\[
 \sum_{m=1}^{N-1}(2m+1)
 =(N-1)(N+1)=N^2-1.
\]
Writing $s=m+1$ gives one field at every spin $s=2,\ldots,N$.

\chapter{Low-arity mixed signs}
Let $x|k|y$ be a suspended interface word.  The two length-lowering contractions are
\[
 b_L[x|k|y]=(-1)^{|x|}(xk)|y,
 \qquad
 b_R[x|k|y]=(-1)^{|x|+|k|}x|(ky).
\]
Applying the remaining contraction gives the two composites $(xk)y$ and $x(ky)$ with opposite determinant-line signs.  Their sum vanishes by the interface associativity relation.  A transverse and a chiral contraction acquire the product orientation
\[
 \det(E_{\ch}\oplus E_\perp)
 =\det(E_{\ch})\tensor\det(E_\perp),
\]
so reversing their order contributes the Koszul sign and yields anticommutation.

\backmatter
\bibliographystyle{alpha}
\bibliography{references}
\end{document}
